
%%FILE blueprint/src/content.tex
\input{content/central_simple}
\input{content/morita}
\chapter{Results in Noncommutative
  Algebra}\label{sec:results-noncommutative-algebra}
\input{content/unclassified}
\input{content/wed_artin}
\input{content/skolem_noether}
\input{content/double_centralizer}
\input{content/brauer-group}
\input{content/brauer-group2}
\input{content/ReducedCharPoly}
%%% Local Variables:
%%% mode: LaTeX
%%% TeX-master: "print"
%%% End:


%%FILE blueprint/src/content/Brauergroup_over_extension.tex
\chapter{BrauerGroup over Extension}

\section{Splitting Field of CSA}
Let $k$ be a field and $A$ a finite dimensional central simple algebra over $k$. 


%%FILE blueprint/src/content/ReducedCharPoly.tex
% !TEX root = ../print.tex
\chapter{Division Algebra over Local Fields}\label{cha:div_over_local}

\section{Reduced Norm \& Trace}
% \begin{definition}
    
% \end{definition}


%%FILE blueprint/src/content/brauer-group.tex
% !TEX root = ../print.tex
\chapter{Brauer Group}\label{cha:brauer-group}

\section{Construction of Brauer Group}
Let $K$ be a field. We denote the class of finite dimensional central simple $K$-algebras as $\CSA_{K}$. When $K$ is clear, we drop the subscript.

\begin{remark}
  By~\cref{lem:tensor-central} and~\cref{thm:tensor-csa}, $\CSA$ is closed under tensor product, that is if $A, B\in\CSA$, we have $A\otimes_{K} B\in\CSA$ as well.
\end{remark}

\begin{definition}[Brauer Equivalence]\label{def:br-eqv}
  For any two $A, B\in\CSA$, we say $A$ and $B$ are Brauer equivalent, when there exists $m, n \in \mathbb{N}_{\ge0}$ such that $\Mat_{m}(A)\cong \Mat_{n}B$ as $K$-algebras. We denote this relation as $A\sim_{\Br_{K}} B$, when $K$ is clear, we drop the subscript.
\end{definition}

\begin{remark}
  Isomorphic $K$-algebras are Brauer equivalent.
\end{remark}

\begin{lemma}
  $\sim_{\Br}$ is reflexive.
  \leanok
  \lean{IsBrauerEquivalent.refl}
\end{lemma}
\begin{proof}
  Indeed, $A \cong \Mat_{1}(A)$ as $K$-algerbas.
\end{proof}

\begin{lemma}
  $\sim_{\Br}$ is symmetric.
  \leanok
  \lean{IsBrauerEquivalent.symm}
\end{lemma}
\begin{proof}
  Indeed, just exchange $m$ and $n$.
\end{proof}

\begin{lemma}
  $\sim_{\Br}$ is transitive.
  \leanok
  \lean{IsBrauerEquivalent.trans}
  \uses{con:matrix-matrix}
\end{lemma}
\begin{proof}
  Let $A\sim_{\Br}B$ and $B\sim_{\Br}C$; that is for some $m,n,p, q\in\mathbb{N}_{\ge0}$ we have $\Mat_{n}(A)\cong\Mat_{m}(B)$ and $\Mat_{p}(B)\cong \Mat_{q}(C)$ as $K$-algebras. Hence, from~\cref{con:matrix-matrix}, we have the following:
  \[
    \begin{aligned}
      \Mat_{np}(A)&\cong \Mat_{p}\left(\Mat_{n}(A)\right)\cong\Mat_{p}\left(\Mat_{m}(B)\right)\\
                  &\cong\Mat_{mp}(B)\cong\Mat_{m}\left(\Mat_{p}(B)\right)\\
      &\cong\Mat_{m}\left(\Mat_{q}(C)\right)\cong\Mat_{mq}(C).
    \end{aligned}
  \]
  In another word, $A\sim_{\Br}C$.
\end{proof}

Hence $\sim_{\Br}$ is really an equivalence relation, we denote the quotient ${}^{\CSA}/_{\sim_{\Br}}$ as $\Br(K)$.

\begin{lemma}\label{lem:br-mul-wd}
  $(\bullet\otimes_{K}\bullet):\CSA\times\CSA \to \CSA$ descends to a function on $\Br(K)$.
  \leanok
  \lean{BrauerGroup.eqv_tensor_eqv}
  \uses{con:matrix-tensor-matrix}
\end{lemma}
\begin{proof}
  We need to prove that for all $A, B, C, D \in \CSA$ such that $A\sim_{\Br} B$ and $C\sim_{\Br}D$, $A\otimes_{R}C \sim_{\Br} B\otimes_{R}D$ as well.
  Suppose $\Mat_{m}(A)\cong \Mat_{n}(B)$ as $K$-algebras and $\Mat_{p}(C)\cong\Mat_{q}(D)$, by~\cref{con:matrix-tensor-matrix}, we have
  \[
    \begin{aligned}
      \Mat_{mp}\left(A\otimes_{R} C\right)&\cong \Mat_{m}(A)\otimes_{R}\Mat_{p}(C) \\
                                          &\cong \Mat_{n}(B)\otimes_{R}\Mat_{q}(D)\\
      &\cong \Mat_{nq}\left(B\otimes_{R}D\right).
    \end{aligned}
  \]

\end{proof}

\begin{construction}[Brauer Group]\label{con:br}
  $\Br(K)$ forms a group under $[A]_{\sim_{\Br}}\cdot[B]_{\sim_{\Br}}=[A\otimes_{K}B]_{\sim_{\Br}}$ with neutral element $[K]_{\sim_{\Br}}$ where $A, B\in\CSA$ and $[A]_{\sim_{Br}}^{-1}=[A^{\opp}]_{\sim_{\Br}}$. We need to prove the following properties:
  \begin{enumerate}
    \item associativity: for all $A, B, C\in\CSA$, $[A]_{\sim_{\Br}}\cdot\left([B]_{\sim_{\Br}}\cdot[C]_{\sim_{\Br}}\right)=\left([A]_{\sim_{\Br}}\cdot[B]_{\sim_{\Br}}\right)\cdot [C]_{\sim_{\Br}}$ because $A\ox_{R}\left(B\ox_{R}C\right)\cong \left(A\ox_{R}B\right)\ox_{R}C$ as $K$-algebras.
    \item neutral element: for all $A\in \CSA$, $[K]_{\sim_{\Br}}\cdot[A]_{\sim_{\Br}}=[A]_{\sim_{\Br}}=[A]_{\sim_{\Br}}\cdot [K]_{\sim_{\Br}}$. Since $[K]_{\sim_{\Br}}\cdot[A]_{\sim_{\Br}}=[K\ox_{K} A]_{\sim_{\Br}}$, in~\cref{con:matrix-tensor-matrix}, we see that $\Mat_{n}(A)\cong A\ox_{K}\Mat_{n}(K)$, by~\cref{lem:br-mul-wd}, $A\ox_{K}\Mat_{n}(K)$ is Brauer equivalent to $A\ox_{K} K$ since $K\sim_{\Br}\Mat_{n}(K)$.
    \item cancellation: for all $A \in \CSA$, we need $[A]_{\sim_{\Br}}\cdot[A^{\opp}]_{\sim_{\Br}}$, that is we want $A\ox_{K}A^{\opp}\sim_{\Br} K$. By~\cref{con:self-tensor-opp-iso-end}, we have $A\ox_{K}A^{\opp}\cong \End_{K}A$ which is isomorphic to $\Mat_{\dim_{K}A}(K)$ as $K$-algebras.
  \end{enumerate}
  \leanok
  \lean{BrauerGroup.Bruaer_Group}
  \uses{con:matrix-tensor-matrix,lem:br-mul-wd,con:self-tensor-opp-iso-end}
\end{construction}

\begin{theorem}
  \label{thm:br-triv-alg-closed}
  If $K$ is algebraically closed, $\Br(K)$ is trivial; in particular $\br_{n}(\mathbb{C})$ is trivial.
  \leanok
  \lean{BrauerGroup.Alg_closed_eq_one}
  \uses{thm:wed-artin-algebra}
\end{theorem}
\begin{proof}
  We need to show that every $A\in\CSA$ is isomorphic to $\Mat_{n}(K)$ for some $K$ when $K$ is algebraically closed. Indeed, by~\cref{thm:wed-artin-algebra}, $A \cong \Mat_{n}(D)$ for some division algebra $D$ and $n\in\nat_{\ge0}$. Since $K$ is algebraically closed and $D$ is an integral domain and finite dimensional, the structure morphism $\rho: K\to D$ is a isomorphism; therefore $A\cong \Mat_{n}(K)$.
\end{proof}

\begin{lemma}
  \label{lem:common-div-alg}
Let $A, B \in \csa_{K}$. There exists a division $K$-algebra $D$ and non-zero $m,n\in\mathbb{N}$ such that $A\cong\mat_{m}(D)$ and $B\cong\mat_{n}(D)$ as $K$-algebras.
  \leanok
  \lean{IsBrauerEquivalent.exists_common_division_algebra}
  \uses{thm:wed-artin-algebra,thm:wed-artin-uniq}
\end{lemma}

\begin{proof}
  By~\cref{thm:wed-artin-algebra}, we can find division algebras $S_{A}, S_{B}$ and non-zero $m, n\in\mathbb{N}$ such that $A\cong\mat_{n}\left(S_{A}\right)$ and $B\cong\mat_{m}\left(S_{B}\right)$ as $K$-algebras. Hence $B\sim_{\br}A\sim_{\br}\mat_{n}\left(S_{A}\right)\sim_{\br}S_{A}$, in another word, for some non-zero $a, a'\in\mathbb{N}$, we have $\mat_{a}(B)\cong\mat_{a'}\left(S_{A}\right)$ as $K$-algebras. Hence, by~\cref{thm:wed-artin-uniq}, we have that $S_{A}\cong S_{B}$ as $K$-algebras and the lemma is proved.
\end{proof}

\section{Base Change}
In this section let $E/K$ be a field extension. We have seen in~\cref{cor:csa-basechange} that if $A\in\CSA_{K}$ then $E\otimes_{K}A\in\CSA_{E}$; therefore we have a set-theoretic function $\CSA_{K}\to\CSA_{E}$. In this section we prove that this descends to a group homomorphism $\Br(K)\to\Br(E)$. For brevity, if $A\in \CSA_{K}$, we denote $E\ox_{K}A$ as $A_{E}$ when this causes no confusion.
\begin{construction}
  \label{con:br-base-change}
  We will construct a series of isomorphisms (either over $K$ or $E$) to arrive at the conclusion that $A\sim_{\Br_{K}}B$ implies $A_{E}\sim_{\Br_{E}}B_{E}$. Assume $m,n\in\nat_{\ge0}$ are such that $\Mat_{m}(A)\cong\Mat_{n}(B)$ are $K$-algebras. Then we do the following calculation: as $E$-algebras
  \[
    \begin{aligned}
      \Mat_{m}\left(A_{E}\right)
      &\cong A_{E}\otimes_{E}\Mat_{m}(E) &\text{see~\cref{con:matrix-tensor-matrix}}\\
      &\cong A_{E}\otimes_{E}\left(E\ox_{K} \Mat_{m}(K)\right) & \text{see}~\dagger\\
      &\cong E \ox_{K}\left(A\ox_{K}\Mat_{m}(K)\right) &\text{see}~\ddagger\\
      &\cong E\ox_{K}\Mat_{m}(A) &\text{see~\cref{con:matrix-tensor-matrix} and}~\dagger\!\!\dagger \\
      \Mat_{n}\left(B_{E}\right) &\cong E\ox_{K}\Mat_{n}(B) &\text{same as the case of}~A\\
      Mat_{m}\left(A_{E}\right)&\cong \Mat_{n}\left(B_{E}\right) &\text{see}~\dagger\!\!\dagger
    \end{aligned}.
  \]

   \noindent$\dagger$: Wee need to check $\Mat_{m}(E)\cong E\otimes_{K}\Mat_{m}K$ as $E$-algebras since~\cref{con:matrix-tensor-matrix} only gives a $K$-algebra isomorphism. If $e \in E$, then its image in $E\ox_{K}\Mat_{m}(K)$ is $e\otimes 1$ and its image in $\Mat_{m}(E)$ is $\diag(e)$ which under the $K$-algebra isomorphism is mapped to $\sum_{ij}\diag(e)_{ij}\cdot \delta_{ij}=e\ox1$.

   \noindent$\ddagger$: This is defined by combining two $E$-algebra homomorphisms
   \[
     A_{E}\to A_{E}\ox_{K}\Mat_{m}(K)\to E\ox_{K}\left(A\ox_{K}\Mat_{m}(K)\right)\]
     and
     \[
       E\ox_{K}\Mat_{m}(K)\to \left(E\ox_{K}\Mat_{m}(K)\right)\ox_{K}A \to E\ox_{K}(A\ox_{K}\Mat_{m}(K)).
    \]
   Since $\left(E\ox_{K}A\right)\ox_{E}\left(E\ox_{K}\Mat_{m}(K)\right)$ is a simple ring, this morphism is automatically injective. It is surjective as well: let $x\in E\ox_{K}\left(A\ox_{K}\Mat_{m}(K)\right)$, without loss of generality, assume $x=e\ox (a\ox\delta_{ij})$ for some $e\in E$, $a\in A$. Then precisely $\left(e\ox a\right)\ox\left(1\ox\delta_{ij}\right)$ is mapped to $x$.

   \noindent$\dagger\!\dagger$: a $K$-algebra isomorphism $A\cong B$ gives an $E$-algebra isomorphism $E\ox_{K}A \cong E\ox_{K} B$.

   Thus we have a well defined function $\Br(K)\to\Br(E)$. We now check that this is a group homomorphism. $[K]_{\sim_{\Br_{K}}}$ is mapped to $[E\ox_{K}K]_{\sim_{\Br_{E}}}$ but $E\ox_{K}K\cong E$ as $E$-algebra. For $A, B\in\CSA_{K}$, we have that $[AB]_{\sim_{\Br_{K}}}$ is mapped to $\left(A\ox_{K}B\right)_{E}\cong A_{E}\ox_{E} B_{E}$ as $E$-algebras; hence $[AB]_{\sim_{\Br_{K}}}$ and $[A]_{\sim_{\Br_{K}}}\cdot[B]_{\sim_{\Br_{K}}}$ have the same image under base change.
   \leanok
   \lean{BrauerGroupHom.BaseChange}
   \uses{con:matrix-tensor-matrix}
 \end{construction}
 Denote the base change morphism in~\cref{con:br-base-change} as $\Br_{K}^{E}$.
 \begin{lemma}\label{lem:br-base-change-self}
   $\Br_{K}^{K}$ is identity.
   \leanok
 \end{lemma}
 \begin{proof}
   If $A\in\CSA$, then $A\sim_{\Br}K\ox_{K}A$.
 \end{proof}

 \begin{lemma}\label{lem:br-base-change-tower}
   Consider the tower of field extension $E/F/K$,
   \[
     \Br_{K}^{E} = \Br_{E}^{F}\circ\Br_{K}^{E}.
   \]
 \end{lemma}

 \begin{proof}
   If $A\in\CSA_{K}$, then $E\ox_{F}\left(F\ox_{K}A\right)$ is isomorphic to $E \ox_{F} A$ as $E$-algebras.
   \leanok
   \lean{BrauerGroupHom.baseChange_idem}
 \end{proof}

 \begin{corollary}
   $\Br$ forms a functor from category of field to category of abelian groups.
   \leanok
   \lean{BrauerGroupHom.Br}
   \uses{lem:br-base-change-self, lem:br-base-change-tower}
 \end{corollary}
 \begin{proof}
   This is the categorical version of~\cref{lem:br-base-change-self} and~\cref{lem:br-base-change-tower}.
 \end{proof}

 \begin{definition}[Relative Brauer Group]
   Let $E/K$ be a field extension, we define the relative Brauer group $\Br(E/K)$ to be the kernel of the base change morphism $\Br_{K}^{E}$.
   \leanok
   \lean{RelativeBrGroup}
   \uses{con:br-base-change}
 \end{definition}

 \begin{remark}
   Unpacking the definition of the relative Brauer group, we see that for any $A \in \CSA_{K}$, if $E\ox_{K}A\cong\mat_{n}(E)$ as $E$-algebras, then $\br^{E}_{K}\left([A]_{\sim_{\br}}\right)=1$.
 \end{remark}

 \begin{definition}[Splitting Field]
   For any field extension $E/K$ and any $K$-algebra $A$, we say $E$ is a splitting field of $A$ if and only if $E\ox_{K}A \cong \mat_{n}(E)$ as $E$-algebras for some non-zero $n$. We also say $E$ splits $A$ or $A$ is splited by $E$
   \leanok
   \lean{isSplit}
 \end{definition}

  \begin{theorem}
   \label{thm:split-iff-mem-relative}
   Let $E/K$ be a field extension and $A\in\csa_{K}$, $E$ splits $A$ if and only if $[A]_{\sim_{\br}}\in\br(E/K)$.
   \leanok
   \lean{BrauerGroup.split_iff}
   \uses{thm:wed-artin-algebra,thm:wed-artin-uniq}
 \end{theorem}

 \begin{proof}
   The ``only if'' part is by definition. For the other direction, we know by definition that $\mat_{n}(E\ox_{K}A)\cong\mat_{m}(E)$ as $E$-algebras for some non-zero $m, n$. By~\cref{thm:wed-artin-algebra}, we find some division algebra $D$ and non-zero natural number $p$ such that $E\ox_{K}A\cong \mat_{p}(D)$ as $E$-algebras. Thus $\mat_{pm}(E)\cong\mat_{pn}\left(E\ox_{K}A\right)\cong\mat_{p^{2}n}(D)$ as $E$-algebras. By~\cref{thm:wed-artin-uniq}, we conclude that $E\cong D$ as $E$-algebras. Hence $E\ox_{K}A\cong \mat_{p}(E)$, in another word, $E$ splits $A$.
 \end{proof}

 \begin{remark}
   In light of~\cref{lem:br-base-change-tower}, if $K$ is algebraic closed then $K$ splits any $K$-algebra $A$. Indeed, $K$ splits $A$ if and only if $[A]_{\sim\br}$ but $[A]_{\sim\br}$ is equal to $1$.
 \end{remark}

 \begin{remark}
   If two $\csa_{K}$ are Brauer equivalent, in another word, $A \sim_{\br_{K}} B$, then $E$ splits $A$ if and only if $E$ splits $B$. Indeed, if $A$ and $B$ are equivalent, then $[A]_{\sim\br} \in \br(E/K)$ if and only if $[B]_{\sim\br}\in\br(E/K)$.
 \end{remark}


 \section{Good Representative Lemma}
 In this section, let $K/F$ be a finite dimensional field extension.

 \begin{lemma}
   \label{lem:good-rep-inv}
   Let $A\in\csa_{F}$ splitted by $K$. There exists a $B\in\csa_{F}$ such that
   \begin{itemize}
     \item $[A]_{\sim_{\br}}[B]_{\sim_{\br}}=1$
     \item there exists $F$-algebra map $K\hookrightarrow B$
     \item ${\left(\dim_{F}K\right)}^{2}=\dim_{F}B$.
   \end{itemize}
   \leanok
   \lean{exists_embedding_of_isSplit}
   % \uses{thm:double-centralizer, lem:simple-centralizer, cor:self-tensor-centralizer, lem:dim-centralizer}
 \end{lemma}
\begin{proof}
   Since $K$ splits $A$, we find a non-zero natural number $n$ such that $K\ox_{F}A\cong\mat_{n}K\cong\End_{K}\left(K^{n}\right)$ as $K$-algebras.
   We define an $F$-algebra map $\iota: A\to\End_{F}\left(K^{n}\right)$ by
   \begin{center}
     \begin{tikzcd}
       A \arrow{r} & K \ox_{F} A \arrow{r}{\cong} & \End_K\left(K^n\right) \arrow{r}{|_{F}} & \End_F\left(K^n\right)
     \end{tikzcd},
   \end{center}
   where $|_{F}$ is restriction of scalars. Since $A$ is simple, $\iota$ is injective, therefore $A \cong \iota(A)$ as $F$-algebras. Define $B$ as $C_{\End_{F}\left(K^{n}\right)}(\iota(A))$, the centralizer of the range of $\iota$ in $\End_{F}\left(K^{n}\right)$.
   We construct an embedding $K \hookrightarrow B$ by $r \mapsto (r \cdot \bullet)$

   $B$ is a central $F$-algebra: if $x \in Z(B)$, then $x \in \iota(A)$ because by~\cref{thm:double-centralizer}, it is sufficient to prove that $x$ is in $C_{\End_{F}\left(K^{n}\right)}\left(B\right)$ which follows from the fact that $x \in Z(B)$. In fact, $x \in Z(\iota(A))$: suppose $a \in A$, we need to check $x \cdot \iota(a) = \iota(a)\cdot x$, this is the case because $B$ is defined as the centralizer of $\iota(A)$. Since $\iota(A) \cong A$ as $F$-algebras, $\iota(A)$ is $F$-central, hence $x \in F$.

   $B$ is a simple ring: by~\cref{lem:simple-centralizer}, it is sufficient to prove that $\iota(A)$ is a simple ring which comes from $A \cong \iota(A)$ as $F$-algebras.

   By~\cref{cor:self-tensor-centralizer}, we have $F$-algebra isomorphism $\End_{F}\left(K^{n}\right)\cong \iota(A)\ox_{F}B \cong A \ox_{F}B$. Since $\End_{F}\left(K^{n}\right)\cong\mat_{\dim_{F}\left(K^{n}\right)}\left(F\right)$ as $F$-algebras, we see that $[A]_{\sim_{\br}}$ and $[B]_{\sim_{\br}}$ are inverses.

   By~\cref{lem:dim-centralizer}, $\dim_{F}B\cdot \dim_{F}\iota(A) = \dim_{F}B\cdot\dim_{F}A = \dim_{F}\End_{F}\left(K^{n}\right) = {\left(\dim_{F}\left(K^{n}\right)\right)}^{2}={\left(\dim_{F}K\cdot \dim_{K}\left(K^{n}\right)\right)}^{2}=n^{2}\cdot{\left(\dim_{F}K\right)}^{2}$. On the other hand, since $K\ox_{F}A\cong\mat_{n}K$, we have $\dim_{F}K\ox_{F}A=\dim_{F}K\cdot\dim_{F}A=\dim_{F}\mat_{n}K=\dim_{F}K\dim_{K}\mat_{n}K=n^{2}\dim_{F}K$. Since $\dim_{F}K\ne 0$ , we conclude $\dim_{F}A=n^{2}$. Since $n\ne 0$ and $\dim_{F}(B)\cdot \dim_{F}(A)=n^{2}\dim_{F}(B)=n^{2}{\left(\dim_{F}K\right)}^{2}$, we get the desired result.

 \end{proof}

 \begin{corollary}\label{cor:good-rep}
   Let $A\in\csa_{F}$ splitted by $K$. There exists a $B \in \csa_{F}$ such that
   \begin{itemize}
     \item $[B]_{\sim_{\br}}=[A]_{\sim_{\br}}$
     \item there exists an $F$-algebra map $K\hookrightarrow B$
     \item  ${\left(\dim_{F}K\right)}^{2}=\dim_{F}B$.
   \end{itemize}
   % \uses{lem:good-rep-inv}
 \end{corollary}

 \begin{proof}
   Let $B$ and $\iota : K \hookrightarrow B$ be as in~\cref{lem:good-rep-inv}. Consider $B^{\opp}$ and $K \hookrightarrow B \to B^{\opp}$. This works.
 \end{proof}

 \begin{theorem}\label{thm:good-rep-iff-split}
   Let $A\in\csa_{F}$. $K$ splits $A$ if and only if there exists a $B\in\csa_{F}$ such that
   \begin{itemize}
     \item $[B]_{\sim_{\br}}=[A]_{\sim_{\br}}$
     \item there exists an $F$-algebra map $K \hookrightarrow B$
     \item ${\left(\dim_{F}K\right)}^{2}=\dim_{F}B$.
   \end{itemize}
   \leanok
   \lean{isSplit_iff_dimension}
   % \uses{cor:good-rep, cor:alghom-bijective-of-dim-eq}
 \end{theorem}
 \begin{proof}
   The ``if'' direction is~\cref{cor:good-rep}. For the ``only if'' direction, let $B \in\csa_{F}$ and $\iota : K \hookrightarrow B$ be given. We give $B$ a $K$-module structure by right multiplication, that is for any $a \in K$ and $b \in B$, we define $a \cdot b := b \cdot \iota(a)$.
   Since $B$ is a finite dimensional $F$-vector space and $K/F$ is a finite dimensional field extension, $B$ is a finite dimensional $K$-vector space as well. Since $[B]_{\sim_{\br}}=[A]_{\sim_{\br}}$, it is sufficient to show that $K$ splits $B$.
   We define an $F$-bilinear map $\mu : K \to B \to \End_{K}B$ by $(c, a) \mapsto (c \cdot a \cdot \bullet)$ which induce an $F$-linear map $\mu' : K \ox_{F} B \to \End_{K}B$. Since for any $r, c\in K$ and $a \in B$, we have $\mu'\left(r \cdot c\ox a\right)(a') = a a' \iota(rc)=aa'\iota(c)\iota(r)=r\cdot \mu'(c\ox a)$, that is $\mu'$ is $K$-linear as well. Note that \[\mu'(1)=\mu'(1\ox 1) = (1\cdot 1\cdot \bullet) = 1\] and that
   \[
     \begin{aligned}
       \mu'(c\ox a \cdot c'\ox a')(a'') &= \mu'(cc'\ox aa')(a'') \\
                                        &= cc' \cdot aa' \cdot a'' \\
                                        &= aa' a'' \iota(c c') \\
                                        &= a(a' a'' \iota(c'))\iota(c) \\
                                        &=\mu'(c \ox a)(a' a'' \iota(c'))\\
                                        &=\mu'(c\ox a)\left(\mu'(c'\ox a')(a'')\right) \\
                                        &= \left(\mu'(c\ox a)\circ\mu'(c'\ox a')\right)(a'')
     \end{aligned},
     \]
     that is, $\mu'$ is an $K$-algebra map.

   If we can show that $\mu'$ is a bijection, we will prove the result for $K\ox_{F} B \cong \End_{K} B \cong \mat_{\dim_{K} B}K$ as $K$-algebras. By~\cref{cor:alghom-bijective-of-dim-eq}, it is sufficient to show $\dim_{K}K\ox_{F}B = \dim_{K}\End_{K}B$. Let $n$ denote $\dim_{F}K$. Since, $\dim_{F}K\dim_{K}K\ox_{F}B =\dim_{F}K\ox_{F}B=\dim_{F}K \dim_{F}B$. we have $\dim_{K}K\ox_{F}B = \dim_{F}B = {\left(\dim_{F}K\right)}^{2}$. On the other hand, since ${\left(\dim_{F}K\right)}^{2}=\dim_{F} B = \dim_{F}K\dim_{K}B$, we have $\dim_{K} B = \dim_{F}K$; thus $\dim_{K}\End_{K} B = {\left(\dim_{K}B\right)}^{2}={\left(\dim_{F}K\right)}^{2}$ and the result is proved.
 \end{proof}

 In light of~\cref{thm:good-rep-iff-split}, we isolate the following useful definition:
 \begin{definition}[Good Representation]\label{def:good-rep}
   For any $X \in \br(F)$, a $K$-good representation of $X$ is an $A \in \csa_{F}$ and an $F$-algebra map $K \hookrightarrow A$ such that $[A]_{\sim_{\br}} = X$ and $\dim_{F}A={\left(\dim_{F}K\right)}^{2}$. We often denote the $F$-algebra map $K \hookrightarrow A$ as $\iota$ or $\iota_{A}$.

   \leanok
   \lean{GoodRep}
 \end{definition}

 When $K$ is clear from context, we will simply say good representation instead of $K$-good representation

 \begin{corollary}
   \label{cor:mem-relative-br-iff-good-rep}
   For any $X\in\br(F)$, $X\in\br(K/F)$ if and only if $X$ admits a good representation.
   \leanok
   \lean{mem_relativeBrGroup_iff_exists_goodRep}
   % \uses{thm:good-rep-iff-split, thm:split-iff-mem-relative}
 \end{corollary}
 \begin{proof}
   Rephrase of~\cref{thm:good-rep-iff-split} and~\cref{thm:split-iff-mem-relative}.
 \end{proof}

\subsection{Basic Properties}
We observe the following easy result about good representations. Let $X \in \br(F)$ and $A$ be a good representation of $X$.

\begin{lemma}
  The range $\iota_{A}(A)$ is a simple ring.
  \leanok
  \lean{GoodRep.ιRange}
\end{lemma}

\begin{proof}
  Because $K$ is a simple ring, $\iota_{A}$ is injective therefore $\iota_{A}(A)\cong K$.
\end{proof}

\begin{lemma}\label{lem:centralizer-range-good-rep}
  $C_{A}\left(\iota_{A}(A)\right) = \iota_{A}(A)$.
  \leanok
  \lean{GoodRep.centralizerιRange}
  % \uses{lem:tfae-subfield}
\end{lemma}
\begin{proof}
  In the language of~\cref{sec:subfield}, $\iota_{A}(A)$ is a subfield of $A$, hence by~\cref{lem:tfae-subfield}, we only need to show $\dim_{F}A = {\left(\dim_{F}\iota_{A}(A)\right)}^{2}$. But $\dim_{F}A={\left(\dim_{F}K\right)}^{2}$ and $\iota(A)\cong K$.
\end{proof}

\begin{construction}\label{con:good-rep-mod}
  We give $A$ a $K$-module structure by {\em left\/} multiplication, that is for any $c \in K$ and $a \in A$, we define $c\cdot a$ to be $\iota_{A}(c)a$. Note that if $c \in F$ then $\iota_{A}(c)a = c \cdot a$, in another word, the $K$-action and the $F$-action on $A$ are compatible. Then $A$ is a finite dimensional $K$-vector space and $\dim_{K}A=\dim_{F}K$: indeed $\dim_{F}K\cdot\dim_{K}A = \dim_{F}K\cdot \dim_{F}K = \dim_{F}A$.
  \leanok
  \lean{GoodRep.dim_eq'}
\end{construction}

\begin{lemma}\label{lem:good-rep-iso}
  If $A$ and $B$ are two good representations of $X$, then $A \cong B$ as $F$-algebras.
  \leanok
  \lean{GoodRep.exists_iso}
  % \uses{lem:common-div-alg}
\end{lemma}
\begin{proof}
  By~\cref{lem:common-div-alg}, we find a division $F$-algebra $D$ and non-zero natural numbers $m, n$ such that $A \cong \mat_{m}(D)$ and $B \cong \mat_{n}(D)$ as $F$-algebras. Therefore
  \[
    \begin{aligned}
      {\left(\dim_{F}K\right)}^{2} &= \dim_{F}A &= m^{2}\dim_{F}D \\
      &= \dim_{F}B &= n^{2}\dim_{F}D.
    \end{aligned}
  \]
  Therefore $m = n$ and $A \cong \mat_{m}D =\mat_{n}D\cong B$.
\end{proof}

\subsection{Conjugation Factors and Conjugation Sequences}
In this section, let $K/F$ be a field extension, $X \in \br(F)$ and $A$ be a $K$-good representation of $X$.
\begin{remark}
  Since $\gal(K/F)$ acts on $K^{\star}$, for $x\in K^{\star}$,we feel free to write $\sigma\cdot x$ when it feels more readable than $\sigma(x)$, for example when there are nested brackets.
\end{remark}


\begin{definition}[Conjugation Factor]
  \label{def:conj-factor}
  With respect to $A$, a conjugation factor of $\sigma$ is a unit $x_{\sigma} \in A^{\star}$ such that for all $c \in K$,
  \[
    x_{\sigma} \iota_{A}(c)x_{\sigma}^{-1} = \iota_{A}(\sigma \cdot c).
  \]

  A conjugation sequence is a sequence $x:\gal(K/F) \to A^{\star}$ such that for all $\sigma\in\gal(K/F)$, $x_{\sigma}$ is a conjugation factor of $\sigma$. When we want to stress $A$, we say $A$-conjugation factor and $A$-conjugation sequence.
  \leanok
  \lean{GoodRep.conjFactor}
\end{definition}

\begin{remark}\label{rem:conj-factor-alternative-eq}
  When $x_{\sigma}$ is a conjugation factor of $\sigma$, the equalities $x_{\sigma}\iota_{A}(c) = x_{\sigma}\iota_{A}(\sigma(c))$ and $\iota_{A}(c)x_{\sigma}^{-1}=x_{\sigma}^{-1}\iota_{A}(\sigma(c))$ are also useful.
\end{remark}

\begin{construction}\label{con:exists-conj-seq}
  $A$ has a conjugation sequence:
  let $\sigma \in \gal(K/F)$, we have two $F$-algebra homomorphisms $K \to A$ given by $\iota_{A}$ and $\iota_{A}\circ \sigma$. Applying~\cref{thm:skolem-noether} to $\iota_{A}$ and $\iota_{A}\circ \sigma$ gives us the desired conjugation factor.
  \leanok
  \lean{GoodRep.arbitraryConjFactor}
  % \uses{thm:skolem-noether}
\end{construction}

\begin{construction}
  If $x$ is a conjugation factor of $\sigma$ and $y$ of $\tau$, then $xy$ is a conjugation factor of $\sigma\tau$. For any $c \in K$
  \[
    \iota_{A}(\sigma \cdot\tau(c)) = x\iota_{A}(\tau\cdot c)x^{-1} = xy\iota_{A}(c)y^{-1}x^{-1}=\left(xy\right)\iota_{A}{\left(xy\right)}^{-1}.
  \]
  \leanok
  \lean{GoodRep.mul'}
\end{construction}

\begin{theorem}\label{thm:conj-seq-basis}
  If $x$ is an $A$-conjugation sequence, then $\{x_{\sigma}|\sigma\in \gal(K/F)\}$ is an $K$-linearly independent set. When $K/F$ is finite dimensional and Galois, $\{x_{\sigma}|\sigma\in \gal(K/F)\}$ is a $K$-basis for $A$.
  \leanok
  \lean{GoodRep.conjFactor_linearIndependent, GoodRep.conjFactorBasis}
  % \uses{con:good-rep-mod, def:conj-factor}
\end{theorem}

\begin{proof}
  Suppose $\{x_{\sigma}\}$ is linearly dependent. Let $J \subseteq \gal(K/F)$ be such that $\{x_{\sigma}|\sigma\in J\}$ is a maximally linearly independent subset.
  Then $J \ne \gal(K/F)$, let $\sigma \in\gal(K/F)$ be an arbitrary automorphism that is not in $J$. Since $\{x_{\tau}|\tau\in J\}$ is maximally linearly independent, $x_\sigma \in \left\langle x_{\tau}|\tau\in J\right\rangle$. Hence, by~\cref{con:good-rep-mod} we have
  \[
    x_\sigma = \sum_{\tau\in J'}\lambda_{\tau}\cdot x_{\tau} = \sum_{\tau\in J'}\iota_{A}\left(\lambda_{\tau}\right)x_{\tau},
  \]
  for some non-zero $\lambda_{\tau}\in K$ and $J' \subseteq J$.
  For each $c \in K$, we have the following equality
  \[
    \begin{aligned}
      \iota_{A}\left(\sigma \cdot c\right) x_{\sigma}
      &= x_{\sigma} \iota_{A}(c) &\text{by~\cref{def:conj-factor}}\\
      &= \sum_{\tau\in J'}\lambda_{\tau}\cdot x_{\tau}\iota_{A}(c) & \\
      &= \sum_{\tau\in J'}\lambda_{\tau}\cdot \iota_{A}\left(\tau\cdot c\right)x_{\tau} &\text{by~\cref{def:conj-factor} again}\\
      &=\sum_{\tau\in J'}\iota_{A}\left(\lambda_{\tau}\tau(c)\right)x_{\tau};\\
      \iota_{A}\left(\sigma\cdot c\right)x_{\sigma}
      &= \sum_{\tau\in J'}\iota_{A}\left(\lambda_{\tau}\right) x_{\tau}\iota_{A}(c)\\
      &= \sum_{\tau\in J'}\iota_{A}\left(\sigma(c)\lambda_{\tau}\right)x_{\tau}.
    \end{aligned}
  \]
  Since $\{x_{\tau}|\tau\in J'\}$ is linearly independent, we have that for each $\tau \in J'$, $\lambda_{\tau}\tau(c) = \sigma(c)\lambda_{\tau} = \lambda_{\tau}\sigma(c)$. Note that $J'$ is not empty, for otherwise $x_{\sigma}=\sum_{\tau\in\emptyset}\lambda_{\tau}\cdot x_{\tau}=0$ but $x_{\sigma}$ is invertible. Since for any $\tau\in J'$, $\lambda_{\tau}$ is not zero, we have that for all $c \in K$, $\sigma(c) = \tau(c)$, i.e. $\sigma = \tau$. Hence $\sigma$ is in $J' \subseteq J$ after all; contradiction.

  If $K/F$ is finite dimensional and Galois, then $\dim_{F}K$ is equal to the cardinality of $\gal(K/F)$, then by the linear independence of $\{x_{\sigma}|\sigma\in\gal(K/F)\}$, we conclude that it is indeed a $K$-basis for $A$.
\end{proof}

\section{The Second Galois Cohomology}

In this section, we construct a group isomorphism between $\br(K/F) \cong \HH^{2}\left(\gal(K/F), K^{\star}\right)$ where $K/F$ is a finite dimensional Galois extension. To keep alignment of the Brauer group, let us use the multiplicative notation for group cohomology. Recall:

\begin{definition}[the Second Group Cohomology]
  \label{def:second-group-coh}
  Let $G$ be a group and $M$ an abelian group (written multiplicatively) with a $G$-action.

  A function $f : G \times G \to M$ is a {\em 2-cocycle\/} if for all $g,h,j\in G$,
  \[
    f(gh, j)f(g, h) = \left(g\cdot f(h, j)\right) f(g, hj).
  \]
  We denote the subgroup of 2-cocycles as $\mathcal{Z}^{2}(G, M)$.

  A function $f : G\times G \to M$ is a {\em 2-coboundary\/} if there exists an $x : G \to M$ such that for all $g, h \in G$
  \[
    \frac{g\cdot x(h)}{x(gh)} x(g) = f(g, h).
  \]
  We denote the subgroup of 2-coboundaries as $\mathcal{B}^{2}(G, M)$.

  The second group cohomology $\hh^{2}\left(G, M\right)$ is defined to be the quotient group of 2-cocycles modulo 2-coboundaries $^{\mathcal{Z}^{2}(G, M)}/_{\mathcal{B}^{2}(G, M)}$. If $s, t \in \mathcal{Z}^{2}(G, M)$, we say $s$ and $t$ are cohomologous if their equivalence class $[s], [t] \in \HH^{2}(G, M)$ are the same; in another word $st^{-1}\in\mathcal{B}^{2}(G, M)$.
  \leanok
  \lean{groupCohomology.H2, groupCohomology.IsMulTwoCoboundary, groupCohomology.IsMulTwoCocycle}
\end{definition}

\begin{lemma}\label{lem:2-cocycle-one-one}
  If $f \in \mathcal{B}^{2}(G, M)$ is a 2-cocycle and $x\in G$, we have
  \[
    \begin{aligned}
    f(1, x) &= f(1, 1) \\
    f(x, 1) &= x\cdot f(1, 1).
    \end{aligned}
  \]
\end{lemma}
\begin{proof}
  Indeed:
  \[
    \begin{aligned}
      f(1\cdot 1, x)f(1, 1) &= (1\cdot f(1,x))f(1, 1\cdot x)  \\
      f(1, x)f(1, 1) &= f(1, x) f(1,x) \\
      f(1, x) &= f(1,1)
    \end{aligned}
  \] and
  \[
    \begin{aligned}
      f(x\cdot 1, 1)f(x, 1) &= (x\cdot f(1, 1))f(x, 1\cdot 1)\\
      f(x, 1)f(x, 1) &= (x\cdot f(1, 1))f(x, 1)\\
      f(x, 1) &= x \cdot f(1, 1).
    \end{aligned}
  \]
  \leanok
  \lean{CrossProductAlgebra.a_one_left, CrossProductAlgebra.a_one_right}
\end{proof}

In the following sections of this chapter, we assume that $X \in\br(F)$ and $A$ is a good representation of $X$. We use $\rho, \sigma, \tau$ to denote elements of $\gal(K/F)$. To improve typographic aesthetics of our proofs, we sometimes use subscript to mean function application.


\subsection{From $\br(K/F)$ to $\hh^{2}\left(\gal(K/F),K^{\star}\right)$}

\begin{lemma}[Twisting Conjugation Factors]\label{lem:twist-spec1}
  If $x$ and $y$ are two conjugation factors of $\sigma$, then there exists a unique $c \in K$ such that $x = y\iota_{A}(c)$.
  \leanok
  \lean{GoodRep.conjFactor_rel, GoodRep.conjFactorTwistCoeff}
  % \uses{rem:conj-factor-alternative-eq, lem:centralizer-range-good-rep}
\end{lemma}

\begin{proof}
  The uniqueness is clear: suppose $x = y\iota_{A}(c)=y\iota_{A}(c')$, then $c = c'$ because $x, y$ are units and $\iota_{A}$ is injective. We first observe that $y^{-1}x\in C_{A}(\iota(A))$: for any $z \in K$, $y^{-1}x\iota_{A}(z)=y^{-1}\iota_{A}(\sigma(z))x=\iota_{A}(z)y^{-1}x$ (by~\cref{rem:conj-factor-alternative-eq}). By~\cref{lem:centralizer-range-good-rep}, $y^{-1}x\in \iota(A)$, that is for some $z \in K$, we have that $y^{-1}x = \iota_{A}(z)$ and the claim is proved.
\end{proof}

We denote such $c$ by $\twist^{\sigma}(x, y)$ or $\twist^{\sigma}_{x, y}$, when $\sigma$ is clear from context, we often omit the superscript. With this notation, $x = y\iota_{A}(\twist_{x,y})$.

\begin{remark}
  $\twist(x, x)$ is equal to $1$ by uniqueness.
\end{remark}

\begin{remark}
  In fact, $\twist(x, y)$ is in $K^{\star}$ and $\twist(x, y)^{-1}=\twist(y, x)$.
\end{remark}

\begin{lemma}\label{lem:twist-spec2}
  If $x$ and $y$ are conjugation factors for $\sigma$, $x=\iota_{A}(\sigma(\twist_{x, y}))y$.
  \leanok
  \lean{GoodRep.conjFactorTwistCoeff_spec'}
\end{lemma}
\begin{proof}
  \[
    \begin{aligned}
      x &= x \iota_{A}\left(\twist_{x,y}\right) x^{-1} x \iota_{A}(\twist_{y,x}) \\
        &= \iota_{A}\left(\sigma\cdot\twist_{x,y}\right)x\iota_{A}\left(
          \twist_{y,x}\right)\\
      &=\iota_{A}\left(\sigma\cdot \twist_{x,y}\right)y
    \end{aligned}.
  \]
\end{proof}

\begin{construction}[Comparing Conjugation Factors]\label{con:compare-conj-factors}
  Let $x$ be a conjugation factor for $\sigma$, $y$ for $\tau$ and $z$ for $\sigma\tau$. Since $xy$ is a also a conjugation factor, we define the comparison coefficient to be $\comp^{\sigma, \tau}_{x,y,z}:=\sigma\left(\tau\left(\twist_{xy, z}\right)\right)$. We often omit superscript when the context is clear. Note that $\comp_{x,y,z}$ is a unit in $K$ with inverse $\sigma\left(\tau\left(\twist_{z, xy}\right)\right)$. By~\cref{lem:twist-spec1} and~\cref{lem:twist-spec2}, we have the following useful equalities
  \[
    \begin{aligned}
      xy &= \iota_{A}\left(\comp_{x,y,z}\right)z \\
      \iota_{A}\left(\comp_{x,y,z}^{-1}\right)xy &= z \\
      \iota_{A}\left(\comp_{x,y,z}\right) &= xyz^{-1}\\
      \iota_{A}\left(\comp_{x,y,z}^{-1}\right) &= zy^{-1}x^{-1}\\
      \dots &=\dots
    \end{aligned}.
  \]
  %If we have a conjugation sequence $x : \gal(K/F)\to A^{\star}$, instead of $\comp_{x_{\sigma},x_{\tau},x_{\sigma\tau}}^{\sigma,\tau}$, we simply write $\comp_{\sigma,\tau}^{x}$.
  \leanok
  \lean{GoodRep.conjFactorCompCoeff}
  % \uses{lem:twist-spec1, lem:twist-spec2}
\end{construction}

\begin{lemma}\label{lem:compare-conj-factor-spec}
  Let $x : \gal(K/F) \to A^{\star}$ be a conjugation sequence. We have
  \[
    \comp_{x_{\rho}, x_{\sigma}, x_{\rho\sigma}}\,\comp_{x_{\rho\sigma},x_{\tau},x_{\rho\sigma\tau}} =
    \rho\left(\comp_{x_{\sigma},x_{\tau},x_{\sigma\tau}}\right)\,\comp_{x_{\rho}, x_{\sigma\tau}, x_{\rho\sigma\tau}}.
  \]
  \leanok
  \lean{GoodRep.conjFactorCompCoeff_comp_comp}
  % \uses{con:compare-conj-factors, def:conj-factor, con:br-base-change}
\end{lemma}
\begin{proof}
  It is sufficient to make the following calculations:
  \begin{align}
   x_{\rho}\,x_{\sigma}x_{\tau} &= \iota_{A}\left(\comp_{x_{\rho},x_{\sigma},x_{\rho\sigma}}\right) \, \iota_{A}\left(\comp_{x_{\rho\sigma}, x_{\tau}, x_{\rho\sigma\tau}}\right)\, x_{\rho\sigma\tau}\label{eq:1}\\
   x_{\rho}\left(x_{\sigma}x_{\tau}\right) &= \iota_{A}\left(\rho\cdot\comp_{x_{\sigma},x_{\tau},x_{\sigma\tau}}\right)\,\iota_{A}\left(\comp_{x_{\rho},x_{\sigma\tau},x_{\rho\sigma\tau}}\right) x_{\rho\sigma\tau} \label{eq:2}
 \end{align}
 Then since $x_{\rho\sigma\tau}$ is invertible and $\iota_{A}$ is injective, we proved the desired result.

 \Cref{eq:1} is because: by the first equality in~\cref{con:compare-conj-factors} (twice)
 \[
   x_{\rho}x_{\sigma}x_{\tau}=\iota_{A}\left(\comp_{x_{\rho},x_{\sigma}, x_{\rho\sigma}}\right)x_{\rho\sigma}x_{\tau}=\iota_{A}\left(\comp_{x_{\rho},x_{\sigma},x_{\rho\sigma}}\right)\iota_{A}\left(\comp_{x_{\rho\sigma},x_{\tau},x_{\rho\sigma\tau}}\right)x_{\rho\sigma\tau}.
 \]
 \Cref{eq:2} is because: by~\cref{def:conj-factor}, we have
 \[\iota_{A}\left(\rho\cdot\comp_{x_{\sigma},x_{\tau},x_{\sigma\tau}}\right)x_{\rho}=x_{\rho}\iota_{A}\left(\comp_{x_{\sigma},x_{\tau},x_{\sigma\tau}}\right),\]
 therefore by~\cref{con:br-base-change}
 \[
   \begin{aligned}
     x_{\rho}\left(x_{\sigma}x_{\tau}\right) &= x_{\rho}\iota_{A}\left(\comp_{x_{\sigma},x_{\tau},x_{\sigma\tau}}\right)x_{\sigma\tau} \\
                                             &= \iota_{A}\left(\rho\cdot \comp_{x_{\sigma},x_{\tau},x_{\sigma\tau}}\right)x_{\rho}x_{\sigma\tau}\\
     &= \iota_{A}\left(\rho\cdot\comp_{x_{\sigma},x_{\tau},x_{\sigma\tau}}\right)\,\iota_{A}\left(\comp_{x_{\rho},x_{\sigma\tau},x_{\rho\sigma\tau}}\right)x_{\rho\sigma\tau}
   \end{aligned}.
 \]
\end{proof}

\begin{construction}[from good representation to 2-cocycle]
  \label{con:good-rep-to-2-cocycles}
  Let $x$ be an $A$-conjugation sequence. We associate with $x$ a function $\mathcal{B}^{2}(x) : \gal(K/F)\times\gal(K/F)\to K^{\star}$ defined by
  \[
    (\sigma,\tau)\mapsto \comp_{x_{\sigma},x_{\tau},x_{\sigma\tau}}.
  \]
  We will write $\mathcal{B}^{2}(x)$ as $\mathcal{B}^{2}_{A, x}$, $\mathcal{B}^{2}_{A}(x)$ or $\mathcal{B}^{2}_{x}$ as well.
  \leanok
  \lean{GoodRep.toTwoCocycles}
\end{construction}

\begin{lemma}
  For any $A$-conjugation sequence $x$, $\mathcal{B}^{2}_{x} \in \mathcal{B}^{2}\left(\gal(K/F), K^{\star}\right)$, that is $\mathcal{B}_{x}$ is indeed a 2-cocycle.
  \leanok
  \lean{GoodRep.isTwoCocyles}
  % \uses{lem:compare-conj-factor-spec}
\end{lemma}

\begin{proof}
  We need to prove
  \[
    \mathcal{B}_{x}(\rho\sigma, \tau)\, \mathcal{B}_{x}(\rho,\sigma) = \rho\left(\mathcal{B}_{x}(\sigma, \tau)\right)\, \mathcal{B}_{x}(\rho, \sigma\tau).
  \]
  But this is exactly~\cref{lem:compare-conj-factor-spec}.
\end{proof}

For any good representation $A$ of $X \in \br(K/F)$ and any $A$-conjugation sequence $x$, we have constructed a 2-cocycle $\mathcal{B}^{2}_{A}(x)$. But to obtain a well-defined function from $\br(K/F)$ to $\HH^{2}\left(\gal(K/F), K^{\star}\right)$, we need to verify that for any other good representation $B$ of $X$ and $B$-conjugation sequence $y$, $\mathcal{B}^{2}_{A}(x)$ and $\mathcal{B}^{2}_{B}(y)$ are cohomologous. Let us fix another good representation $B$ of $X \in \br(K/F)$ and a $B$-conjugation sequence $y$.

\begin{construction}\label{con:good-rep-iso-coeff}
  By~\cref{lem:good-rep-iso}, $A$ and $B$ are isomorphic as $F$-algebras, we use $e_{A,B}$ to denote an arbitrary $F$-algebra isomorphism between $A$ and $B$. When there is no confusion, we write $e$ instead of $e_{A, B}$ Since $e \circ \iota_{A}$ and $\iota_{B}$ are two $F$-algebra homomorphism from $K$ to $B$, by~\cref{thm:skolem-noether}, there exists some $u \in B^{\star}$ such that for all $r \in K$, we have $\iota_{B}(r) = ue\left(\iota_{A}(r)\right)u^{-1}$ (or equivalently, $u^{-1}\iota_{B}(r)u = e\left(\iota_{A}(r)\right)$). When there is confusion, we write $u_{A, B}$ instead of $u$.
  % \uses{lem:good-rep-iso}
\end{construction}

\begin{lemma}\label{lem:iso-conj-coeff-spec2}
  For any $c \in K$, $\sigma\in\gal(K/F)$ and $A$-conjugation factor $x$ of $\sigma$, we have
  \[
    \iota_{B}\left(\sigma\cdot c\right) = ue(x)u^{-1}\,\iota_{B}(c)\,ue\left(x^{-1}\right)u^{-1}.
  \]
  \leanok
  \lean{GoodRep.isoConjCoeff_spec''}
  % \uses{def:conj-factor, con:good-rep-iso-coeff}
\end{lemma}
\begin{proof}
  From~\cref{def:conj-factor}, we have $e\left(\iota_{A}(\sigma \cdot c)\right) = e\left(x\iota_{A}(c)x^{-1}\right)$. Substituting it in~\cref{con:good-rep-iso-coeff}, we get
  \[
    \begin{aligned}
      \iota_{B}(\sigma \cdot c) &= ue\left(x\iota_{A}(c)x^{-1}\right)u^{-1}\\
                                &= ue(x)e\left(\iota_{A}(c)\right)e\left(x^{-1}\right)u^{-1}\\
      &= ue(x)u^{-1}\iota_{B}(c)ue\left(x^{-1}\right)u^{-1}.
    \end{aligned}
  \]
\end{proof}

\begin{construction}
  \label{con:push-forward-conj-factor}
  If $x$ is an $A$-conjugation factor for $\sigma$, we can obtain a $B$-conjugation factor for $\sigma$ by defining $B_{\star}x := ue(x)u^{-1}$ with inverse $ue\left(x^{-1}\right)u^{-1}$. We use~\cref{lem:iso-conj-coeff-spec2} to check that $B_{\star}x$ is indeed a conjugation factor for $\sigma$. If $y$ is a $B$-conjugation factor for $\sigma$, another useful constant is $v := \sigma\left(\twist_{y, B_{\star}x}\right)$. We have
  \[
    \begin{aligned}
    y &= \iota_{B}\left(v\right)\, B_{\star}x \\
      \iota_{B}\left(v\right) &= u e\left(\iota_{A}\left(v\right)\right)u^{-1}\\
      v^{-1} &= \sigma\left(\twist_{B_{\star}x, y}\right)
    \end{aligned}.
  \]
  We also write $v_{x,y}$ or even $v^{A,B}_{x, y}$ when we stress the importance of good representation $A$ and $B$ and their conjugation factor $x$ and $y$.
  % \uses{lem:iso-conj-coeff-spec2}
\end{construction}

\begin{lemma}\label{lem:compare-to-2-cocycle-aux}
  Let $x$ be an $A$-conjugation sequence and $y$ a $B$-conjugation sequence.
  We have
  \[
    \comp_{y_{\sigma}, y_{\tau}, y_{\sigma\tau}} v_{x_{\sigma\tau},y_{\sigma\tau}} = v_{x_{\sigma},y_{\sigma}} \sigma\left(v_{x_{\tau},y_{\tau}}\right) \comp_{x_{\sigma},x_{\tau},x_{\sigma\tau}}.
  \]
  \leanok
  \lean{GoodRep.compare_conjFactorCompCoeff}
  % \uses{con:compare-conj-factors, con:push-forward-conj-factor, con:good-rep-iso-coeff, def:conj-factor}
\end{lemma}
\begin{proof}
  By~\cref{con:compare-conj-factors}, we have $y_{\sigma}y_{\tau} = \iota_{B}\left(\comp_{y_{\sigma},y_{\tau},y_{\sigma\tau}}\right)y_{\sigma\tau}$.
  By repeated application of~\cref{con:push-forward-conj-factor} and~\cref{con:good-rep-iso-coeff}, we have
  \[
    \begin{aligned}
      y_{\sigma\tau}&= u\,e\left(\iota_{A}\left(v_{x_{\sigma\tau},y_{\sigma\tau}}\right)x_{\sigma\tau}\right)\,u^{-1}\\
      y_{\sigma}y_{\tau} &= u\,e\left(\iota_{A}\left(v_{x_{\sigma},y_{\sigma}}\right)x_{\sigma}\iota_{A}\left(v_{x_{\tau},y_{\tau}}\right)x_{\tau}\right)\,u^{-1}\\
                    &= \iota_{B}\left(\comp_{y_{\sigma},y_{\tau},y_{\sigma\tau}}\right)y_{\sigma\tau}\\
                    &= \iota_{B}\left(\comp_{y_{\sigma},y_{\tau},y_{\sigma\tau}}\right)u\,e\left(\iota_{A}\left(v_{x_{\sigma\tau},y_{\sigma\tau}}\right)x_{\sigma\tau}\right)\,u^{-1}\\
                    &= u\,e\left(\iota_{A}\left(\comp_{y_{\sigma},y_{\tau},y_{\sigma\tau}}\right)\right)\,u^{-1}\,\,\,ue\left(\iota_{A}\left(v_{x_{\sigma\tau},y_{\sigma\tau}}\right)x_{\sigma\tau}\right)\,u^{-1}\\
      &= ue\left(\iota_{A}\left(\comp_{y_{\sigma},y_{\tau},y_{\sigma\tau}}v_{x_{\sigma\tau},y_{\sigma\tau}}\right)x_{\sigma\tau}\right)u^{-1}.
    \end{aligned}
  \]
  Hence
  \[
    \iota_{A}\left(v_{x_{\sigma},y_{\sigma}}\right)x_{\sigma}\iota_{A}\left(v_{x_{\tau},y_{\tau}}\right)x_{\tau} = \iota_{A}\left(\comp_{y_{\sigma},y_{\tau},y_{\sigma\tau}}v_{x_{\sigma\tau},y_{\sigma\tau}}\right)x_{\sigma\tau}.
  \]
  We also have by~\cref{def:conj-factor}
  \[
    x_{\sigma} \iota_{A}\left(v_{x_{\tau},y_{\tau}}\right) x_{\tau} =
    \iota_{A}\left(\sigma\cdot v_{x_{\tau}y_{\tau}}\right) x_{\sigma}x_{\tau}.
  \]
  Hence
  \[
    \begin{aligned}    \iota_{A}\left(v_{x_{\sigma},y_{\sigma}}\right)x_{\sigma}\iota_{A}\left(v_{x_{\tau},y_{\tau}}\right)x_{\tau} &=
                                                                                                                                      \iota_{A}\left(v_{x_{\sigma},y_{\sigma}}\, \sigma\left(v_{x_{\tau},y_{\tau}}\right)\right) x_{\sigma}x_{\tau} \\
      &= \iota_{A}\left(v_{x_{\sigma},y_{\sigma}}\sigma\left(v_{x_{\tau},y_{\tau}}\right)\comp_{x_{\sigma},x_{\tau},x_{\sigma\tau}}\right)x_{\sigma\tau}\\
      &= \iota_{A}\left(\comp_{y_{\sigma},y_{\tau},y_{\sigma\tau}}v_{x_{\sigma\tau},y_{\sigma\tau}}\right)x_{\sigma\tau}.
    \end{aligned}
  \]
  Cancelling $x_{\sigma\tau}$ and by injectivity of $\iota_{A}$, the result is proved.
\end{proof}

\begin{lemma}\label{lem:compare-to-2-cocycle}
  Let $x$ be an $A$-conjugation sequence and $y$ a $B$-conjugation sequence. We have
  \[
    \mathcal{B}^{2}_{B,y}(\sigma,\tau) v_{x_{\sigma\tau},y_{\sigma\tau}} =
    v_{x_{\sigma},y_{\sigma}} \sigma\left(v_{x_{\tau},y_{\tau}}\right) \mathcal{B}^{2}_{A,y}(\sigma,\tau).
  \]
  \leanok
  \lean{GoodRep.compare_toTwoCocycles}
  % \uses{con:good-rep-to-2-cocycles,lem:compare-to-2-cocycle-aux}
\end{lemma}

\begin{proof}
  If we unfold~\cref{con:good-rep-to-2-cocycles}, we discover the lemma is saying exactly~\cref{lem:compare-to-2-cocycle-aux}.
\end{proof}

We finally arrive at our main conclusion for this section.
\begin{corollary}\label{cor:to-2-cocycle-wd}
  Let $x$ be an $A$-conjugation sequence and $y$ a $B$-conjugation sequence. $\mathcal{B}^{2}_{A,x}$ and $\mathcal{B}^{2}_{B,y}$ are 2-cohomologous.
  \leanok
  \lean{GoodRep.compare_toTwoCocycles'}
  % \uses{def:second-group-coh, lem:compare-to-2-cocycle}
\end{corollary}
\begin{proof}
  By~\cref{def:second-group-coh}, we need to find a function $f : \gal(K/F)\to K^{\star}$ such that for all $\sigma, \tau\in\gal(K/F)$,
  \[
    \frac{\sigma\left(f(\tau)\right)}{f\left(\sigma\tau\right)}f(\sigma)=\frac{\mathcal{B}^{2}_{B,y}}{\mathcal{B}^{2}_{A,x}}.
  \]
  Let $f(\rho) := v_{x_{\rho},y_{\rho}}$, by~\cref{lem:compare-to-2-cocycle} we see the equality holds.
\end{proof}

\begin{construction}[from $\br(K/F)$ to $\HH^{2}\left(\gal(K/F), K^{\star}\right)$]\label{con:br-to-snd-coh}
  Let $X \in \br(K/F)$, by~\cref{cor:mem-relative-br-iff-good-rep}, $X$ admits a good representation $A$; by~\cref{con:exists-conj-seq}, $A$ admits a conjugation sequence $x$. We associate with $X$ an element $\HH^{2}(X) := \left[\mathcal{B}^{2}_{A,x}\right]$in $\HH^{2}\left(\gal(K/F),K^{\star}\right)$. By~\cref{cor:to-2-cocycle-wd}, for any other good representation $B$ and $B$-conjugation sequence $y$, we have $\left[\mathcal{B}^{2}_{A,x}\right]=\left[\mathcal{B}^{2}_{B,y}\right]$, hence we have a well-defined function $\HH^{2} : \br(K/F) \to \HH^{2}\left(\gal(K/F),K^{\star}\right)$.
  \leanok
  \lean{RelativeBrGroup.toSnd, RelativeBrGroup.toSnd_wd}
  % \uses{con:exists-conj-seq, cor:mem-relative-br-iff-good-rep,cor:to-2-cocycle-wd}
\end{construction}

\subsection{Cross Product as a Central Simple Algebra}

Let $\mathfrak{a} \in \mathcal{B}^{2}\left(\gal(K/F), K^{\star}\right)$ be any 2-cocycle. In this section, we construct the {\em cross product\/} associated with $\mathfrak{a}$ which we prove to be $F$-central simple. Finally, we show that if $\mathfrak{a}, \mathfrak{b}\in\mathcal{B}^{2}\left(\gal(K/F), K^{\star}\right)$ are cohomologous, the cross products associated with $a$ and $b$ are Brauer equivalent.

\begin{construction}[Cross product]\label{con:cross-product}
  Denote $\cross_{\mathfrak{a}}$ to be $\gal(K/F) \to K$, i.e. functions from $\gal(K/F)$ to $K$. Notationally, elements of $\cross_{\mathfrak{a}}$ are sequences in $K$ indexed by $\gal(K/F)$; we denote $\Delta^{\mfa}_{\sigma, c}$ to be the sequence with value $c$ at $\sigma$-th index and zero elsewhere. When $\mfa$ is clear from context, we will omit the superscript. We give $\cross_{\mathfrak{a}}$ the usual zero, addition, negation, that is, we give $\cross_{\mathfrak{a}}$ the normal additive abelian group structure. Since for each $c \in \cross_{\mathfrak{a}}$,
  \[
    c = \sum_{\sigma\in\gal(K/F)}\Delta_{\sigma, c(\sigma)},
    \]
    it is often, if not always, sufficient to consider the special cases of $\Delta_{\sigma, c}$ and extend the result linearly.
  For multiplications, we define the result of multiplying $\Delta_{\sigma, c},\Delta_{\tau, d} \cross_{\mathfrak{a}}$ to be $\Delta_{\sigma\tau, c\sigma(d)\, \mathfrak{a}(\sigma,\tau)}$. Immediately, if either $c$ or $d$ is $0$, the result of multiplication is also zero. That is, for all $c \in \cross_{\mfa}$, we have $c\cdot 0=0\cdot c =0$. For any $r \in F$ and $\Delta_{\sigma, c} \in \cross_{\mathfrak{a}}$, we define $r \cdot \Delta_{\sigma, c}$ to be $\Delta_{\sigma, r\cdot c}$.
  \leanok
  \lean{CrossProductAlgebraMul, CrossProductAlgebraSMul, CrossProductAlgebra.addCommGroup}
\end{construction}

\begin{remark}
  When $K/F$ is infinite dimensional, the correct definition of $\cross_{\mathfrak{a}}$ is perhaps $\bigoplus_{\sigma\in\gal(K/F)}K$. But in Lean4, function type is easier to manipulate than direct sums. Since our scope is finite dimensional Galois extension, our definition is still accurate.
\end{remark}

\begin{lemma}
  The cross product $\cross_{\mathfrak{a}}$ is a ring with the multiplicative unit $\Delta_{\mathsf{id}, \mathfrak{a}(\mathsf{id}, \mathsf{id})^{-1}}$. The $F$-action on $\cross_{\mathfrak{a}}$ defined by $r \cdot \Delta_{\sigma, c} := \Delta_{\sigma, r\cdot c}$ makes it an $F$-algebra.
  \leanok
  \lean{CrossProductAlgebra.algebra}
  % \uses{def:second-group-coh, lem:2-cocycle-one-one}
\end{lemma}
\begin{proof}
  We verify the axioms of rings on elements of the form $\Delta_{\sigma, c}$. Let $\sigma,\tau,\rho\in\gal(K/F)$ and $a,b,c\in K$.
  \begin{itemize}
    \item associativity of multiplication. We need to check that $\Delta_{\sigma,a}\left(\Delta_{\tau,b}\Delta_{\rho,c}\right) = \left(\Delta_{\sigma,a}\Delta_{\tau,b}\right)\Delta_{\rho,c}$:
          \[
          \begin{aligned}
            \Delta_{\sigma, a}\left(\Delta_{\tau,b}\Delta_{\rho,c}\right)
            &= \Delta_{\sigma,a}\Delta_{\tau\rho,b\tau(c) \mathfrak{a}(\sigma,\tau)}\\
            &= \Delta_{\sigma\tau\rho,a\sigma(b)\sigma(\tau(c)) \sigma(\mathfrak{a}(\sigma, \tau))};\\
            %
            \left(\Delta_{\sigma,a}\Delta_{\tau,b}\right)\Delta_{\rho,c}
            &= \Delta_{\sigma\tau,a\sigma(b)\mathfrak{a}(\sigma,\tau)}\Delta_{\rho,c}\\
            &= \Delta_{\sigma\tau\rho,a\sigma(b)\mathfrak{a}(\sigma, \tau)\sigma(\tau(c))\mathfrak{a}(\sigma\tau, \rho)}.
          \end{aligned}
          \]
          Hence it is sufficient to check
          \[
          \sigma(\tau(c))\sigma(\mathfrak{a}(\sigma,\tau)) =
          \mathfrak{a}(\sigma,\tau)\sigma(\tau(c))\mathfrak{a}(\sigma\tau, \rho).
          \]
          This is the 2-cocycle condition in~\cref{def:second-group-coh} (modulo commutativity of $K$).
    \item multiplicative unit: we need to check $\Delta_{\sigma, a}\Delta_{\mathsf{id},\mathfrak{a}(\mathsf{id}, \mathsf{id})} = \Delta_{\mathsf{id},\mathfrak{a}(\mathsf{id},\mathsf{id})}\Delta_{\sigma, a} = \Delta_{\sigma, a}$. By multiple applications of~\cref{lem:2-cocycle-one-one}
          \[
          \begin{aligned}
            \Delta_{\mathsf{id},\mathfrak{a}(\mathsf{id}, \mathsf{id})^{-1}}\Delta_{\sigma, a}
            &= \Delta_{\sigma, \mathfrak{a}(\mathsf{id}, \mathsf{id})^{-1}a\mathfrak{a}(\mathsf{id},\sigma)}\\
            &= \Delta_{\sigma, \mathfrak{a}(\mathsf{id}, \mathsf{id}) a\mathfrak{a}(\mathsf{id},\mathsf{id})}\\
            &= \Delta_{\sigma, a} \\
            \Delta_{\sigma, a}\Delta_{\mathsf{id},\mathfrak{a}(\mathsf{id},\mathsf{id})^{-1}}
            &= \Delta_{\sigma, a\sigma(\mathfrak{a}(\id,\id))^{-1}\mathfrak{a}(\sigma,\id)}\\
            &=\Delta_{\sigma, a\sigma(\mfa(\id,\id))^{-1}\sigma(\mfa(\id, id))}\\
            &=\Delta_{\sigma, a}
          \end{aligned}
          \]
    \item distributivity: We need to check left-distributivity $\Delta_{\sigma, a}\left(\Delta_{\tau, b} + \Delta_{\rho, c}\right) = \Delta_{\sigma,a}\Delta_{\tau,b} + \Delta_{\sigma, a}\Delta_{\rho, c}$ and right distributivity $\left(\Delta_{\tau,b} + \Delta_{\rho,c}\right)\Delta_{\sigma,a} = \Delta_{\tau,b}\Delta_{\sigma,a} + \Delta_{\rho,c}\Delta_{\sigma,a}$. This is precisely what ``extend linearly'' means.
    \item $F$-algebra: We need to check for all $r \in F$, $\left(r\cdot \Delta_{\id,\mfa(\id,\id)^{-1}}\right)\Delta_{\sigma,c} = \Delta_{\sigma,c}\left(r\cdot\Delta_{\id,\mfa(\id,\id)^{-1}}\right)$. By~\cref{lem:2-cocycle-one-one}
          \[
          \begin{aligned}
            \left(r\cdot \Delta_{\id,\mfa(\id,\id)^{-1}}\right)\Delta_{\sigma,c}
            &= \Delta_{\id,r\cdot\mfa(\id,\id)^{-1}}\Delta_{\sigma,c}\\
            &= \Delta_{\sigma,\left(r\cdot\mfa(\id,\id)^{-1}\right)c\mfa(\id,\sigma)} \\
            &=\Delta_{\sigma,\left(r\cdot\mfa(\id,\id)^{-1}\right)c\mfa(\id,\id)} \\
            &= \Delta_{\sigma, r\cdot c}\\
            %
            \Delta_{\sigma,c}\left(r\cdot\Delta_{\id,\mfa(\id,\id)^{-1}}\right)
            &= \Delta_{\sigma,c}\Delta_{\id,r\cdot \mfa(\id,\id)^{-1}} \\
            &= \Delta_{\sigma,c\sigma\left(r\cdot\mfa(\id,\id)^{-1}\right)\mfa(\sigma,\id)}\\
            &=\Delta_{\sigma,c\left(r\cdot\sigma\left(\mfa(\id,\id)^{-1}\right)\right)\mfa(\id,\id)^{-1}}\\
            &= \Delta_{\sigma, c\left(r\cdot 1\right)}\\
            &=\Delta_{\sigma,r\cdot c}.
          \end{aligned}
          \]
  \end{itemize}
\end{proof}

From now on, we feel free to write $1 \in \cross_{\mfa}$ instead of $\Delta_{\id,\mfa(\id,\id)^{-1}}$. Then the algebra map $F \hookrightarrow \cross_{\mfa}$ is the map $r \mapsto r \cdot 1$.

\begin{construction}[$K$-embedding]\label{con:cross-product-iota}
  The map $\iota_{\cross_{\mfa}}: K \to \cross_{\mfa}$ defined by
  \[
    b \mapsto \Delta_{\id, b\mfa(\id,\id)^{-1}}
  \]
  is an $F$-algebra map.
  Checking that $\iota_{\cross_{\mfa}}$ preserves $1$, multiplication and addition uses nothing but axioms of ring.
  For any $r \in F$, we need to check $\iota_{\cross_{\mfa}}(r) = r \cdot 1$. Indeed $\iota_{\cross_{\mfa}}(r) = \Delta_{\id, r\cdot \mfa(\id,\id)^{-1}}$ and $r\cdot 1 = r \cdot \Delta_{\id,\mfa(\id,\id)^{-1}} =\Delta_{\id,r\cdot\mfa(\id,\id)^{-1}}$.
  When the context is clear, we also write $\iota_{\mfa}$ instead of $\iota_{\cross_{\mfa}}$. We give $\cross_{\mfa}$ a $K$-module structure by left-multiplication, that is for any $b \in K$ and $c\in\cross_{\mfa}$, we define $b\cdot c := \iota_{\mfa}(b)c$.


  We note the following useful equality: for any $b \in K$

  \[
    b\cdot \Delta_{\sigma,c} = \iota_{\mfa}(b)\Delta_{\sigma, c} = \Delta_{\sigma, bc},
  \]
  indeed:
  $\iota_{\mfa}(b)\Delta_{\sigma, c} = \Delta_{\id,b \mfa(\id,\id)^{-1}}\Delta_{\sigma, c} = \Delta_{\sigma, b\mfa(\id,\id)^{-1}c\mfa(\id,\sigma)} = \Delta_{\sigma,b\mfa(\id,\id)^{-1}c\mfa(\id,\id)} = \Delta_{\sigma, bc}$ by~\cref{lem:2-cocycle-one-one}. In another word, for any $b \in F$ and $c \in \cross_{\mfa}$, the $K$-action of $b$ on $c$ and the $F$-action of $b$ on $c$ agree.

  \leanok \lean{CrossProductAlgebra.ι,CrossProductAlgebra.identity_double_cross'}
  % \uses{lem:2-cocycle-one-one}
 \end{construction}

 \begin{lemma}
   \label{lem:cross-product-delta-invertible}
   For every $\sigma \in \gal(K/F)$, $\Delta_{\sigma,1}$ is invertible.
   \leanok
   \lean{CrossProductAlgebra.x_}
 \end{lemma}

 \begin{proof}
   It is sufficient to prove that $\Delta_{\sigma, 1}$ has a left inverse and right inverse.
   The left inverse of $\Delta_{\sigma, 1}$ is
   \[
     \Delta_{\sigma^{-1}, \mfa\left(\sigma^{-1},\sigma\right)^{-1}\mfa(\id,\id)^{-1}}.
   \]
   Indeed, for any $a \in K$, we have
   \[
     \Delta_{\sigma^{-1}, a}\,\Delta_{\sigma, 1} = \Delta_{\id, a\mfa\left(\sigma^{-1},\sigma\right)},
   \]
   hence substitute $a = \mfa\left(\sigma^{-1},\sigma\right)^{-1}\mfa(\id,\id)^{-1}$, we see the right hand side is $\Delta_{
     \id,\mfa(\id,\id)^{-1}}$ which is precisely $1 \in \cross_{\mfa}$. The right inverse is
   \[
     \Delta_{\sigma^{-1}, \sigma^{-1}\left(\mfa\left(\sigma,\sigma^{-1}\right)^{-1} \mfa(\id,\id)^{-1}\right)}.
   \]
   Indeed, for any $a \in K$, we have
   \[
     \Delta_{\sigma,1}\,\Delta_{\sigma^{-1},a} = \Delta_{\id, \sigma(a)\mfa\left(\sigma,\sigma^{-1}\right)},
   \]
   hence substitude $a = \sigma^{-1}\left(\mfa\left(\sigma,\sigma^{-1}\right)^{-1}\mfa(\id,\id)^{-1}\right)$, the right hand side is again $\Delta_{\id,\mfa(\id,\id)^{-1}}$ which is precisely $1 \in \cross_{a}$.
 \end{proof}
 \begin{lemma}\label{lem:cross-product-basis-conj}
   For any $c \in K$, we have
   \[
     \Delta_{\sigma,1}\,\iota_{\mfa}(c) = \iota_{\mfa}(\sigma\cdot c)\Delta_{\sigma, 1} = \Delta_{\sigma,\sigma\cdot c}
   \]
   and consequently,
   \[
     \Delta_{\sigma,1}\,\iota_{\mfa}(c)\,\Delta_{\sigma, 1}^{-1} = \iota_{\mfa}(\sigma \cdot c).
   \]
   \leanok
   \lean{CrossProductAlgebra.x__conj}
 \end{lemma}

 \begin{proof}
   We calculate
   \[
     \begin{aligned}
       \Delta_{\sigma,1}\iota_{\mfa}(c)
       &= \Delta_{\sigma,1}\,\Delta_{\id,c\mfa(\id,\id)^{-1}}\\
       &= \Delta_{\sigma, \sigma(c\mfa(\id,\id)^{-1})\mfa(\sigma,1)}\\
       &= \Delta_{\sigma,\sigma(c)\sigma(\mfa(\id,\id))^{-1}\sigma(\mfa(\id,\id))}\\
       &= \Delta_{\sigma,\sigma(c)}.
     \end{aligned}
   \]
 \end{proof}


 \begin{lemma}\label{lem:cross-product-basis-mul}
   We have $\Delta_{\sigma,1}\,\Delta_{\tau,1} = \iota_{\mfa}(\mfa(\sigma,\tau))\Delta_{\sigma\tau,1} = \mfa(\sigma,\tau)\cdot\Delta_{\sigma\tau,1}$
   Consequently we have for any $c, d \in K$,
   \[
     \Delta_{\sigma, c}\Delta_{\tau,d} = \left(c \sigma(d)\mfa(\sigma,\tau)\right) \cdot \Delta_{\sigma\tau, 1}.
   \]
   \leanok
   \lean{CrossProductAlgebra.x__mul,CrossProductAlgebra.mul_single_in_xAsBasis}
   % \uses{con:cross-product-iota,lem:cross-product-basis-conj}
 \end{lemma}
 \begin{proof}
   The first equality is in~\cref{con:cross-product-iota}.
   For the second equality, by~\cref{lem:cross-product-basis-conj}, we have
   \[
     \begin{aligned}
       \Delta_{\sigma,c}\Delta_{\sigma,d}
       &= \left(c\cdot\Delta_{\sigma, 1}\right)\left(d\cdot\Delta_{\tau,1}\right) \\
       &= \iota_{\mfa}(c)\left(\Delta_{\sigma, 1}\iota_{\mfa}(d)\right)\Delta_{\tau, 1} \\
       &= \iota_{\mfa}\Delta_{\sigma, \sigma \cdot d}\Delta_{\tau, 1}\\
       &= c \cdot \sigma(d) \cdot \Delta_{\sigma, 1}\Delta_{\tau, 1} \\
       &= c \cdot \sigma(d)\cdot\mfa(\sigma,\tau)\cdot\Delta_{\sigma\tau,1}.\\
     \end{aligned}
   \]
 \end{proof}


 \begin{lemma}
   \label{lem:cross-product-basis}
   The set $\{\Delta_{\sigma, 1}|\sigma\in\gal(K/F)\}$ forms a $K$-basis for $\cross_{\mfa}$.
   \leanok
   \lean{CrossProductAlgebra.x_AsBasis}
   % \uses{con:cross-product-iota}
 \end{lemma}
 \begin{proof}
   Suppose some linear combination $\sum_{\sigma}\lambda_{\sigma}\cdot\Delta_{\sigma, 1}$ is $0$ for some $\lambda_{\sigma}$'s in $K$. We have, by the equality in~\cref{con:cross-product-iota}
   \[
       \sum_{\sigma\in\gal(K/F)}\lambda_{\sigma}\cdot\Delta_{\sigma,1}
       = \sum_{\sigma\in\gal(K/F)}\Delta_{\sigma,\lambda_{\sigma}} = 0.\\
     \]
     Thus, for any $\tau \in \gal(K/F)$, we have
     \[
\left(\sum_{\sigma\in\gal(K/F)}\lambda_{\sigma}\cdot \Delta_{\sigma,1}\right)(\tau) = 0 = \lambda_{\tau},
\]
which proves linear independence. The fact that $\{\Delta_{\sigma, 1}|\sigma\in\gal(K/F)\}$ spans $\cross_{\mfa}$ is easy to see because every $\Delta_{\tau, a} = a \cdot \Delta_{\tau, 1}$ is certainly in the span.
\end{proof}

\begin{corollary}
  \label{cor:dim-cross-product}
  When $K/F$ is a finite dimensional Galois extension, the $K$-dimension of $\cross_{\mfa}$ is $\dim_{F}K$ and the $F$-dimension of $\cross_{\mfa}$ is ${\left(\dim_{F}K\right)}^{2}$.
  \leanok
  \lean{CrossProductAlgebra.dim_eq_square}
\end{corollary}

Now we see that cross product, like a good representation, is a $K$-module and $F$-algebra with a $K$-embedding and correct $F$-dimension. In the next sections, we prove that $\cross_{a}$ is in fact a central simple $F$-algebra.

\subsubsection{Central Algebra}

We will assume $K/F$ is a finite dimensional Galois extension.

\begin{theorem}[Centrality]\label{thm:cross-product-central}
  $\cross_{\mfa}$ is a central $F$-algebra.
  \leanok
  \lean{CrossProductAlgebra.is_central}
  % \uses{lem:cross-product-basis-mul,lem:cross-product-basis-conj, lem:cross-product-basis}
\end{theorem}

\begin{proof}
  Let $z \in \cross_{\mfa}$ that is in the centre. We want to prove that $z$ is in $F$. We write $z$ as $\sum_{\sigma}\lambda_{\sigma}\cdot\Delta_{\sigma,1}$. We see that, for any $\tau \in \gal(K/F)$, we have
  \[
    z = \sum_{\sigma\in\gal(K/F)} \lambda_{\tau^{-1}\sigma\tau}\cdot\Delta_{\tau^{-1}\sigma\tau,1}.
  \]
  Therefore for any $d \in K$ and $\tau\in\gal(K/F)$, by~\cref{lem:cross-product-basis-mul} and~\cref{lem:cross-product-basis-conj}, we have
  \[
    \begin{aligned}
      z\, \Delta_{\tau,d}
      &= \sum_{\sigma\in\gal(K/F)}\lambda_{\sigma}\cdot\Delta_{\sigma, 1}\Delta_{\tau,d}\\
      &= \sum_{\sigma\in\gal(K/F)}\lambda_{\sigma}\cdot\sigma(d)\cdot \mfa(\sigma,\tau) \cdot \Delta_{\sigma\tau, 1}\\
      &= \sum_{\sigma\in\gal(K/F)}\left(\lambda_{\sigma}\sigma(d)\mfa(\sigma,\tau)\right)\cdot \Delta_{\sigma\tau,1}\\
      %
      %
      %
      \Delta_{\tau,d}\, z
      &= \sum_{\sigma\in\gal(K/F)}\Delta_{\tau,d}\left(\lambda_{\tau^{-1}\sigma\tau}\cdot \Delta_{\tau^{-1}\sigma\tau, 1}\right)\\
      &= \sum_{\sigma\in\gal(K/F)}\Delta_{\tau,d}\iota_{\mfa}\left(\lambda_{\tau^{-1}\sigma\tau}\right)\Delta_{\tau^{-1}\sigma\tau, 1} \\
      &=\sum_{\sigma\in\gal(K/F)} d \cdot \Delta_{\tau, 1}\iota_{\mfa}\left(\lambda_{\tau^{-1}\sigma, \tau}\right)\Delta_{\tau^{-1}\sigma\tau, 1}\\
      &=\sum_{\sigma\in\gal(K/F)} d \cdot \Delta_{\tau, \tau\cdot \lambda_{\tau^{-1}\sigma\tau}}\Delta_{\tau^{-1}\sigma\tau, 1}\\
      &=\sum_{\sigma\in\gal(K/F)} d\cdot \tau\left(\lambda_{\tau^{-1}\sigma\tau}\right)\cdot \Delta_{\tau,1}\Delta_{\tau^{-1}\sigma\tau, 1}\\
      &= \sum_{\sigma\in\gal(K/F)}d\cdot \tau\left(\lambda_{\tau^{-1}\sigma\tau}\right)\cdot \mfa\left(\tau,\tau^{-1}\sigma\tau\right)\cdot \Delta_{\sigma\tau, 1}\\
      &= \sum_{\sigma\in\gal(K/F)}\left(d\tau\left(\lambda_{\tau^{-1}\sigma\tau}\right)\mfa\left(\tau,\tau^{-1}\sigma\tau\right)\right) \cdot \Delta_{\sigma\tau,1}.
    \end{aligned}
  \]
  By~\cref{lem:cross-product-basis}, for any $\sigma,\tau\in\gal(K/F)$ and $d\in K$, we have that
  \begin{equation}
    \label{eq:cross-product-central-eqn}
    \lambda_{\sigma}\sigma(d)\mfa(\sigma,\tau) = d\tau\left(\lambda_{\tau^{-1}\sigma\tau}\right)\mfa\left(\tau,\tau^{-1}\sigma\tau\right).
  \end{equation}
  In particular, with $d = 1$, we have
  \[
    \lambda_{\sigma}\mfa(\sigma,\tau) = \tau\left(\lambda_{\tau^{-1}\sigma\tau}\right)\mfa\left(\tau,\tau^{-1}\sigma\tau\right),
  \]
  we substitute back into~\cref{eq:cross-product-central-eqn} and get
  \[
    \lambda_{\sigma}\sigma(d)\mfa(\sigma,\tau) = d\lambda_{\sigma}\mfa(\sigma,\tau).
  \]
  With $\tau = \id$, we have
  \[
    \lambda_{\sigma}\sigma(d)\mfa(\sigma,\id)=d\lambda_{\sigma}\mfa(\sigma,\id),
  \]
  Hence for all $d \in K$ with $\lambda_{\sigma} \ne 0$, we have $\sigma(d) = d$. We immediately deduce that for all $\sigma \ne \id$, $\lambda_{\sigma} = 0$ by contraposition. Thus $z = \lambda_{\id}\Delta_{\id, 1} = \Delta_{\id,\lambda_{\id}} = \iota_{\mfa}\left(\lambda_{\id}\mfa(\id,\id)\right) = \left(\lambda_{id}\mfa(\id,\id)\right) \cdot 1$. Consequently, to prove $z$ is in $F$, it is sufficient to prove that $\lambda_{\id}\mfa(\id,\id)$ is in $F$. Since $K/F$ is finite dimensional and Galois, we only need to prove that $\lambda_{\id}\mfa(\id,\id)$ is fixed by every $\tau\in\gal(K/F)$.Indeed, with $d = 1$ and $\sigma = \id$ in~\cref{eq:cross-product-central-eqn}, we have
  \[
    \begin{aligned}
      \lambda_{\id}\mfa(\id,\tau) &= \tau\left(\lambda_{\id}\right)\mfa(\tau,\id)\\
      &= \lambda_{\id}\mfa(\id,\id) \\ &= \tau\left(\lambda_{\id}\right)\tau\left(\mfa(\id,\id)\right)\\
      &= \tau\left(\lambda_{\id}\mfa(\id,\id)\right).
    \end{aligned}
  \]

\end{proof}


\subsubsection{Simple Ring}
In this section we assume $K/F$ is a finite dimensional field extension. Let $I \subseteq \cross_{\mfa}$ be a two sided ideal, we aim to show that either $I = \{0\}$ or $I = \cross_{\mfa}$. In this section, we use $\pi$ to denote the canonical ring homomorphism $\cross_{\mfa}\to{}^{\cross_{\mfa}}/_{I}$. We restrict $\pi$ to $\pi|_{\im\left(\iota_{\mfa}\right)} : \im\left(\iota_{\mfa}\right) \to {}^{\cross_{\mfa}}/_{I}$ and denote the range of $\pi|_{\im\left(\iota_{\mfa}\right)}$ to be $\Pi$.

\begin{construction}\label{con:cross-product-simple-ring-mod}
  The quotient ring ${}^{\cross_{\mfa}}/_{I}$ is a $\Pi$-module defined by $\pi\left(\iota_{\mfa}(a)\right)\cdot \pi(y) := \pi(a \cdot y)$. We first check that the $\Pi$-action is well-defined:
  \begin{itemize}
    \item Independence of $a$: Let $a, b \in K$ be such that $\pi\left(\iota_{\mfa}(a)\right) = \pi\left(\iota_{\mfa}(b)\right)$, that is, $\iota_{\mfa}(a-b)\in I$. Since $I$ is a two sided ideal, $a\cdot y - b \cdot y = (a - b)\cdot y = \iota_{\mfa}(a - b) y$ is also in $I$. This proves $\pi(a\cdot y) = \pi(b\cdot y)$.
    \item Independence of $y$: Let $y_{1},y_{2}\in\cross_{\mfa}$ be such that $y_{1}-y_{2}\in I$, then for any $a \in K$, $a \cdot y_{1} - a \cdot y_{2} = \iota_{\mfa}(a)\left(y_{1}-y_{2}\right)$ is in $I$ because $I$ is a two sided ideal. This proves $\pi\left(\iota_{\mfa}(a)\right)\cdot \pi(y_{1}) = \pi\left(\iota_{\mfa}(a)\right)\cdot \pi(y_{2})$.
  \end{itemize}
  Then we check the axioms of module:
  \begin{itemize}
    \item Let $y \in \cross_{\mfa}$, we check that $1 \cdot \pi(y) = \pi(y)$ and $0 \cdot \pi(y) = 0$. This is because $\Pi\ni 1 = \pi\left(\iota_{\mfa}(1)\right)$ and $\Pi\ni 0=\pi\left(\iota_{\mfa}(0)\right)$. Let $a \in K$, $\pi\left(\iota_{\mfa}(a)\right)\cdot 0 = 0$ because $0 \in {}^{\cross_{\mfa}}/_{I}$ is equal to $\pi(0)$.
    \item Let $a, b \in K$ and $x, y \in \cross_{\mfa}$, we check $\left(\pi\left(\iota_{\mfa}(a)\right) + \pi\left(\iota_{\mfa}(b)\right)\right)\cdot \pi(x) = \pi\left(\iota_{\mfa}(a)\right)\cdot\pi(x) + \pi\left(\iota_{\mfa}(b)\right)\cdot \pi(x)$ and $\pi\left(\iota_{\mfa}(a)\right)\cdot(\pi(x)+\pi(y)) = \pi\left(\iota_{\mfa}(a)\right)\cdot\pi(x) + \pi\left(\iota_{\mfa}(a)\right)\cdot\pi(y)$. These are true because $\pi$ preserves addition. Similarly $\pi\left(\iota_{\mfa}(a)\right)\cdot \pi\left(\iota_{\mfa}(b)\right) \cdot \pi(x) = \pi\left(\iota_{\mfa}(ab)\right)\cdot \pi(x)$ because $\pi$ preserves multiplication as well.
  \end{itemize}

  Hence ${}^{\cross_{\mfa}}/_{I}$ is also a $K$-module by pulling back the $\Pi$-module structure along $K \to \Pi$ given by $a \mapsto \pi\left(\iota_{\mfa}(a)\right)$. Note that $\pi$ is a $K$-linear map between $\cross_{\mfa}$ and ${}^{\cross_{\mfa}}/_{I}$ by this construction.
  \leanok
  \lean{CrossProductAlgebra.is_simple_proofs.smul_def_quot, CrossProductAlgebra.is_simple_proofs.K_smul_quot}
\end{construction}

\begin{lemma}\label{lem:cross-product-quotient-basis}
  If $I \ne \cross_{\mfa}$, the set $\{\pi\left(\Delta_{\sigma, 1}\right)|\sigma\in\gal(K/F)\}$ forms a $K$-basis for ${}^{\cross_{\mfa}}/_{I}$.
  \leanok
  \lean{CrossProductAlgebra.is_simple_proofs.basis}
  % \uses{lem:cross-product-basis,thm:conj-seq-basis,con:cross-product-simple-ring-mod,con:cross-product-iota, lem:cross-product-basis-conj, lem:cross-product-delta-invertible}
\end{lemma}
\begin{proof}
  It is easy to see that the set spans ${}^{\cross_{\mfa}}/_{I}$ because $\{\Delta_{\sigma, 1}|\sigma\in\gal(K/F)\}$ spans $\cross_{\mfa}$ (\cref{lem:cross-product-basis}). For linear-independence, the idea is the same as in the proof of~\cref{thm:conj-seq-basis}. We repeat the argument here.

  Suppose that $\{\pi\left(\Delta_{sigma, 1}\right)|\sigma\in\gal(K/F)\}$ is linearly dependent. Let $J \subseteq \gal(K/F)$ be such that $\{\pi\left(\Delta_{\sigma, 1}\right)|\sigma\in J\}$ is the maximally linearly independent set. Let $\sigma$ be an arbitrary automorphism that is not in $J$. Therefore, we have $\pi\left(\Delta_{\sigma, 1}\right) \in \left\langle \pi\left(\Delta_{\tau, 1}\right) | \tau\in J \right\rangle$. Hence we have, by~\cref{con:cross-product-simple-ring-mod} and~\cref{con:cross-product-iota}
  \[
    \pi\left(\Delta_{\sigma,1}\right) = \sum_{\tau\in J'}\lambda_{\tau}\cdot\pi\left(\Delta_{\tau, 1}\right) = \sum_{\tau\in J'}\pi\left(\iota_{\mfa}\left(\lambda_{\tau}\right)\right)\pi\left(\Delta_{\tau, 1}\right) = \sum_{\tau\in J'}\pi\left(\iota_{\mfa}\left(\lambda_{\tau}\right)\Delta_{\tau,1}\right) =\sum_{\tau\in J'}\pi\left(\lambda_{\tau}\cdot \Delta_{\tau, 1}\right).
  \]
  for some non-zero $\lambda_{\tau} \in K$ and some $J' \subseteq J$.
  Hence, for any $c \in K$, we have
  \[
    \begin{aligned}
      \pi\left(\iota_{\mfa}\left(\sigma \cdot c\right)\right) \pi\left(\Delta_{\sigma, 1}\right)
      &= \pi\left(\Delta_{\sigma, 1}\right)\pi\left(\iota_{\mfa}(c)\right) &\text{by~\cref{lem:cross-product-basis-conj}} \\
      &= \sum_{\tau\in J}
        \pi\left(\iota_{\mfa}\left(\lambda_{\tau}\right)\right)
        \pi\left(\Delta_{\tau, 1}\right)
        \pi\left(\iota_{\mfa}(c)\right)\\
      &= \sum_{\tau\in J}
        \pi\left(
        \iota_{\mfa}\left(\lambda_{\tau}\right)\Delta_{\tau, 1} \iota_{\mfa} c
        \right) \\
      &= \sum_{\tau\in J}
        \pi\left(
        \iota_{\mfa}\left(\lambda_{\tau}\right)\iota_{\mfa}(\tau\cdot c)\Delta_{\tau, 1}
        \right) & \text{by~\cref{lem:cross-product-basis-conj} again}\\
      &= \sum_{\tau\in J}
        \pi\left(\iota_{\mfa}\left(\lambda_{\tau}\tau(c)\right)\right)
        \pi\left(\Delta_{\tau,1}\right)\\
      &= \sum_{\tau\in J}
        \left(\lambda_{\tau}\tau(c)\right)\cdot \pi\left(\Delta_{\tau, 1}\right)\\
        %%%%
        %%%%
      \pi\left(\iota_{\mfa}\left(\sigma \cdot c\right)\right) \pi\left(\Delta_{\sigma, 1}\right)
      &= \sum_{\tau\in J}
        \pi\left(\iota_{\mfa}(\sigma\cdot c)\right)
        \pi\left(\iota_{\mfa}\left(\lambda_{\tau}\right)\right)
        \pi\left(\Delta_{\tau, 1}\right) \\
      &= \sum_{\tau\in J}
        \pi\left(\iota_{\mfa}\left(\sigma(c)\lambda_{\tau}\right)\right)
        \pi\left(\Delta_{\tau, 1}\right) \\
      &= \sum_{\tau\in J}
        \left(\sigma(c)\lambda_{\tau}\right)\cdot \pi\left(\Delta_{\tau, 1}\right).
    \end{aligned}
  \]
  Since, $\{\pi\left(\Delta_{\tau, 1}\right) | \tau \in J\}$ is linearly independent, for all $c \in K$ and $\tau\in J$, we have that $\lambda_{\tau}\tau(c) = \sigma(c)\lambda_{\tau}$. Note that $J' \ne \emptyset$, otherwise, $\pi\left(\Delta_{\sigma, 1}\right) = 0$ implying that $\Delta_{\sigma, 1} \in I$ which by~\cref{lem:cross-product-delta-invertible} is invertible but $I$ does not equal to $\cross_{\mfa}$. Hence for each $\tau \in J'$, we have that for all $c \in K$, since $\lambda_{\tau}$ is not zero, $\sigma(c) = \tau(c)$, i.e. $\sigma = \tau$. Therefore, $\sigma$ is in $J' \subseteq J$ after all.
\end{proof}

\begin{corollary}\label{cor:cross-product-quot-iso}
  If $I \ne \cross_{\mfa}$, the quotient ring ${}^{\cross_{\mfa}}/_{I}$ is isomorphic to $\cross_{\mfa}$ as $K$-modules. In particular $\pi$ is a $K$-linear isomorphism between $\cross_{\mfa}$ and the quotient ring ${}^{\cross_{\mfa}}/_{I}$.
  \leanok
  \lean{CrossProductAlgebra.is_simple_proofs.π₁}
  % \uses{lem:cross-product-basis,lem:cross-product-quotient-basis}
\end{corollary}

\begin{proof}
  Indeed, by~\cref{lem:cross-product-basis}, $\{\Delta_{\sigma, 1}|\sigma \in \gal(K/F)\}$ is a $K$-basis for $\cross_{\mfa}$; and by~\cref{lem:cross-product-quotient-basis}, $\{\pi\left(\Delta_{\sigma, 1}\right) | \sigma\in\gal(K/F)\}$ is a $K$-basis for ${}^{\cross_{\mfa}}/_{I}$. The two sets obviously biject. Hence we can define a $K$-linear isomorphism by $\Delta_{\sigma, 1} \mapsto \pi\left(\Delta_{\sigma, 1}\right)$. This isomorphism is equal to $\pi$ everywhere.
\end{proof}

\begin{corollary}[Simple Ring]\label{cor:cross-product-simple}
  $\cross_{\mfa}$ is a simple ring.
  \leanok
  \lean{CrossProductAlgebra.is_simple}
  % \uses{cor:cross-product-quot-iso}
\end{corollary}

\begin{proof}
  For any two-sided-ideal $I$ that is not equal to $\cross_{\mfa}$, by~\cref{cor:cross-product-quot-iso}, $\pi : cross_{\mfa} \to {}^{\cross_{\mfa}}/_{I}$ is a $K$-linear isomorphism, therefore $I$ is equal to $0$.
\end{proof}

\begin{theorem}\label{thm:cross-product-csa}
  Let $K/F$ be a finite dimensional and Galois field extension and $\mfa$ be a 2-cocycle in $\mathcal{B}^{2}\left(\gal(K/F), K^{\star}\right)$, $\cross_{a}$ is a finite dimensional central simple $F$-algebra.
  \leanok
  \lean{CrossProductAlgebra.asCSA}
  % \uses{thm:cross-product-central, lem:cross-product-basis, cor:cross-product-simple}
\end{theorem}
\begin{proof}
  \Cref{thm:cross-product-central}, \cref{lem:cross-product-basis} and \cref{cor:cross-product-simple}.
\end{proof}

\subsection{From $\HH^{2}\left(\gal(K/F), K^{\star}\right)$ to $\br(K/F)$}

For every 2-cocycle $\mfa$, we have defined the cross product $\cross_{\mfa}$ and proved that it is indeed a finite dimensional central simple $F$-algebra in~\cref{thm:cross-product-csa}; that is we have a function from $\mathcal{B}^{2}\left(\gal(K/F), K^{\star}\right)$ to $\csa_{F}$. If we want a function from $\HH^{2}\left(\gal(K/F), K^{\star}\right)$ to $\br(K/F)$, we need to show that if $\mfa$ and $\mfb$ are cohomologous, $\cross_{\mfa}$ and $\cross_{\mfb}$ are Brauer equivalent. We state it as a theorem:

\begin{theorem}\label{thm:snd-coh-to-relative-br}
  If $K/F$ is a finite dimensional and Galois field extension,
  the function $\cross : \HH^{2}\left(\gal(K/F), K^{\star}\right) \to \br(K/F)$ defined by
  \[
    \mfa \mapsto \left[\cross_{\mfa}\right]_{\sim_{\br}}
  \]
  is well-defined.
  \leanok
  \lean{RelativeBrGroup.fromSnd, RelativeBrGroup.fromSnd_wd}
  % \uses{def:second-group-coh, con:cross-product-iota, lem:cross-product-basis-mul, lem:cross-product-basis-conj}
\end{theorem}

\begin{proof}
  Let $\mfa$ and $\mfb$ be two cohomologous 2-cocycles. By~\cref{def:second-group-coh}, for some $\mfc: \gal(K/F) \to K^{\star}$, for all $\sigma, \tau\in\gal(K/F)$, we have
    \begin{equation}
      \label{eq:cross-product-cohomologous}
    \frac{\sigma(\mfc(\tau))}{\mfc(\sigma\tau)}\mfc(\sigma) = \frac{\mfa(\sigma, \tau)}{\mfb(\sigma, \tau)}.
    \end{equation}
    Let us denote $\mathsf{A}$ to be the $K$-basis $\{\Delta^{\mfa}_{\sigma, 1} | \sigma\in\gal(K/F)\}$ for $\cross_{\mfa}$ and $\mathsf{B}$ to be the $K$-basis $\{\mfc(\sigma)\cdot \Delta^{\mfb}_{\sigma, 1}|\sigma\in\gal(K/F)\}$ for $\cross_{\mfb}$. We immediately have a $K$-linear isomorphism $\phi : \cross_{\mfa}\cong \cross_{\mfb}$ by mapping $\mathsf{A}$ to $\mathsf{B}$. Since the $K$-action on $\cross_{a}$ and $\cross_{b}$ agrees with the $F$-action on them (\cref{con:cross-product-iota}), $\phi$ is also an $F$-linear isomorphism. We check that $\phi(1) = 1$ and $\phi(xy) = \phi(x)\phi(y)$ for all $x, y \in \cross_{\mfa}$:
    \begin{enumerate}
      \item preservation of one: with $\sigma = \tau =\id$ in \cref{eq:cross-product-cohomologous}, we have $\mfc(\id) = \mfa(\id,\id)\mfb(\id,\id)^{-1}$, thus
  \[
    \begin{aligned}
      \phi(1)
      &= \phi\left(\Delta^{\mfa}_{\id,\mfa(\id,\id)^{-1}}\right)\\
      &= \phi\left(
        \mfa(\id,\id)^{-1}\cdot \Delta^{\mfa}_{\id, 1}
        \right) \\
      &= \mfa(\id,\id)^{-1}\cdot \mfc(\id) \cdot \Delta^{\mfb}_{\id, 1} \\
      &= \mfa(\id,\id)^{-1}\cdot \mfc(\id)\cdot \mfb(\id,\id) \cdot \mfb(\id,\id)^{-1}\cdot \Delta^{\mfb}_{\id, 1} \\
      &= \left(
        \mfa(\id,\id)^{-1}\mfc(\id)\mfb(\id,\id)
        \right)\cdot
        \left(
        \mfb(\id,\id)^{-1}\cdot \Delta_{\id, 1}
        \right) \\
      &= \left(\mfa(\id,\id)^{-1}\mfa(\id,\id)\right)\cdot \Delta_{\id,\mfb(\id,\id)^{-1}} \\
      &= \Delta_{\id,\mfb(\id,\id)^{-1}}.
    \end{aligned}
  \]
      \item preservation of multiplication:
            let $\sigma,\tau\in\gal(K/F)$ and $a, b \in K$, we need to prove that $\phi\left(\Delta^{\mfa}_{\sigma, a}\Delta^{\mfa}_{\tau, b}\right) = \phi\left(\Delta^{\mfa}_{\sigma, a}\right)\phi\left(\Delta^{\mfb}_{\tau, b}\right)$.
            From~\cref{eq:cross-product-cohomologous}, we see that
            \[
            \sigma(\mfc(\tau))\mfc(\sigma)\mfb(\sigma,\tau) = \mfc(\sigma\tau)\mfa(\sigma,\tau).\]
            Hence, by~\cref{lem:cross-product-basis-mul} and~\cref{lem:cross-product-basis-conj}, we have
            \[
            \begin{aligned}
              \phi\left(\Delta^{\mfa}_{\sigma, a}\Delta^{\mfa}_{\tau, b}\right)
              &= \phi\left(
                \Delta^{\mfa}_{\sigma\tau, a \sigma(b) \mfa(\sigma,\tau)}
                \right)\\
              &= \phi\left(
                % \left(
                a\sigma(b) \mfa(\sigma,\tau) \,
                % \right)
                \cdot \,
                \Delta^{\mfa}_{\sigma\tau, 1}
                \right)\\
              &=
                % \left(
                a \sigma(b) \mfa(\sigma,\tau)
                % \right)
                \, \cdot \,
                \phi\left(\Delta^{\mfa}_{\sigma\tau, 1}\right) \\
              &=
                % \left(
                a \sigma(b) \mfa(\sigma,\tau) \mfc(\sigma\tau)
                % \right)
                \, \cdot \,
                \Delta^{\mfb}_{\sigma\tau, 1}.\\
              &=
                % \left(
                a\sigma(b)\sigma(\mfc(\tau))\mfc(\sigma)\mfb(\sigma,\tau)
                % \right)
                \, \cdot \,
                \Delta^{\mfb}_{\sigma\tau, 1}\\
              &=
                % \left(
                \mfc(\sigma) \sigma(\mfc(\tau))
                % \right)
                \,\cdot\,
                \Delta^{\mfb}_{\sigma\tau,a\sigma(b)\mfb(\sigma\tau)}\\
              &=
                % \left(
                \mfc(\sigma)\sigma(\mfc(\tau))
                % \right)
                \, \cdot \,
                % \left(
                \Delta^{\mfb}_{\sigma, a}\Delta^{\mfb}_{\tau, b}
                % \right)
              \\
              %
              %
              %
              \phi\left(\Delta^{\mfa}_{\sigma, a}\right)\phi\left(\Delta^{\mfa}_{\tau, b}\right)
              &= \phi\left(a \cdot \Delta^{\mfa}_{\sigma, 1}\right) \phi\left(b\cdot \Delta^{\mfa}_{\tau, 1}\right) \\
              &= \left(a \cdot \phi\left(\Delta^{\mfa}_{\sigma, 1}\right)\right)
                \left(b \cdot \phi\left(\Delta^{\mfa}_{\tau, 1}\right)\right) \\
              &= \left(
                a\mfc(\sigma)
                \, \cdot \,
                \Delta^{\mfb}_{\sigma, 1}
                \right)
                \left(
                b \mfc(\tau)
                \, \cdot \,
                \Delta^{\mfb}_{\tau, 1}
                \right)\\
              &= a\mfc(\sigma) \, \cdot \,
                \left(\Delta^{\mfb}_{\sigma, 1}
                \iota_{\mfb}(b\mfc(\tau))\right) \,
                \Delta^{\mfb}_{\tau, 1} \\
              &=
                a\mfc(\sigma) \, \cdot \,
                \Delta^{\mfb}_{\sigma, \sigma(b\mfc(\tau))}\,
                \Delta^{\mfb}_{\tau, 1} \\
              &= a\mfc(\sigma)\sigma(b)\sigma(\mfc(\tau)) \, \cdot \,
                \Delta^{\mfb}_{\sigma, 1}\Delta_{\tau, 1}^{\mfb}\\
              &= a\mfc(\sigma)\sigma(b)\sigma(\mfc(\tau))\mfb(\sigma,\tau) \, \cdot \, \Delta^{\mfb}_{\sigma\tau, 1}\\
              &= \mfc(\sigma)\sigma(\mfc(\tau)) \, \cdot \, \Delta^{\mfb}_{\sigma\tau, a\sigma(b)\mfb(\sigma,\tau)}\\
              &= \mfc(\sigma)\sigma(\mfc(\tau)) \, \cdot \, \Delta^{\mfb}_{\sigma, a}\Delta^{\mfb}_{\tau, b}.
            \end{aligned}
            \]
    \end{enumerate}
    Hence $\phi$ is actually an $F$-algebra isomorphism between $\cross_{\mfa}$ and $\cross_{\mfb}$ and isomorphic central simple $F$-algebras are certainly Brauer equivalent.
  \end{proof}



%%% Local Variables:
%%% mode: LaTeX
%%% TeX-master: "../print"
%%% End:


%%FILE blueprint/src/content/brauer-group2.tex
\subsection{$\HH^{2} \circ \cross$ and $\cross \circ \HH^{2}$}

For a finite dimensional Galois extension of field $K/F$, we have constructed two functions $\HH^{2}$ and $\cross$ between the second cohomology group $\HH^{2}\left(\gal(K/F), K^{\star}\right)$ and the relative Brauer group $\br(K/F)$. In this section, we prove that they are mutual inverse to one another,

\begin{lemma}\label{lem:relative-br-snd-inverse-1}
  The composition of $\cross$ and $\HH^{2}$ is the identity:
  \begin{center}
    \begin{tikzcd}
      \displaystyle \HH^2\left(\gal(K/F), K^\star\right) \arrow{r}{\cross} \arrow[bend right = 15, swap]{rr}{\id} &
      \br(K/F) \arrow{r}{\HH^2} & \displaystyle \HH^2\left(\gal(K/F), K^\star\right).
    \end{tikzcd}
  \end{center}
  \leanok
  \lean{RelativeBrGroup.toSnd_fromSnd}
  % \uses{lem:cross-product-basis-conj,con:good-rep-to-2-cocycles,cor:to-2-cocycle-wd, thm:snd-coh-to-relative-br, con:compare-conj-factors}
\end{lemma}
\begin{proof}
  Let $\mfa$ be any 2-cocycle, by~\cref{lem:cross-product-basis-conj}, we notice that $x : \sigma \mapsto \Delta_{\sigma, 1}$ is a conjugation sequence for $\cross_{\mfa}$. Hence by~\cref{con:good-rep-to-2-cocycles}, ~\cref{cor:to-2-cocylce-wd} and~\cref{thm:snd-coh-to-relative-br}, we evaluate the composition at $\mfa$ as:
  \begin{center}
    \begin{tikzcd}
      {[\mfa]} \arrow[mapsto]{r} & {\left[\cross_\mfa\right]_{\sim_{\br}}} \arrow[mapsto]{r}& {\left[(\sigma,\tau)\mapsto\comp^{x}_{\Delta_{\sigma, 1}, \Delta_{\tau, 1}, \Delta_{\sigma\tau, 1}}\right]}
    \end{tikzcd}.
  \end{center}
  That is, we need to show that $\mfa$ and $(\sigma, \tau) \mapsto \comp_{\Delta_{\sigma, 1}, \Delta_{\tau, 1}, \Delta_{\sigma\tau, 1}}$ are 2-cohomologous. In fact, they are equal.
  By~\cref{con:compare-conj-factors}, we have that $\iota_{\cross_{\mfa}}\left(\comp^{x}_{\Delta_{\sigma, 1}, \Delta_{\tau, 1}, \Delta_{\sigma\tau, 1}}\right) = \Delta_{\sigma, 1}\Delta_{\tau, 1}\Delta_{\sigma\tau, 1}^{-1} = \mfa(\sigma,\tau)\cdot \Delta_{\sigma\tau, 1}\Delta_{\sigma\tau, 1}^{-1} = \mfa(\sigma,\tau)\cdot 1 = \Delta_{\id,\mfa(\sigma,\tau)}$ which is precisely $\iota_{\cross_{\mfa}}(\mfa(\sigma,\tau))$.
\end{proof}

\begin{lemma}\label{lem:relative-br-snd-inverse-2}
  The composition of $\HH^{2}$ and $\cross$ is the identity:
  \begin{center}
    \begin{tikzcd}
      \br(K/F) \arrow{r}{\HH^2} \arrow[bend right = 15, swap]{rr}{\id} &
      \HH^2\left(\gal(K/F), K^\star\right) \arrow{r}{\cross} &
      \br(K/F)
    \end{tikzcd}.
  \end{center}
  \leanok
  \lean{RelativeBrGroup.fromSnd_toSnd}
  % \uses{cor:mem-relative-br-iff-good-rep, con:exists-conj-seq, def:good-rep, cor:to-2-cocycle-wd, thm:snd-coh-to-relative-br, con:good-rep-mod, con:compare-conj-factors, def:conj-factor}
\end{lemma}
\begin{proof}
  Let $X \in \br(K/F)$, $A$ be an arbitrary good representation of $X$ and $x$ be an arbitrary $A$-conjugation sequence which exists by~\cref{cor:mem-relative-br-iff-good-rep} and~\cref{con:exists-conj-seq}. By~\cref{def:good-rep}, $X = \left[A\right]_{\sim_{\br}}$.
  Hence by~\cref{cor:to-2-cocylce-wd} and~\cref{thm:snd-coh-to-relative-br}, we evaluate the composition at $X$ as:
  \begin{center}
    \begin{tikzcd}
      {[A]_{\sim_{\br}}} \arrow[mapsto]{r} &
      {\left[\mathcal{B}^2_x\right]} \arrow[mapsto]{r} &
      {\left[\cross_{\mathcal{B}^2_x}\right]}
    \end{tikzcd}.
  \end{center}
  Hence we need to prove that $A$ and $\cross_{\mathcal{B}^{2}_{x}}$ are Brauer equivalent. We will show that they are isomorphic as $F$-algebras. since $\{x_{\sigma}|\sigma\in\gal(K/F)\}$ is a $K$-basis for $A$ and $\{\Delta_{\sigma, 1}|\sigma \in \gal(K/F)\}$ is a $K$-basis for $\cross_{\mathcal{B}^{2}_{x}}$, they are certainly isomorphic as $K$-modules. Let $\phi : \cross_{\mathcal{B}^{2}_{x}} \cong A$ be the $K$-linear isomorphism defined by $\Delta_{\sigma, 1} \mapsto x_{\sigma}$, since the $K$-action on $A$ and the $F$-action on $A$ are compatible (\cref{con:good-rep-mod}), $\phi$ is also an $F$-linear isomorphism. Like in~\cref{thm:snd-coh-to-relative-br}, we check that $\phi(1) = 1$ and $\phi(xy)=\phi(x)\phi(y)$ for all $x,y \in A$:
  \begin{enumerate}
    \item preservation of one: by~\cref{con:compare-conj-factors}, we have
          \[
          \begin{aligned}
            \phi(1)
            &= \phi\left(\Delta_{\id,\mathcal{B}^{2}_{x}(\id,\id)^{-1}}\right) \\
            &= \mathcal{B}^{2}_{x}(\id,\id)^{-1} \phi\left(\Delta_{\id, 1}\right) \\
            &= \mathcal{B}^{2}_{x}(\id,\id)^{-1} x_{\id} \\
            &= \comp_{x_{\id},x_{\id},x_{\id}}^{-1} x_{\id} \\
            &= \comp_{x_{\id},x_{\id},x_{\id}} x_{\id} x_{\id} x_{\id}^{-1} \\
            &= x_{\id} x_{\id}^{-1} \\
            &= 1.
          \end{aligned}
          \]
    \item preservation of multiplication: let $\sigma,\tau\in\gal(K/F)$ and $c,d \in K$, by~\cref{con:compare-conj-factors} and~\cref{def:conj-factor}, we have
          \[
          \begin{aligned}
            \phi\left(\Delta_{\sigma,c}\Delta_{\tau,d}\right)
            &= \phi\left(\Delta_{\sigma\tau, c\sigma(d)\mathcal{B}^{2}_{x}(\sigma,\tau)}\right) \\
            &= c\sigma(d)\mathcal{B}^{2}_{x}(\sigma,\tau) \,\cdot\, \phi\left(\Delta_{\sigma\tau, 1}\right) \\
            &= c\sigma(d)\mathcal{B}^{2}_{x}(\sigma,\tau) \,\cdot\, x_{\sigma\tau} \\
            &= c\sigma(d)\comp_{x_{\sigma},x_{\tau},x_{\sigma\tau}} \,\cdot\, x_{\sigma\tau} \\
            &= c\sigma(d)\,\cdot\, \iota_{A}\left(\comp_{x_{\sigma},x_{\tau},x_{\sigma\tau}}\right)x_{\sigma\tau} \\
            &= c\sigma(d) \,\cdot\, x_{\sigma}x_{\tau}\\
            %%%%
            %%%%
            \phi\left(\Delta_{\sigma,c}\right)\phi\left(\Delta_{\tau,d}\right)
            &=\left(c \cdot \phi\left(\Delta_{\sigma, 1}\right)\right)
              \left(d \cdot \phi\left(\Delta_{\tau, 1}\right)\right) \\
            &= \left(c \cdot x_{\sigma}\right) \left(d \cdot x_{\tau}\right)\\
            &= c \,\cdot\, x_{\sigma}\iota_{A}(d) x_{\tau} \\
            &= c\sigma_{d} \,\cdot\, x_{\sigma}x_{\tau}.
          \end{aligned}
          \]
  \end{enumerate}
\end{proof}

\begin{corollary}\label{cor:relative-br-2nd-coh-bij}
  For a finite dimensional and Galois extension of field $K/F$, the relative Brauer group $K/F$ bijects to the second cohomology group $\HH^{2}\left(\gal(K/F), K^{\star}\right)$ by the following commutative diagram
  \begin{center}
  \begin{tikzcd}
    {\br(K/F)} \arrow{r}{\HH^2} \arrow[equal]{d} & {\HH^2\left(\gal(K/F), K^\star\right)} \arrow[equal]{d} \\
    {\br(K/ F)} & {\HH^2\left(\gal(K/F), K^{\star}\right)} \arrow{l}{\cross}
  \end{tikzcd}.
\end{center}
\leanok
\lean{RelativeBrGroup.equivSnd}
% \uses{lem:relative-br-snd-inverse-1,lem:relative-br-snd-inverse-2}
\end{corollary}

\begin{proof}
  Exactly~\cref{lem:relative-br-snd-inverse-1} and~\cref{lem:relative-br-snd-inverse-2}.
\end{proof}

\subsection{Group Homomorphism}

In previous sections, when $K/F$ is a finite dimensional Galois extension, we have set up a bijection between the relative Brauer group $\br(K/F)$ and the second cohomology group $\HH^{2}\left(\gal(K/F), K^{\star}\right)$. But both functions $\HH^{2}$ and $\cross$ are only set-theoretical function. In this section, we aim to upgrade them to group homomorphisms. Technically, we only need to prove either one of them preserves multiplication; we provide a proof that $\HH^{2}$ preserves one anyway because we found the proof to be entertaining.

\subsubsection{$\cross_{1} = 1$ and $\HH^{2}(1) = 1$}

\begin{theorem}
  \label{thm:map_one}
  The function $\cross:\HH^{2}\left(\gal(K/F), K^{\star}\right) \to \br(K/F)$ preserves one, that is $\cross_{}$
 \leanok
 \lean{map_one_proof.fromSnd_zero}
 % \uses{lem:cross-product-basis, cor:alghom-bijective-of-dim-eq}
\end{theorem}

\begin{proof}
  Since $\{\Delta_{\sigma, 1}|\sigma\in\gal(K/F)\}$ is a $K$-basis for $\cross_{1}$ where $1 \in \mathcal{B}^{2}\left(\gal(K/F),K^{\star}\right)$ is the constant function $1$ (\cref{lem:cross-product-basis}), we construct a $K$-linear map $\phi : \cross_{1} \to \End_{F} K$ by $\Delta_{\sigma, 1} \mapsto \sigma$; note that $\phi$ is $F$-linear as well. In fact, $\phi$ is also an $F$-algebra homomorphism:
  \begin{enumerate}
    \item $\phi(1) = 1$: indeed $\phi\left(\Delta_{\id,1}\right) = \id$.
    \item $\phi(xy) = \phi(x)\phi(y)$: indeed, let $\sigma,\tau\in\gal(K/F)$ and $c,d\in K$, we need to check that $\phi\left(\Delta_{\sigma,c}\Delta_{\tau,d}\right) = \phi\left(\Delta_{\sigma,c}\right)\phi\left(\Delta_{\tau,d}\right)$. The left hand side is equal to
          \[\phi\left(\Delta_{\sigma\tau,c\sigma(d)}\right) = \phi\left(c\sigma(d)\,\cdot\, \Delta_{\sigma\tau,1}\right) = c\sigma(d)\,\cdot\, \sigma\tau;\]
          and the right hand side is equal to
          \[
        \phi\left(c\cdot \Delta_{\sigma,1}\right)\phi\left(d\cdot \Delta_{\tau, 1}\right) = (c \cdot \sigma) (d \cdot \tau).
          \]
          For any $x \in K$, applying left hand side to $x$ will result in $c\sigma(d)\sigma(\tau(x))$ while right hand side will result in $c\sigma(d\tau(x))$, hence both sides are equal.
  \end{enumerate}
  Hence, $\phi$ is an $F$-algebra isomorphism by~\cref{cor:alghom-bijective-of-dim-eq}; that is we have $\cross_{1} \cong \End_{F}K \cong\mat_{\dim_{F}K}(F)$. We conclude that $\cross_{1}$ is Brauer equivalent to $F$ and consequently $\HH^{2}(1) = 1$.
\end{proof}


\begin{corollary}
  The function $\HH^{2}:\br(K/F) \to \HH^{2}\left(\gal(K/F), K^{\star}\right)$ preserves one, that is $\HH^{2}(1) = 1$.
  \leanok
  \lean{RelativeBrGroup.toSndAddMonoidHom}
  % \uses{lem:relative-br-snd-inverse-2,thm:map_one}
\end{corollary}
\begin{proof}
  Apply $\cross$ then use~\cref{lem:relative-br-snd-inverse-2} and~\cref{thm:map_one}.
\end{proof}

\subsubsection{$\cross_{\mfa\mfb}\sim_{\br}\cross_{\mfa}\otimes_{F}\cross_{\mfb}$}
The argument in this section is more complicated, because, unlike before, the left hand side and the right hand side are not isomorphic as $F$-algebras --- left hand side has $F$-dimension $\left(\dim_{F}K\right)$ while the right hand side has $F$-dimension $\left(\dim_{F}K\right)^{4}$. Let $\mfa$ and $\mfb$ be two 2-cocycles in $\mathcal{B}^{2}\left(\gal(K/F), K^{\star}\right)$, we denote $\mfc$ to be the 2-cocycle $\mfa \mfb$. Intuitively, $\cross_{\mfa} \ox_{F} \cross_{\mfb}$ is too ``big'', to address this issue we introduce a quotient module.

\begin{construction}[$M$]
  Consider the quotient module
  \[M := {}^{\cross_{\mfa}\ox_{F}\cross_{\mfb}}/_{\left\langle (k\cdot a)\ox b - a \ox (k\cdot b) | k \in K, a \in \cross_{\mfa}, b \in \cross_{\mfb} \right\rangle}.\]
  For any $a' \in \cross_{\mfa}$ and $b' \in \cross_{\mfb}$, we can define an $F$-linear map $M \to M$ by descending the $F$-linear map
$\cross_{\mfa}\ox_{F}\cross_{\mfb} \to \cross_{\mfa}\ox \cross_{\mfb}$
  \[
      a\ox b \mapsto aa' \ox bb';
    \]
    we need to check that for all $k \in K, a \in \cross_{\mfa}, b \in \cross_{\mfb}$, the image of $(k\cdot a)\ox b - a\ox(k\cdot b)$ is in $\left\langle (k\cdot a)\ox b - a \ox (k\cdot b) | k \in K, a \in \cross_{\mfa}, b \in \cross_{\mfb} \right\rangle$:
    the image is $(k\,\cdot\, aa')\ox b - a\ox(k\,\cdot\, bb')$ which is in the generating set with $k \in K, aa' \in\cross_{\mfa}, $ and $bb'\in\cross_{\mfb}$.
  This map is in fact $F$-linear in both $a'$ and $b'$, hence we have an $F$-bilinear map $\cross_{\mfa}\ox\cross_{\mfb} \to M \to M$.
  This gives $M$ a $\left(\cross_{\mfa}\ox_{F}\cross_{\mfb}\right)^{\opp}$-module structure given by
  \[
    (a' \ox b') \cdot [a\ox b] = aa' \ox bb'
  \]
  for any $a,a' \in \cross_{\mfa}$ and $b, b' \in\cross_{\mfb}$. All of the module axioms in this case follows from $F$-bilinearity.


  For any $c \in \cross_{\mfc}$, we can define another $F$-linear map $M \to M$ by descending the $F$-linear map $\cross_{\mfa}\ox_{F}\cross_{\mfb} \to \cross_{\mfa}\ox \cross_{\mfb}$
  \[
      a \ox b \mapsto \sum_{\sigma\in\gal(K/F)} \Delta^{\mfa}_{\sigma,c(\sigma)}a \, \ox \, \Delta^{\mfb}_{\sigma, 1}b;
    \]
    we need check that for all $k \in K, a \in\cross_{\mfa}, b \in \cross_{\mfb}$, the image of $(k\cdot a)\ox b - a \ox (k\cdot b)$ is in
    $\left\langle (k\cdot a)\ox b - a \ox (k\cdot b) | k \in K, a \in \cross_{\mfa}, b \in \cross_{\mfb} \right\rangle$: by~\cref{lem:cross-product-basis-conj} the image is
    \[
\begin{aligned}
      & \sum_{\sigma\in\gal(K/F)}
      \Delta^{\mfa}_{\sigma, c(\sigma)}(k \cdot a) \,\ox\, \Delta^{\mfb}_{\sigma, 1}b
      -
        \Delta^{\mfa}_{\sigma, c(\sigma)}a \,\ox\, \Delta^{\mfb}_{\sigma, 1}(k \cdot b) \\
  = & \sum_{\sigma\in\gal(K/F)}
      \Delta^{\mfa}_{\sigma,c(\sigma)}\iota_{\mfa}(k)a \,\ox\, \Delta^{\mfb}_{\sigma,1}b
      -
      \Delta^{\mfa}_{\sigma,c(\sigma)} \,\ox\, \Delta^{\mfb}_{\sigma, 1}\iota_{\mfb}(k)b \\
  =& \sum_{\sigma\in\gal(K/F)}
     \sigma(k)\cdot\Delta^{\mfa}_{\sigma, c(\sigma)}a \,\ox\, \Delta^{\mfb}_{\sigma, 1}b
     - \Delta^{\mfa}_{\sigma, c(\sigma)} \,\ox\, \sigma(k)\cdot\Delta^{\mfb}_{\sigma,1}b,
  \end{aligned}
\]
which is in $\left\langle (k\cdot a)\ox b - a \ox (k\cdot b) | k \in K, a \in \cross_{\mfa}, b \in \cross_{\mfb} \right\rangle$ because for each $\sigma \in \gal(K/F)$, the summand is in the generating set with $\sigma(k) \in K, \Delta^{\mfa}_{\sigma, c(\sigma)}a \in \cross_{\mfa}$ and $\Delta^{\mfb}_{\sigma, 1}b \in \cross_{\mfb}$. This map is in fact $F$-linear in $c$, therefore we have an $F$-bilinear map $\cross_{\mfc}\to M \to M$. (In the above calculation ``$\ox$'' symbol has low precedence.) This gives $M$ a $\cross_{\mfc}$-module structure given by
\[
  c \cdot [a \ox b] = \left[\sum_{\sigma\in\gal(K/F)}\Delta^{\mfa}_{\sigma, c(\sigma)}a \ox \Delta^{\mfb}_{\sigma, 1}b\right].
\]
In particular, if $c$ is of the form $k \cdot \Delta^{\mfc}_{\tau, 1}$, then $\left(k\cdot \Delta^{\mfc}_{\tau, 1}\right)\cdot [a \ox b]$ is equal to $\left[\left(k\cdot \Delta^{\mfa}_{\tau,1}\right)a \ox \Delta^{\mfb}_{\tau, 1} b\right]$ because $\Delta_{\tau, 1}^{\mfc}(\sigma) = 0$ for all $\sigma \ne \tau$. Two of the module axioms need more than $F$-bilinearity:
\begin{itemize}
  \item $c = 1$: note that $c = 1 = \mfb(\id,\id)^{-1}\mfa(\id,\id)^{-1}\,\cdot\,\Delta^{\mfc}_{\id, 1}$, hence
        \[
        \begin{aligned}
          1 \cdot [a\ox b]
          &= [\mfb(\id,\id)^{-1}\mfa(\id,\id)^{-1}\Delta^{\mfa}_{\id,1}a \ox \Delta^{\mfb}_{\id, 1}b] \\
          &= [\Delta^{\mfa}_{\id,\mfa(\id,\id)^{-1}}a \ox \Delta^{\mfb}_{\id,\mfb(\id,\id)^{-1}}b] \\
          &= [a\ox b].
        \end{aligned}
        \]
  \item $c_{1}c_{2}\cdot [a \ox b] = c_{1}\cdot c_{2}\cdot [a\ox b]$: assume $c_{1} = k_{1}\cdot\Delta^{\mfc}_{\tau_{1}, 1}$ and $c_{2}=k_{2}\cdot\Delta^{\mfc}_{\tau_{2},1}$. Then $c_{1}c_{2} = k_{1}\tau_{1}\left(k_{2}\right)\cdot \Delta^{\mfc}_{\tau_{1}\tau_{2},\mfc(\tau_{1},\tau_{2})} = k_{1}\tau_{1}\left(k_{2}\right)\mfa(\tau_{1},\tau_{2})\mfb(\tau_{1},\tau_{2})\cdot \Delta^{\mfc}_{\tau_{1}\tau_{2},1}$.
        Therefore, the left hand side is equal to
        \[
        \begin{aligned}
          & \left[\left(k_{1}\tau_{1}(k_{2})\mfa(\tau_{1},\tau_{2})\mfb(\tau_{1},\tau_{2})\cdot \Delta^{\mfa}_{\tau_{1}\tau_{2}, 1}\right)a \ox \Delta^{\mfb}_{\tau_{1}\tau_{2},1}b\right] \\
          =& \left[
             k_{1}\tau_{1}(k_{2})\mfa(\tau_{1},\tau_{2})\cdot \Delta^{\mfa}_{\tau_{1}\tau_{2},1} a \ox \mfb(\tau_{1},\tau_{2})\Delta^{\mfb}_{\tau_{1}\tau_{2}, 1}b
             \right]\\
          =& \left[
             \Delta^{\mfa}_{\tau_{1},k_{1}}\Delta^{\mfa}_{\tau_{2},k_{2}} a \ox
             \Delta^{\mfb}_{\tau_{1},1}\Delta^{\mfb}_{\tau_{2},1}b
             \right];
        \end{aligned}
        \]
        and the right hand side is also equal to
        \[
        \begin{aligned}
        &\left(k_{1}\cdot \Delta^{\mfc}_{\tau_{1}, 1}\right)
        \left[
        k_{2}\cdot\Delta^{\mfa}_{\tau_{2},1}a \ox \Delta^{\mfb}_{\tau_{2}, 1} b
          \right] \\
          =& \left[
             \left(k_{1}\cdot \Delta^{\mfa}_{\tau_{1}, 1}\right)\left(k_{2}\cdot\Delta^{\mfa}_{\tau_{2}, 1}\right)a \ox
             \Delta^{\mfb}_{\tau_{1},1}\Delta^{\mfb}_{\tau_{2},1}b
             \right]\\
          =&\left[
           k_{1}\tau_{1}(k_{2})\cdot \Delta^{\mfa}_{\tau_{1}\tau_{2},\mfa(\tau_{1},\tau_{2})}a \ox \Delta^{\mfb}_{\tau_{1}, 1}\Delta^{\mfb}_{\tau_{2},1}b
             \right]\\
          =&\left[
             \left(k_{1}\cdot \Delta^{\mfa}_{\tau_{1}, 1}\right)\left(k_{2}\cdot\Delta^{\mfa}_{\tau_{2}, 1}\right)a \ox
             \Delta^{\mfb}_{\tau_{1},1}\Delta^{\mfb}_{\tau_{2},1}b
             \right].
          \end{aligned}
        \]
\end{itemize}

Expanding everything out and checking on the basic elements, we see that for any $x \in \left(\cross_{\mfa}\ox_{F}\cross_{\mfb}\right)^{\opp}$, $y \in \cross_{\mfc}$ and $z \in M$, $x \cdot y \cdot z = y \cdot x \cdot z$. In another word, we gave $M$ a $\left(\cross_{\mfc}, \cross_{\mfa} \ox_{F} \cross_{\mfb}\right)$-bimodule structure.
  \leanok
  \lean{map_mul_proof.M, map_mul_proof.Aox_FB_op_tmul_smul_mk_tmul, map_mul_proof.C_smul_calc}
  % \uses{lem:cross-product-basis-conj}
\end{construction}


\begin{lemma}\label{lem:M-iso-tensor-over-K}
  $M$ is isomorphic to $\cross_{\mfa}\ox_{K}\cross_{\mfb}$ as $F$-modules.
  \leanok
  \lean{map_mul_proof.Aox_KBEquivM}
\end{lemma}
\begin{proof}
  The map $M \to \cross_{\mfa}\ox_{K}\cross_{\mfb}$ is obtained by descending the obvious $F$-linear map $\cross_{\mfa}\ox_{F}\cross_{\mfb}\to\cross_{\mfa}\ox_{K}\cross_{\mfb}$.
  By universal property of tensor product, there is an additive group homomorphism $\cross_{\mfa}\ox_{K}\cross_{\mfb} \to M$ given by $a \ox b \mapsto [a\ox b]$, this map is in fact $F$-linear. The two maps are inverse to each other.
\end{proof}

\begin{corollary}\label{cor:dim-finite-M}
  The $F$-dimension of $M$ is equal to ${\left(\dim_{F}K\right)}^{3}$, consequently $M$ is a finitely generated $\cross_{\mfc}$-module.
  \leanok
  \lean{map_mul_proof.M_F_dim}
  % \uses{lem:M-iso-tensor-over-K, lem:cross-product-basis}
\end{corollary}
\begin{proof}
  By~\cref{lem:M-iso-tensor-over-K}, the dimension of $M$ is equal to $\dim_{F}\cross_{\mfa}\ox_{K}\cross_{\mfb} = \dim_{K}\cross_{\mfa}\ox_{K}\cross_{\mfb}\dim_{F}K$. By~\cref{lem:cross-product-basis}, $\dim_{K}\cross_{\mfa} = \dim_{K}\cross_{\mfb} = \dim_{F}K$.
\end{proof}

\begin{construction}\label{con:cross-product-mul-iso}
  By~\cref{lem:direct-sum-simple-mod}, there exists some simple $\cross_{\mfc}$-module $S$ such that $\cross_{\mfc}$ is isomorphic to $\bigoplus_{i\in J}S$ as $\cross_{\mfc}$-module for some indexing set $J$. If we give $S$ the $F$-module structure by pulling back the $\cross_{\mfc}$-module structure, by restricting scalars $\cross_{\mfc}$ is isomorphic to $\bigoplus_{i\in J} S$ as $F$-module as well. Since $\cross_{\mfc}$ is a finite dimensional $F$-vector space, $J$ must be finite as well. Note that $S$ must be a finite dimensional $F$-vector space, because $S$ is finitely generated as $\cross_{\mfc}$-module and $\cross_{\mfc}$ has finite $F$-dimension.
  The indexing set $J$ must be nonempty, otherwise $\cross_{\mfc}$ being isomorphic to $\bigoplus_{\emptyset} S$ is a trivial ring; but simple rings are non-trivial. Since $J$ is finite, direct sum over $J$ and direct product over $J$ agree.
  Recall~\cref{con:self-tensor-opp-iso-end} and~\cref{con:end-vec-iso-matrix}, for all non-zero $m \in\mathbb{N}$, we have
  \[
    \End_{\cross_{\mfc}}\left(S^{m}\right) \cong \mat_{m}\left(\End_{\cross_{\mfc}}S\right)
  \]
  as $F$-algebras, hence
  \[
    \cross_{\mfc}^{\opp} \cong \End_{\cross_{\mfc}}\cross_{\mfc}\cong \End_{\cross_{\mfc}}S^{|J|} \cong \mat_{|J|}\left(\End_{\cross_{\mfc}}S\right)
  \]
  as $F$-algebras.
  Finally
  \[
    \cross_{\mfc}\cong \mat_{|J|}\left(\End_{\cross_{\mfc}}\left(S\right)^{\opp}\right)
  \]
  as $F$-algebras.
  \leanok
  \lean{map_mul_proof.exists_simple_module_directSum,map_mul_proof.isoDagger, map_mul_proof.mopEquivEnd', map_mul_proof.C_iso}
  % \uses{lem:direct-sum-simple-mod, con:self-tensor-opp-iso-end, con:end-vec-iso-matrix}
\end{construction}

\begin{corollary}\label{cor:dim-eq-1}
  \[
    \left(\dim_{F}K\right)^{2} = |J|^{2}\dim_{F}\End_{\cross_{\mfc}}S.
  \]
  \leanok
  \lean{map_mul_proof.dim_endCSM}
  % \uses{con:cross-product-mul-iso}
\end{corollary}
\begin{proof}
  They are both equal to $\dim_{F}\cross_{\mfc}$ by~\cref{con:cross-product-mul-iso}.
\end{proof}

\begin{corollary}\label{cor:dim-eq-2}
  \[
    |J|\dim_{F}S = \left(\dim_{F}K\right)^{2}
  \]
  \leanok
  \lean{map_mul_proof.SM_F_dim}
  % \uses{con:cross-product-mul-iso}
\end{corollary}
\begin{proof}
  They are all equal to $\dim_{F}\cross_{\mfc} = \dim_{F}S^{|J|}$ by~\cref{con:cross-product-mul-iso}.
\end{proof}

\begin{lemma}\label{lem:M-iso-simple-pow}
  There exists a $\cross_{\mfc}$-linear isomorphism between $M$ and $S^{|J|\dim_{F}K}$.
  % Consequently there also exists an $F$-linear isomorphism between $M$ and $S^{|J|\dim_{F}K}$ and $F$-algebra isomorphism between $\End_{\cross_{\mfc}}M$ and $\End_{\cross_{\mfc}}\left(S^{|J|\dim_{F}K}\right)\cong\mat_{|J|\dim_{F}K}\left(S\right)$.
  \leanok
  \lean{map_mul_proof.M_iso_pow}
  % \uses{lem:iso-of-dim-eq, cor:dim-finite-M, cor:dim-eq-2}
\end{lemma}

\begin{proof}
  By~\cref{lem:iso-of-dim-eq}, we only need to show that $\dim_{F}M = \dim_{F}S^{|J|\dim_{F}K}$. We already have $\dim_{F}M = \left(\dim_{F}K\right)^{3}$ by~\cref{cor:dim-finite-M}. We also have $\dim_{F}S^{|J|\dim_{F}K} = |J|\dim_{F}K\dim_{F}S = \dim_{F}K \left(|J|\dim_{F}S\right) = \dim_{F}K\left(\dim_{F}K\right)^{2}$ by~\cref{cor:dim-eq-2}.
\end{proof}

\begin{corollary}\label{cor:M-iso-simple-pow}
  As $F$-vector spaces, $M \cong S^{|J|\dim_{F}K}$.
  \leanok
  \lean{map_mul_proof.M_iso_pow'}
  % \uses{lem:M-iso-simple-pow}
\end{corollary}
\begin{proof}
  Restricting scalars on the $\cross_{\mfc}$-linear isomorphism in~\cref{lem:M-iso-simple-pow}
\end{proof}

\begin{corollary}\label{cor:end-M-iso-mat}
  As $F$-algebras, $\End_{\cross_{\mfc}}M \cong \mat_{|J|\dim_{F}K}(\End_{\cross_{\mfc}}S)$.
  \leanok
  \lean{map_mul_proof.endCMIso'}
  % \uses{cor:M-iso-simple-pow, con:end-vec-iso-matrix}
\end{corollary}
\begin{proof}
  From~\cref{cor:M-iso-simple-pow}, we have $\End_{\cross_{\mfc}}M\cong \End_{\cross_{\mfc}}\left(S^{|J|\dim_{F}K}\right)$. By~\cref{con:end-vec-iso-matrix}, they are isomorphic to $\mat_{|J|\dim_{F}K}(\End_{\cross_{\mfc}}S)$.
\end{proof}

\begin{corollary}\label{cor:dim-eq-3}
  \[
    \dim_{F}\End_{\cross_{\mfc}}M = {\left(\dim_{F}K\right)}^{4}.
  \]
  \leanok
  \lean{map_mul_proof.dim_endCM}
  % \uses{cor:dim-eq-1}
\end{corollary}

\begin{proof}
  \[
    \begin{aligned}
      \dim_{F}\End_{\cross_{\mfc}M}
      &= \dim_{F}\mat_{|J|\dim_{F}K}(\End_{\cross_{\mfc}}S) \\
      &= \dim_{F}\left(\End_{\cross_{\mfc}}S\ox_{F}\mat_{|J|\dim_{F}K}(F)\right)\\
      &= |J|^{2}\left(\dim_{F}K\right)^{2}\dim_{F}\End_{\cross_{\mfc}}S\\
      &= \left(\dim_{F}K\right)^{2}\left(|J|^{2}\dim_{F}\End_{\cross_{\mfc}}S\right) \\
      &= \left(\dim_{F}K\right)^{2}\left(\dim_{F}K\right)^{2},
    \end{aligned}
  \]
  where the last equality is by~\cref{cor:dim-eq-1}.
\end{proof}

\begin{theorem}\label{thm:cross-mul}
  The cross product $\cross_{\mfc}$ and the tensor product $\cross_{\mfa}\ox_{F}\cross_{\mfb}$ are Brauer equivalent.
  \leanok
  \lean{map_mul_proof.isBrauerEquivalent}
  % \uses{cor:dim-eq-3, cor:alghom-bijective-of-dim-eq, cor:end-M-iso-mat, con:cross-product-mul-iso}
\end{theorem}

\begin{proof}
  We define an $F$-algebra homomorphism $\phi : \left(\cross_{\mfa}\ox_{F}\cross_{\mfb}\right)^{\opp} \to \End_{\cross_{\mfc}}M$ by
  $x \mapsto (x \cdot \bullet)$. By~\cref{cor:dim-eq-3}, both sides has $F$-dimension $\left(\dim_{F}K\right)^{4}$, therefore, $\phi$ is an $F$-algebra isomorphism by~\cref{cor:alghom-bijective-of-dim-eq}. Hence we have another $F$-algebra isomorphism by composing the isomorphism in~\cref{cor:end-M-iso-mat}:
  \[
    \phi^{\opp}:\cross_{\mfa}\ox_{F}\cross_{\mfb}\cong \left(\End_{\cross_{\mfc}}M\right)^{\opp}\cong \mat_{|J|\dim_{F}K}\left(\left(\End_{\cross_{\mfc}}S\right)^{\opp}\right).
  \]
  In the meantime, by~\cref{con:cross-product-mul-iso}, we have $\cross_{\mfc}\cong\mat_{|J|}\left(\left(\End_{\cross_{\mfc}}S\right)\opp\right)$; hence $\mat_{\dim_{F}K}\left(\cross_{\mfc}\right)$ is isomorphic to $\cross_{\mfa}\ox_{F}\cross_{\mfb}$.
\end{proof}

\begin{corollary}[group isomorphism]
  For a finite dimensional Galois field extension $K/F$, the relative Brauer group $\br(K/F)$ is isomorphic to the second group cohomology $\HH^{2}\left(\gal(K/F),K^{\star}\right)$.
  \leanok
  \lean{RelativeBrGroup.isoSnd}
  % \uses{cor:relative-br-2nd-coh-bij, thm:cross-mul}
\end{corollary}
\begin{proof}
  In~\cref{cor:relative-br-2nd-coh-bij}, we have seen that $\HH^{2}$ and $\cross$ form a bijection, thus it is sufficient to check either one of them preserves multiplication. The function $\cross: \HH^{2}\left(\gal(K/F), K^{\star}\right)$ preserves multiplication: let $[\mfa], [\mfb]$ be two elements in $\HH^{2}\left(\gal(K/F), K^{\star}\right)$, by~\cref{thm:cross-mul}$, \cross(\mfa\mfb)$ is indeed Brauer equivalent to $\cross(\mfa)\ox_{F}\cross(\mfb)$ that is
  \[
    \left[\cross_{\mfa}\right]_{\sim_{\br}}\left[\cross_{\mfb}\right]_{\sim_{\br}} = \left[\cross_{\mfa\mfb}\right].
  \]
\end{proof}



%%% Local Variables:
%%% mode: LaTeX
%%% TeX-master: "../print"
%%% End:


%%FILE blueprint/src/content/central_simple.tex
\chapter{Central Simple Algebras}\label{sec:csa}

\section{Basic Theory}

In this chapter we define central simple algebras. We used some results
in~\cref{sec:unclassified}.

\begin{definition}[Simple Ring] A ring $R$ is simple if the only
  two-sided-ideals of $R$ are ${0}$ and $R$. An algebra is simple if it is
  simple as a ring.
  \leanok
  \lean{IsSimpleRing}
\end{definition}

\begin{remark}
  Division rings are simple.
\end{remark}

\begin{lemma}
  \label{lem:center-simple-ring}
  Let $A$ be a simple ring, then centre of $A$ is a field. \leanok
  \lean{IsSimpleRing.isField_center}
\end{lemma}
\begin{proof}
  Let $0\ne x$ be an element of centre of $A$. Then $I := \{xy | y\in A\}$ is a
  two-sided-ideal of $A$. Since $0\ne x\in I$, we have that $I = A$. Therefore
  $1 \in I$, hence $x$ is invertible.
\end{proof}

\begin{definition}[Central Algebras]
  Let $R$ be a ring and $A$ an $R$-algebra, we say $A$ is central if and only if
  the centre of $A$ is $R$
  \leanok
  \lean{Algebra.IsCentral}
\end{definition}

\begin{remark}
  Every commutative ring is a central algebra over itself.
\end{remark}

\begin{remark}
  Simpleness is invariant under ring isomorphism and centrality is invariant
  under algebra isomorphism.
\end{remark}

\begin{lemma}
  \label{lem:opp-central}
  If $A$ is a central $R$-algerba, $A^{\opp}$ is also central.
  \leanok
  \lean{CSA_op_is_CSA}.
\end{lemma}

\begin{lemma}\label{lem:simple-ring-iff}
  $R$ is a simple ring if and only if any ring homomorphism $f : R \to S$ either
  injective or $S$ is the trivial ring. \leanok
  \lean{IsSimpleRing.iff_injective_ringHom_or_subsingleton_codomain}
\end{lemma}
\begin{proof}
  If $R$ is simple, then the $\ker f$ is either $\{0\}$ or $R$. The former case
  implies that $f$ is injective while the latter case implies that $S$ is the
  trivial ring. Conversely, let $I\subseteq R$ be a two-sided-ideal. Consider
  $\pi: R \to {}^{R}/_{I}$, either $\pi$ is injective implying that $I = \{0\}$
  or that ${}^{R}/_{I}$ is the trivial ring implying that $I = R$.
\end{proof}
\begin{remark}
  If $A$ is a simple $R$-algebra, ``ring homomorphism'' in \cref{lem:simple-ring-iff} can be replaced with $R$-algebra homomorphism.
\end{remark}

\begin{corollary}\label{cor:alghom-bijective-of-dim-eq}
  Assume $R$ is a field. Let $A, B$ be finite dimensional $R$-algebras where $A$
  is simple as well. Then any $R$-algebra homomorphism $f:A\to B$ is bijective
  if $\dim_{R}A=\dim_{R}B$. 
  \leanok
  \lean{bijective_of_dim_eq_of_isCentralSimple}
  % \uses{lem:simple-ring-iff}
\end{corollary}
\begin{proof}
  By~\cref{lem:simple-ring-iff}, $f$is injective. Then
  $\dim_{K}\im f = \dim_{K} B - \dim_{K}\ker f = \dim_{K}B$ meaning that $f$ is
  surjective.
\end{proof}

Let $K$ be a field and $A, B$ be $K$-algebras.
\begin{lemma}
  \label{lem:tensor-central}
  If $A$ and $B$ are central $K$-algebras, $A\otimes_{K}B$ is a central
  $K$-algebra as well.
  \leanok
  \lean{IsCentralSimple.TensorProduct.isCentral}
  % \uses{cor:center-tensor-center}
\end{lemma}
\begin{proof}
  Assume $A$ and $B$ are central algebras, then
  by~\cref{cor:center-tensor-center}
  $Z\left(A\otimes_{R}B\right)=Z\left(A\right)\otimes_{R}Z\left(B\right)=R\otimes_{R}R=R$.
\end{proof}

\begin{theorem}
  \label{thm:tensor-csa}
  If $A$ is a simple $K$-algebra and $B$ is a central simple $K$-algebra,
  $A\otimes_{K}B$ is a central simple $K$-algebra as well. 
  \leanok
  \lean{IsCentralSimple.TensorProduct.simple}
  % \uses{lem:tensor-central, lem:expand-tensor-in-basis}
\end{theorem}
\begin{proof}
  By~\cref{lem:tensor-central}, we need to prove $A\otimes_{K}B$ is a simple
  ring. Denote $f$ as the map $A\to A\otimes_{K}B$. It is sufficient to prove
  that for every two-sided-ideal $I\subseteq A\otimes_{K}B$, we have
  $I = \left\langle f\left(f^{-1}\left(I\right)\right)\right\rangle$. Indeed,
  since $A$ is simple $f^{-1}\left(I\right)$ is either $\left\{0\right\}$ or
  $A$, if it is $\left\{0\right\}$, then $I=\left\{0\right\}$; if it is $A$,
  then $I$ is $A$ as well.

  We will prove that
  $I \le \left\langle f\left(f^{-1}\left(I\right)\right)\right\rangle$, the
  other direction is straightforward. Without loss of generality assume
  $I\ne\left\{0\right\}$. Let $\McA$ be an arbitrary basis of $A$, by
  \cref{lem:expand-tensor-in-basis}, we see that every element
  $x \in A\otimes_{K}B$ can be written as $\sum_{i=0}^{n}\McA_{i}\otimes b_{i}$
  for some natrual number $n$ and some choice of $b_{i}\in B$ and
  $\McA_{i}\in \McA$. Since $I$ is not empty, we see there exists a non-zero
  element $\omega\in I$ such that its expansion
  $\sum_{i=0}^{n}\McA_{i}\otimes b_{i}$ has the minimal $n$. In particular, all
  $b_{i}$ are non-zero and $n\ne0$. We have
  $\omega=\McA_{0}\otimes b_{0}+\sum_{i=1}^{n}\McA_{i}\otimes b_{i}$. Since $B$
  is simple, $1 \in B = \left\langle\langle b_{0} \right\rangle$; hence we write
  $1\in\sum_{j=0}^{m}x_{i}b_{0}y_{i}$ for some $x_{i},y_{i}\in B$. Define
  $\Omega := \sum_{j=0}^{m}(1\otimes x_{i})\omega(1\otimes y_{i})$ which is also
  in $I$. We write
  \[
    \begin{aligned}
      \Omega &= \McA_{0} \otimes \left(\sum_{j=0}^{m}x_{j}b_{0}y_{j}\right) + \sum_{i=1}^{n}\McA_{i}\otimes\left(\sum_{j=0}^{m}x_{j}b_{i}y_{j}\right) \\
             &= \McA_{0}\otimes 1 + \sum_{i=1}^{n}\McA_{i}\otimes\left(\sum_{j=0}^{m}x_{j}b_{i}y_{j}\right)
    \end{aligned}
  \]
  For every $\beta\in B$, we have that
  $\left(1\otimes \beta\right)\Omega - \Omega\left(1\otimes \beta\right)$ is in
  I and is equal to
  \[
    \sum_{i=1}^{n}\McA_{i}\otimes\left(\sum_{j=0}^{m}\beta x_{j}b_{i}y_{j}-x_{j}b_{i}y_{j}\beta\right),
  \]
  which is an expansion of $n-1$ terms, thus
  $\left(1\otimes\beta\right)\Omega-\Omega\left(1\otimes\beta\right)$ must be
  $0$. Hence we conclude that for all $i=1,\dots,n$,
  $\sum_{j=0}^{m}x_{j}b_{i}y_{j}\in Z\left(B\right)=K$. Hence for all
  $i=1,\dots,n$, we find a $\kappa_{i}\in K$ such that
  $\kappa_{i}=\sum_{j=0}^{m}x_{j}b_{i}y_{j}$. Hence we can calculate $\Omega$ as
  \[
    \begin{aligned}
      \Omega &= \McA_{0}\otimes 1 + \sum_{i=1}^{n}\McA_{i}\otimes\left(\sum_{j=0}^{m}\right) \\
             &= \McA_{0}\otimes 1 + \sum_{i=1}^{n}\McA_{i}\otimes\kappa_{i} \\
             &= \left(\McA_{0}+\sum_{i=1}^{n}\kappa_{i}\cdot\McA_{i}\right) \otimes 1
    \end{aligned}.
  \]
  From this, we note that
  $\McA_{0}+\sum_{i}^{n}\kappa_{i}\cdot\McA_{i}\in f^{-1}\left(I\right)$; since
  $A$ is simple, we immediately conclude that $f^{-1}\left(I\right) = A$, once
  we know $\McA_{0}+\sum_{i=1}^{n}\kappa_{i}\cdot \McA_{i}$ is not zero. If it
  is zero, by the fact that $\McA$ is a linearly independent set, we conclude
  that $1,\kappa_{1},\dots,\kappa_{n}$ are all zero; which is a contradiction.
  Since $f^{-1}\left(I\right) = A$, we know
  $\left\langle f\left(f^{-1}I\right)\right\rangle = A \otimes_{K} B$.
\end{proof}

\begin{corollary}\label{cor:csa-basechange}
  Central simple algebras are stable under base change. That is, if $L/K$ is a
  field extension and $D$ is a central simple algebra over $K$, then
  $L\otimes_{K} D$ is central simple over $L$. \leanok
  \lean{IsCentralSimple.baseChange}
  % \uses{thm:tensor-csa, cor:center-tensor-center}
\end{corollary}

\begin{proof}
  By~\cref{thm:tensor-csa}, $L\otimes_{K} D$ is simple. Let
  $x\in Z\left(L\otimes_{K}D\right)$, by~\cref{cor:center-tensor-center}, we
  have $x\in Z\left(L\right)\otimes Z\left(D\right)=Z\left(L\right)$. Without
  loss of generality, we can assume that $x = l \otimes d$ is a pure tensor,
  then $l \in Z\left(L\right)$ and $d \in K$. Therefore $x = d\cdot l \in L$.
\end{proof}

\begin{theorem}\label{thm:tensor-simple-implies-simple}
  If $A\otimes_{K} B$ is a simple ring, then $A$ and $B$ are both simple.
  \leanok
  \lean{IsCentralSimple.TensorProduct.simple}
  % \uses{lem:simple-ring-iff}
\end{theorem}
\begin{proof}
  By symmetry, we only prove that $A$ is simple.
  If $A$ or $B$ is the trivial ring then $A\otimes_{K} B$ is the trivial ring, a contradiction. Thus we assume both $A$ and $B$ are non-trivial. Suppose $A$ is not simple, by~\cref{lem:simple-ring-iff}, there exists a non-trivial $K$-algebra $A'$ and a $K$-algebra homomorphism $f : A \to A'$ such that $\ker f \not= \{0\}$. Let $F : A \otimes_{K} B\to A'\otimes_{K} B$ be the base change of $f$, then since $A\otimes_{K} B$ is simple and $A'\otimes B$ is non-trivial ($A'$ is non-trivial and $B$ is faithfully flat because $B$ is free), we conclude that $F$ is injective. Then we have that
  \begin{center}
    \begin{tikzcd}
      0 \arrow{r}{0} & A \otimes_{K} B \arrow{r}{F} & A'\otimes_K B
    \end{tikzcd}
  \end{center}
  is exact. Since $B$ is faithfully flat as a $K$-module, tensorig with $B$ reflects exact sequences, therefore
  \begin{center}
    \begin{tikzcd}
      0 \arrow{r}{0} & A \arrow{r}{f} & A'
    \end{tikzcd}
  \end{center}
  is exact as well. This is contradiction since $f$ is not injective.
\end{proof}

\section{Subfields of Central Simple Algebras}\label{sec:subfield}

\begin{definition}[Subfield]
  For any field $K$ and $K$-algebra $A$, a subfield $B \subseteq A$ is a commutative $K$-subalgebra of $A$ that is closed under inverse for any non-zero member.
  \leanok
  \lean{SubField}
\end{definition}

\begin{remark}
  Subfields inherit a natural ordering from subalgebras.
\end{remark}

Let $K$ be any field and $D$ a finite dimensional central division $K$-algebra and $A$ a finite dimensional central simple algebra of $A$.
\begin{lemma}\label{lem:dim-maximal-subfield}
  Let $k$ be a maximal subfield of $D$,
  \[
    \dim_{K}D = {\left(\dim_{K}k\right)}^{2}.
  \]
  \leanok
  \lean{dim_max_subfield}
  % \uses{lem:dim-centralizer}
\end{lemma}

\begin{proof}
  By~\cref{lem:dim-centralizer}, we have that $\dim_{K} D = \dim_{K}C_{D}(k)\cdot \dim_{K}k$. Hence it is sufficient to show that $C_{D}(k) = k$. By the commutativity of $k$, we have that $k \le C_{D}(k)$. Suppose $k \ne C_{D}(k)$: let $a \in C_{D}(k)$ that is not in $k$. We see that $L := k(a)$ is another subalgebra of $D$ that is strictly larger than $k$; a contradiction. Therefore $k = C_{D}(k)$ and the theorem is proved.
\end{proof}

\begin{lemma}\label{lem:tfae-subfield}
  Suppose $L$ is a subfield of $A$, the following are equivalent:
  \begin{enumerate}
    \item $L$ = $C_{A}(L)$
    \item $\dim_{K}A = {\left(\dim_{K}L\right)}^{2}$
    \item for any commutative $K$-subalgebra $L' \subseteq A$, $L \subseteq L'$ implies $L = L'$
  \end{enumerate}
  \leanok
  \lean{cor_two_1to2, cor_two_2to3, cor_two_3to1}
  % \uses{lem:dim-centralizer}
\end{lemma}
\begin{proof}
  We prove the following:
  \begin{itemize}
    \item ``1. implies 2.'': this is~\cref{lem:dim-centralizer}.
    \item ``2. implies 1.'': Since $L$ is commutative, we always have $L \subseteq C_{A}(L)$. Hence we only need to show $\dim_{K}L = \dim_{K}C_{A}(L)$. This is because by~\cref{lem:dim-centralizer}, we have that $\dim_{K}A = \dim_{K}L\cdot \dim_{K}C_{A}(L)$ and by 2. we have $\dim_{K}L\cdot \dim_{K}C_{A}(L)= \dim_{K}L\cdot \dim_{K}L$.
    \item ``2. implies 3.'': Since 2. implies 1., we assume $L = C_{A}(L)$, therefore all we need is to prove $L' \subseteq C_{A}(L)$. Let $x \in L'$ and $y \in L \subseteq L'$, we need to show $xy = yx$ which is commutativity of $L'$.
    \item ``3. implies 1.'': By commutativity of $L$, we always have $L \subseteq C_{A}(L)$. For the other direction, suppose $C_{A}(L)\not\subseteq L$, then there exists some $a \in C_{A}(L)$ but not in $L$. Consider $L' = L(a)$, by 3., we have $L' = L$ which is a contradiction.
  \end{itemize}
\end{proof}

%%% Local Variables:
%%% mode: LaTeX
%%% TeX-master: "../print"
%%% End:


%%FILE blueprint/src/content/double_centralizer.tex
\section{Double Centralizer Theorem}\label{sec:double-centralizer}

In this section let $F$ be a field and $A$ an $F$-algebra. Define
$\McL_{A} \subseteq \End_{F}A$ to be
\[\left\{f : A \to A|f(x)=ax~\text{for some}~a\in A\right\},\]
i.e. $F$-linear maps defined by left multiplication; similarly define
$\McR_{A}$. Note that $\McL_{A}$ and $\McR_A$ are $F$-subalgebras of
$\End_{F}A$. When we need to stree the underlying field is $F$, we also write
$\McL_{A}^{F}$ and $\McR_{A}^{F}$. We assume $A$ to be a finite dimensional
central simple $F$-algebra.
\begin{lemma}
  \label{lem:centralizer-mul-left-le}
  The centralizer of $\McL_{A}$ in $\End_{F}A$ is smaller than or equal to
  $\McR_{A}$:
  \[
    C_{\End_{F}A}\left(\McL_{A}\right) \le \McR_{A}.
  \]
  \leanok
  \lean{centralizer_mulLeft_le_of_isCentralSimple}
  % \uses{con:self-tensor-opp-iso-end}
\end{lemma}

\begin{proof}
  Indeed, let $x \in C_{\End_{F}A}\left(\McL_{A}\right)$. Recall
  from~\cref{con:self-tensor-opp-iso-end} that
  $e : A\otimes_{F}A^{\opp}\cong \End_{F}A$ as $F$-algebras. Then $e^{-1}(x)$ is
  in
  $C_{A\otimes_{F}A^{\opp}}\left(\im\left(A\to A\otimes_{F}A^{\opp}\right)\right)$
  (for $e$ sends $a\otimes 1$ to the $F$-linear map $(a\cdot\bullet)$). Since
  $C_{A\otimes_{F}A^{\opp}}\left(\im\left(A\to A\otimes_{F}A^{\opp}\right)\right) = Z(A)\otimes_{F}A^{\opp}=F\otimes_{F}A^{\opp}=\im\left(A^{\opp}\to A\otimes_{F}A^{\opp}\right)$,
  we find some $y\in A^{\opp}$ such that $1 \otimes y = e^{-1}(x)$. Therefore
  $e\left(1\otimes y\right) = x$; but $e\left(1\otimes y\right)$ is in
  $\McR_{A}$ for it is the linear map $(\bullet\cdot y)$.
\end{proof}

\begin{remark}\label{rem:mul-left-center-linear}
  For any $F$-algebra $B$, every element in $C_{\End_{F}B}\left(\McL_{B}\right)$
  is in fact $Z(B)$-linear. Let $x\in C_{\End_{F}B}\left(\McL_{B}\right)$,
  $z\in Z(B)$ and $b \in B$, we have $x(z\cdot b) = z\cdot x(b)$ because $x$
  commutes with the linear map $\left(z\cdot\bullet\right)$.
\end{remark}

\begin{remark}
  $A$ is a $Z(A)$-algebra whose algebra structure is given by
  $Z(A)\hookrightarrow A$. By~\cref{lem:center-simple-ring}, $Z(A)$ is a field.
  $A$ is finite dimensional as a $Z(A)$-module because of the tower $A/Z(A)/F$.
\end{remark}

 \begin{lemma}\label{lem:mul-right-iso-opp}
   As $F$-algebras, we have $\McR_{A}\cong A^{\opp}$. \leanok
   \lean{Module.End.rightMulEquiv}
 \end{lemma}
 \begin{proof}
   We prove the map $A^{\opp} \to \McR_{A}$ is bijective. It is injective
   because if $\left(\bullet \cdot a\right) = \left(\bullet\cdot b\right)$, then
   $a = 1 \cdot a = 1 \cdot b = b$. The map is surjective by the definition of
   $\McR_{A}$.
 \end{proof}
 \begin{lemma}\label{lem:centralizer-mul-left-eq-mul-right}
   Let $B$ be any simple $F$-algebra ({\em not\/} necessarily central). The
   centralizer of $\McL_{B}$ in $\End_{F}B$ is equal to $\McR_{B}$.
   \leanok
   \lean{centralizer_mulLeft}
   % \uses{lem:centralizer-mul-left-le, rem:mul-left-center-linear}
 \end{lemma}
 \begin{proof}
   It is straightforward to show
   $\McR_{B}^{F}\le C_{\End_{F}A}\left(\McL_{B}^{F}\right)$. So we only need to
   prove $C_{\End_{F}A}\left(\McL_{B}^{F}\right)\le \McR_{B}^{F}$.
   By~\cref{lem:centralizer-mul-left-le}, since $B$ is a central simple finite
   dimensional $Z(B)$-algebra, we have that
   \[
     C_{\End_{Z(B)}B}\left(\McL_{B}^{Z(B)}\right) \le \McR_{B}^{Z(B)}.
   \]
   Suppose $f\in\End_{F}B$ is in $C_{\End_{F}B}B$, by~\cref{rem:mul-left-center-linear}, $f$ is $Z(B)$-linear as well. Then $f$ is
   in $\McR_{B}^{Z(B)}$; that is $f$ is equal to $\left(\bullet\cdot b\right)$
   for some $b \in B$ as $Z(B)$-linear maps. Then $f$ is also equal to
   $\left(\bullet\cdot b\right)$ as $F$-linear maps.
 \end{proof}

 \begin{construction}
   Let $B$ be any $F$-algebra and $S\subseteq B$ an $F$-subalgebra. For any
   $x\in B^{\times}$, we have that $x S x^{-1}:=\{xsx^{-1}|s\in S\}$ is an
   $F$-subalgebra of $B$ as well. We have the obvious $F$-algebra isomorphism
   $S \cong xSx^{-1}$ given by $s\mapsto xsx^{-1}$ and $x^{-1}tx\mapsfrom t$.
   Therefore $\dim_{F}S=\dim_{F}xSx^{-1}$ and $S$ is a simple ring if and only
   if $xSx^{-1}$ is a simple ring.
   \leanok
   \lean{Subalgebra.conj}
 \end{construction}

 \begin{lemma}
   \label{lem:centralizer-conj}
   Let $B$ be any $F$-algebra, $x \in B^{\times}$ and $S\subseteq B$ be an
   $F$-subalgebra of $B$, then $C_{B}(xSx^{-1})=x\left(C_{B}(S)\right)x^{-1}$.
   \leanok
   \lean{Subalgebra.conj_centralizer}
 \end{lemma}
 \begin{proof}
   If $a \in C_{B}\left(xSx^{-1}\right)$, then $x^{-1}ax$ is in $C_{B}(S)$.
   Conversely if $a$ is equal to $xbx^{-1}$ with $b\in C_{B}(S)$, then it is in
   $C_{B}\left(xSx^{-1}\right)$ as well.
 \end{proof}

 \begin{remark}
   For any finite dimensional $F$-module $B$, we have isomorphism
   $\End_{F}B \cong \Mat_{\dim_{F}B}F$ as $F$-algebras. Hence $\End_{F} B$ is a
   finite-dimensional central simple algebra over $F$.
 \end{remark}

 \begin{lemma}\label{lem:tensor-mul-right-simple}
   Let $S\subseteq A$ be a simple $F$-subalgebra, then
   $A \otimes_{F} \McR_{S}$ is a simple ring.
   \leanok
   % \uses{thm:tensor-csa, lem:mul-right-iso-opp}
 \end{lemma}

 \begin{proof}
   By~\cref{lem:mul-right-iso-opp}, we have
   $A\otimes \McR_{S}\cong A\otimes S^{\opp}$ as $F$-algebras. The claim follows
   from~\cref{thm:tensor-csa}.
 \end{proof}

 \begin{lemma}\label{lem:centralizer-tensor-end-eq-conj-tensor-mul-right}
   Let $S\subseteq A$ be a simple $F$-subalgebra, then there
   exists an $x\in \left(A\otimes_{F}\End_{F}S\right)^{\times}$ such that
   $C_{A}(S)\otimes_{F}\End_{F}S$ is isomorphic to
   $x\left(A\otimes_{F}\McR_{S}\right)x^{-1}$ as $F$-algebras.
   \leanok
   \lean{centralizer_isSimple.aux.step1}
   % \uses{lem:tensor-central,thm:tensor-csa, thm:skolem-noether, lem:centralizer-conj, lem:centralizer-inclusion-tensor,lem:centralizer-mul-left-eq-mul-right}
 \end{lemma}

 \begin{proof}
   By~\cref{lem:tensor-central} and~\cref{thm:tensor-csa},
   $A \otimes_{F}C_{A}(S)$ is a central simple $F$-algebra. Let
   $f : S \to A\otimes_{F}\End_{F}S$ be an $F$-algebra homomorphism defined by
   $s \mapsto s \otimes \mathbf{1}_{S}$ and $g : S \to A \otimes_{F}\End_{F}S$
   be an $F$-algebra homomorphism defined by
   $\mathbf{1}_{A} \otimes \left(s\cdot\bullet\right)$. Then
   by~\cref{thm:skolem-noether}, we that there exists some
   $x\in \left(A\otimes_{F}\End_{F}S\right)^{\times}$ such that $f = xgx^{-1}$.
   Then we have $S\otimes_{F}\End_{F}S$ is equal to
   $x\left(A\otimes_{F}\McL_{S}\right)x^{-1}$: indeed the left hand side is
   $\im f$ while the right handside is $x\left(\im g\right)x^{-1}$. Therefore
   $C_{A\otimes_{F}\End_{F}S}\left(S\otimes_{F}\End_{F}S\right) = C_{A\otimes_{F}\End_{F}S}\left(x\left(A\otimes_{F}\McR_{S}\right)x^{-1}\right)$.
   By~\cref{lem:centralizer-conj}, the right hand side is equal to
   $xC_{A\otimes_{F}\End_{F}S}\left(A\otimes_{F}\McL_{S}\right)x^{-1}$ which is
   $x\left(A\otimes_{F}C_{\End_{F}S}\left(\McL_{S}\right)\right)x^{-1}$
   by~\cref{lem:centralizer-inclusion-tensor} which is
   $x\left(A\otimes_{F}\McR_{S}\right)x^{-1}$
   by~\cref{lem:centralizer-mul-left-eq-mul-right}.
 \end{proof}

 \begin{lemma}\label{lem:simple-centralizer}
   Let $S\subseteq A$ be a simple $F$-subalgebra, then
   $C_{A}(S)$ is simple as well.
   \leanok
   \lean{centralizer_isSimple}
   % \uses{lem:centralizer-tensor-end-eq-conj-tensor-mul-right, lem:tensor-mul-right-simple, thm:tensor-simple-implies-simple}
 \end{lemma}

 \begin{proof}
   By~\cref{lem:centralizer-tensor-end-eq-conj-tensor-mul-right},
   $C_{A}(S) \otimes_{F}\End_{F}S$ is isomorphic to
   $x\left(A\otimes_{F}\McR_{S}\right)x^{-1}$ as $F$-algebras. Then
   $C(S)\otimes_{F}\End_{F}S$ is simple since $A\otimes_{F}\McR_{S}$ is simple
   by~\cref{lem:tensor-mul-right-simple}.
   By~\cref{thm:tensor-simple-implies-simple}, $C_{A}(S)$ is simple.
 \end{proof}

 \begin{lemma}
   \label{lem:dim-centralizer}
   Let $S\subseteq A$ be a simple $F$-subalgebra. Then
   \[
     \dim_{F}C_{A}(S)\cdot\dim_{F}S=\dim_{F}A.
   \]
   \leanok
   \lean{dim_centralizer}
   % \uses{lem:centralizer-tensor-end-eq-conj-tensor-mul-right, lem:mul-right-iso-opp}
 \end{lemma}
 \begin{proof}
   By~\cref{lem:centralizer-tensor-end-eq-conj-tensor-mul-right},
   $C_{A}(S) \otimes_{F}\End_{F}S$ is isomorphic to
   $x\left(A\otimes_{F}\McR_{S}\right)x^{-1}$ as $F$-algebras. Hence $\dim_{F}\left(C_{A}(S)\otimes_{F}\End_{F}S\right)=\dim_{F}\left(A\otimes_{F}\McR_{S}\right)$ where the left hand side is $\dim_{F}C_{A}(S)\cdot \dim_{F}\End_{F}S$ and the right hand side is $\dim_{F}A\cdot\dim_{F}\McR_{S}$. Since $\dim_{F}\End_{F}S=\dim_{F}S^{2}$ and $\dim_{F}\McR_{S}=\dim S$ (by~\cref{lem:mul-right-iso-opp}), we proved this lemma.
 \end{proof}

\begin{corollary}
  \label{cor:self-tensor-centralizer}
  Let $S \subseteq A$ be a central simple $F$-subalgebra,
  \[
    A \cong B \ox_{F} C_{A}(B).
  \]
  \leanok
  \lean{writeAsTensorProduct}
  % \uses{lem:simple-centralizer, thm:tensor-csa, cor:alghom-bijective-of-dim-eq, lem:dim-centralizer}
\end{corollary}

\begin{proof}
  By~\cref{lem:simple-centralizer}, $C_{A}(B)$ is simple and by~\cref{thm:tensor-csa}, $B \ox_{F} C_{A}(B)$ is simple. Hence the map $B \ox_{F} C_{A}(B) \to A$ induced by $B \hookrightarrow A$ and $C_{A}(B) \hookrightarrow A$ is injective. By~\cref{cor:alghom-bijective-of-dim-eq}, we only need to show $\dim_{F}B\ox_{F}C_{A}(B) =\dim_{F}A$ which is precisely~\cref{lem:dim-centralizer}.
\end{proof}

 \begin{theorem}[Double Centralizer]\label{thm:double-centralizer}
   Let $S\subseteq A$ be a simple $F$-subalgebra, we have
   \[
     C_{A}\left(C_{A}(S)\right) = S.
   \]
   \leanok
   \lean{double_centralizer}
   % \uses{lem:simple-centralizer, lem:dim-centralizer}
 \end{theorem}
 \begin{proof}
  It is straightforward that $S \le C_{A}\left(C_{A}(S)\right)$. By~\cref{lem:simple-centralizer}, $C_{A}(S)$ is simple, hence $\dim_{F}C_{A}\left(C_{A}(S)\right)\cdot \dim_{F}C_{A}(S)=\dim_{F}A=\dim_{F}C_{A}(S)\cdot \dim_{F}S$ (by applying~\cref{lem:dim-centralizer} twice), i.e. $\dim_{F}C_{A}\left(C_{A}(S)\right) = \dim_{F}S$. This equality of dimension gives us the desired result.
\end{proof}

 
%%% Local Variables:
%%% mode: LaTeX
%%% TeX-master: "../print"
%%% End:


%%FILE blueprint/src/content/frob.tex
\chapter{Frobenius Theorem}

In this chapter we will prove the classical result bla

\begin{lemma}
    If an $\mathbb{R}$ algebra has dimension one, then it is isomorphic to $\mathbb{R}$.
    \leanok
    \lean{rank_1_D_iso_R}
\end{lemma}
\begin{proof}
    Trivial.
\end{proof}

%%FILE blueprint/src/content/galois_descent.tex
% \chapter{Galois Descent}

% \section{Preliminary results}
% Let $k$ be a field and $V$ a $k$-vector space. We denote by $V^{\star}$ the dual space of $V$.

% \begin{lemma}\label{lem:dual-tensor-power}
%   There is a map $(V^{\star})^{\otimes p} \to (V^{\otimes p})^\star$
%   \leanok
%   \lean{dualTensorPower}
% \end{lemma}
% \begin{proof}
% Indeed, let $f_1 \otimes, \dots, \otimes f_p \in (V^{\star})^{\otimes p}$ and $v_1 \otimes, \dots, \otimes v_p \in V^{\otimes p}$. We define the map by considering the product $\prod_{i} f_i (v_i)$.
% \end{proof}

% \begin{lemma}[Extending scalars]\label{lem:extend-scalars}
%   \leanok%
%   \lean{PiTensorProduct.extendScalars, Module.Dual.extendScalars}
%   Let $K/k$ be a field extension, then we can extend scalars:
%   \begin{itemize}
%     \item a $k$-linear map $\bigotimes_{i} V \to \bigotimes_{i} K \otimes_{k} V$;
%     \item a $K$-linear map $K \otimes V^{\star} \to \left(K \otimes_{k} V\right)^{\star}$.
%   \end{itemize}
% \end{lemma}

% \begin{proof}
%   $\bigotimes_{i} v_{i} \mapsto \bigotimes_{i} (1\otimes v_{i})$ and $(a \otimes f) \mapsto (x \otimes v) \mapsto f(v)\cdot ax$ will do.
% \end{proof}

% \section{$(p, q)$-tensor}\label{sec:p-q-tensor}

% Let $k$ be a field and $V$ a $k$-vector space.

% \begin{definition}\label{def:tensor-of-type}%
%   For any $p, q \in \mathbb{N}$, a $V$-tensor of type $(p,q)$ or a $(p, q)$-tensor is an element of the tensor product $V^{\otimes p} \otimes {(V^*)}^{\otimes q}$.
%   \leanok%
%   \lean{TensorOfType}
% \end{definition}


% \begin{lemma}
%   If $v$ is a $V$-tensor of type $(p,q)$, then it induces a linear map $V^{\otimes q}\to V^{\otimes p}$.
%   \leanok%
%   \lean{TensorOfType.toHom}
% \end{lemma}
% \begin{proof}
%   Let $v = x \otimes f$ where $f$ can be seen as ${(V^{\otimes q})}^{\star}$ under
%   \cref{lem:dual-tensor-power} and $x \in V^{\otimes p}$, thus if $y \in V^{\otimes q}$, we can obtain an elment in $V^{\otimes p}$ by $f(y)\cdot x$
% \end{proof}

% Let $W$ be another $k$-vector space such that $e : V \cong W$.

% \begin{lemma}\label{lem:tensor-of-type-congr}
%   $e$ induces an isomorphism between $V$-tensor of type $(p,q)$ and $W$-tensor of type $(p, q)$. Furthermore, $\mathsf{1} : V \cong V$ induces theidentity and the linear map induced by $e_{1}\circ e_{2}$ is equal to the composition of linear map induced by $e_{1}$ and $e_{2}$.
%   \leanok%
%   \lean{TensorOfType.congr}
% \end{lemma}
% \begin{proof}
%   Indeed, $e$ induces an automorphism on $V^{\otimes p}$ while $e^{-1}$ induces an automorphism on $W^{\otimes q}$. Checking the functoriality is trivial.
% \end{proof}

% \begin{lemma}[extend scalars]\label{lem:extend-scalars-pq-tensor}
%   \leanok%
%   \lean{TensorOfType.extendScalars}
%   Let $K/k$ be a field extension, $\Phi$ be $V$-tensor of type $(p,q)$ over $k$, then there is a $K\otimes V$-tensor of type $(p,q)$ over $K$.
% \end{lemma}
% \begin{proof}
%   We construct a bilinear map $V^{\otimes p} \to (V^{\star})^{\otimes q} \to (K \otimes_{k} V)^{\otimes p}\otimes_{K} \left({{(K \otimes_{k} V)}^{\star}}\right)^{\otimes q}$
%   We have a map $\alpha : (V^{\star})^{\otimes} q \to ((K \otimes_{k} V)^{\star})^{\otimes q}$ by composing $(V^{\star})^{\otimes q} \to (K \otimes_{k} V^{\star})^{\otimes q}$ and $(K \otimes_{k} V^{\star}) \to (K \otimes_{k} V)^{\star}$ from \cref{lem:extend-scalars}. Thus we can define $v \mapsto f \mapsto \hat{v}\otimes \alpha f$ where $\hat{v}$ is $v$ extended to $(K\otimes_{k} V)^{\otimes p}$ by \cref{lem:extend-scalars} again.
% \end{proof}

% \section{Vector space with tensor of type $(p,q)$}

% Let $k$ be a field.

% \begin{definition}
%   A $k$-vector space with tensor of type $(p, q)$ is a pair of $(V, \Phi)$ where $V$ is a $k$-vector space and $\Phi$ a $V$-tensor of type $(p, q)$.
%   \leanok%
%   \lean{VectorSpaceWithTensorOfType}
% \end{definition}

% \begin{definition}
%   An equivalence between $k$-vector spaces $(V, \Phi)$ and $(W, \Psi)$ with tensor of type $(p, q)$ is a $k$-equivalence $e : V \cong W$ such that $e(\Phi) = \Psi$ where $e$ is seen as an isomorphism between $V$-tensor of type $(p, q)$ and $W$-tensor of type $(p, q)$ under \cref{lem:tensor-of-type-congr}.
%   \leanok%
% \end{definition}

% \begin{lemma}
%   The notion of equivalence between $k$-vector spaces with tensor of type $(p, q)$ is an equivalence relation:
%   \begin{itemize}
%     \item $(V, \Phi)$ is equivalent to itself.
%     \item If $(V, \Phi)$ is equivalent to $(W, \Psi)$, then $(W, \Psi)$ is equivalent to $(V, \Phi)$.
%     \item If $(V_{1}, \Phi_{1})$ is equivalent to $(V_{2}, \Psi_{2})$ and $(V_{2}, \Psi_{2})$ is equivalent to $(V_{3}, \Psi_{3})$, then $(V_{1}, \Phi_{1})$ is equivalent to $(V_{3}, \Psi_{3})$.
%   \end{itemize}
%   \leanok%
%   \lean{VectorSpaceWithTensorOfType.Equiv.refl, VectorSpaceWithTensorOfType.Equiv.symm,  VectorSpaceWithTensorOfType.Equiv.trans}
% \end{lemma}

% \begin{lemma}[extend scalars]\label{lem:extend-scalars-vector-space-with-tensor}
%   \leanok%
%   \lean{VectorSpaceWithTensorOfType.extendScalars}
%   Let $K/k$ be a field extension and $(V, \Phi)$ a $k$-vector space with tensor of type $(p,q)$, then $(V, \Phi)$ can be extended to a $K$-vector space with tensor of type $(p,q)$.
% \end{lemma}
% \begin{proof}
%   Indeed, $(K \otimes_{k} V, \Phi_{K})$ works where $\Phi_{K}$ is obtained via~\cref{lem:extend-scalars-pq-tensor}.
% \end{proof}

%%% Local Variables:
%%% mode: LaTeX
%%% TeX-master: "../print"
%%% End:


%%FILE blueprint/src/content/morita.tex
\chapter{Morita Equivalence}\label{sec:morita}

This chapter intertwine with \cref{sec:wed-artin}: \cref{sec:stacks-074e}
depends on \cref{sec:wed-artin-proof}; while \cref{sec:wed-artin-unique} depends
on \cref{sec:stacks-074e}.

\section{Construction of the equivalence}\label{sec:morita-construction}

Let $R$ be a ring and $0 \ne n\in \mathbb{N}$. In this chapter, we prove that
the category $R$-modules and the category of $\Mat_{n}(R)$-modules are
equivalent. Then we use the equivalence to prove several useful lemmas.

\begin{construction}\label{con:morita-eqv-functor0}
  \leanok
  \lean{matrix_smul_vec_def}
  If $M$ is an $R$-module, we have a natural
  $\Mat_{n}(R)$-module structure on $\widehat{M}:=M^{n}$ given by
  $(m_{ij})\cdot (v_{k})=\sum_{j}m_{ij}\cdot v_{j}$. If $f : M \to N$ is an
  $R$-linear map, then $\widehat{f} : M^{n}\to N^{n}$ given by
  $(v_{i}) \mapsto (f(v_{i}))$ is a $\Mat_{n}(R)$-linear map. Thus we have a
  well-defined functor $\MOD_{R} \Longrightarrow \MOD_{\Mat_{n}(R)}$.
\end{construction}

\begin{remark}
  Note that all modules are assumed to be left modules; when we need to consider
  right $R$-modules, we will consider left $R^{\opp}$-modules instead. We use
  $\delta_{ij}$ to denote the matrix whose $(i,j)$-th entry is $1$ and $0$
  elsewhere. $\delta_{ij}$ forms a basis for matrices.
\end{remark}

\begin{construction}\label{con:morita-eqv-functor1}
  \leanok
  \lean{fromModuleCatOverMatrix.smul_α_coe} If $M$ is a
  $\Mat_{n}(R)$-module, then
  $\widetilde{M} := \{\delta_{ij}\cdot m | m \in M\} \subseteq M$ is an
  $R$-module whose $R$-action is given by
  $r \cdot (\delta_{ij}\cdot m) := (r\cdot \delta_{ij})\cdot m$. More over if
  $f : M \to N$ is a $\Mat_{n}(R)$-linear map,
  $\widetilde{f} : \widetilde{M} \to \widetilde{N}$ given by the restriction of
  $f$ is $R$-linear. Hence, we have a functor
  $\MOD_{\Mat_{n}(R)}\Longrightarrow \MOD_{R}$.
\end{construction}

\begin{theorem}[Morita Equivalence]\label{thm:morita}
  The functors constructed in \cref{con:morita-eqv-functor0} and
  \cref{con:morita-eqv-functor1} form an equivalence of category.
  \leanok
  \lean{moritaEquivalentToMatrix}
  % \uses{con:morita-eqv-functor0, con:morita-eqv-functor1}
\end{theorem}

\begin{proof}
  Let $M$ be an $R$-module, then the unit $\widetilde{\widehat{M}} \cong M$ is
  given by
  \[
    \begin{aligned}
      x \mapsto \sum_{j} x_{j} \\
      (x,0,\dots,0) \mapsfrom x
    \end{aligned}
  \]

  Let $M$ be an $\Mat_{n}(R)$-module, then the counit
  $\widehat{\widetilde{M}}\cong M$ is given by $m \mapsto (\delta_{i0}\cdot m)$.
  This map is both injective and surjective.
\end{proof}

\section{Stacks 074E}\label{sec:stacks-074e}

Let $A$ be a finite dimensional simple $k$-algebra.

\begin{lemma}\label{lem:isomorphic-simple-mod}
  Let $M$ and $N$ be simple $A$-modules, then $M$ and $N$ are isomorphic as
  $A$-modules.
  \leanok
  \lean{linearEquiv_of_isSimpleModule_over_simple_ring}
  % \uses{thm:wed-artin-algebra, thm:morita}
\end{lemma}

\begin{proof}
  By~\cref{thm:wed-artin-algebra}, there exists non-zero $n\in\mathbb N$,
  $k$-division algebra $D$ such that $A\cong \Mat_{n}(D)$ as $k$-algebras. Then
  by~\cref{thm:morita}, we have equivalence of category
  $e : \MOD_{A}\cong\MOD_{D}$. Since simple module is a categorical notion (it
  can be defined in terms monomorphisms), $e(M)$ and $e(N)$ are simple
  $D$-modules. Since $D$ is a division ring, $e(M)$ and $e(N)$ are isomorphic as
  $D$-modules, therefore $M$ and $N$ are isomorphic as $A$-modules.
\end{proof}

\begin{lemma}\label{lem:direct-sum-simple-mod}
  Let $M$ be an $A$-module, there exists a simple $A$-module $S$ such that $M$
  is a direct sum of copies of $S$, i.e. $M \cong \bigoplus_{i \in \iota} S$ for
  some indexing set $\iota$. \leanok
  \lean{directSum_simple_module_over_simple_ring}
  % \uses{thm:wed-artin-algebra, thm:morita, lem:direct-sum-simple-mod2}
\end{lemma}

\begin{proof}
  By~\cref{thm:wed-artin-algebra}, there exists non-zero $n\in\mathbb N$,
  $k$-division algebra $D$ such that $A\cong \Mat_{n}(D)$ as $k$-algebras. Then
  by~\cref{thm:morita}, we have equivalence of category
  $e : \MOD_{A}\cong\MOD_{D}$. Since simple module is a categorical notion (it
  can be defined in terms monomorphisms), $e^{-1}(D)$ is a simple module over
  $A$. Since\label{lem:direct-sum-simple-mod2} $e(M)$ is a free module over $D$, we can write $e(M)$ as
  $\bigoplus_{i\in \iota} D$ for some indexing set $\iota$. By precomposing the
  unit of $e$, we get an isomorphism
  $M \cong e^{-1}\left(\bigoplus_{i\in\iota} D\right)$. We only need to prove
  $e^{-1}\left(\bigoplus_{i\in\iota} D\right)\cong\bigoplus_{i\in\iota}e^{-1}\left(D\right)$.
  This is because direct sum is the categorical coproduct.
\end{proof}

\begin{remark}
  Note that by~\cref{lem:isomorphic-simple-mod}, any two simple $A$-module are
  isomorphic, hence for any $A$-module $M$ and {\em any} simple $A$-module $S$,
  we can write $M$ as a direct sum of copies of $S$.
\end{remark}

\begin{lemma}\label{lem:iso-of-dim-eq}
  Let $M$ and $N$ be two finite $A$-module with compatible $k$-action. Then $M$
  and $N$ are isomorphic as $A$-module if and only if $\dim_{k} M$ and
  $\dim_{k} N$ are equal. \leanok
  \lean{linearEquiv_iff_finrank_eq_over_simple_ring}
  % \uses{lem:direct-sum-simple-mod}
\end{lemma}

\begin{proof}
  The forward direction is trivial as an $A$-linear isomorphism is a $k$-linear
  isomorphism as well. Conversely, suppose $\dim_{k}M = \dim_{k}N$.
  By~\cref{lem:direct-sum-simple-mod}, there exists a simple $A$-module $S$ such
  that $M \cong \bigoplus_{i\in\iota} S$ and $N \cong \bigoplus_{i\in\iota'} S$
  as $A$-modules. Without loss of generality $S \ne 0$, for otherwise we have
  $M \cong N$ anyway. If either of $\iota$ or $\iota'$ is empty, then
  $\dim_{k}M = \dim_{k} N = 0$ implying that $M = N = 0$, we again have
  $M \cong N$. Thus, we assume both $\iota$ and $\iota'$ are non-empty. By
  pulling back the $A$-module structure on $S$ to a $k$-module structure along
  $k \hookrightarrow A$,
  $M, N, S, \bigoplus_{i\in\iota} S, \bigoplus_{i\in\iota'} S$ are all finite
  dimensional $k$-vector spaces. Hence $\iota$ and $\iota'$ are finite. The
  equality $\dim_{k}M=\dim_{k}N$ tells us that $\iota \cong \iota'$ as set,
  hence $M \cong \bigoplus_{i\in\iota}S\cong\bigoplus_{i\in\iota'}S\cong N$ as
  required.
\end{proof}

Let $A \cong \Mat_{n}(D)$ as $k$-algebras for some $k$-division algebra and
$n\ne0$.

\begin{lemma}\label{lem:exists-simple-mod}
  $D^{n}$ is a simple $A$-module where the module structure is given by pulling
  back the $\Mat_{n}(D)$-module structure of $D^{n}$.
  \leanok
  \lean{simple_mod_of_wedderburn}
  % \uses{thm:morita}
\end{lemma}
\begin{proof}
  By~\cref{thm:morita}, we have $\MOD_{A}\cong\MOD_{D}\cong\MOD_{\Mat_{n}(D)}$.
  Since $D$ is a simple $D$-module, $D^{n}$ is a simple $\Mat_{n}(D)$ module and
  consequently, a simple $A$-module.
\end{proof}

\begin{remark}
  Note that any $A$-linear endomorphism of $D^{n}$ is $\Mat_{n}(D)$-linear, and
  vice versa. Thus we have
  $\End_{A}\left(D^{n}\right)\cong \End_{\Mat_{n}(D)}\left(D^{n}\right)$ as
  $k$-algebras.
\end{remark}

\begin{lemma}\label{lem:end-vec-iso}
  $\End_{A}\left(D^{n}\right)$ is isomorphic to $D^{\opp}$ as $k$-algebras.
  \leanok
  \lean{end_simple_mod_of_wedderburn}
  % \uses{thm:morita}
\end{lemma}

\begin{proof}
  Indeed, we calculate:
  \[
    \begin{aligned}
      \End_{A}\left(D^{n}\right) &\cong \End_{\Mat_{n}(D)}\left(D^{n}\right) & \\
                                 &\cong \End_{D}D & \text{by~\cref{thm:morita},}~\MOD_{D}\cong\MOD_{\Mat_{n}D}\\
                                 &\cong D^{\opp}
    \end{aligned}
  \]
\end{proof}

\begin{lemma}\label{lem:end-simple-iso}
  Let $M$ be a simple $A$-module, then $\End_{A} M\cong D^{\opp}$ as
  $k$-algebras.
  \leanok
  \lean{end_simple_mod_of_wedderburn'}
  % \uses{thm:morita, lem:isomorphic-simple-mod, lem:end-vec-iso}
\end{lemma}
\begin{proof}
  By~\cref{thm:morita}, $D^{n}$ is simple as $A$-module; hence
  by~\cref{lem:isomorphic-simple-mod}, $D^{n}$ and $M$ are isomorphic as
  $A$-module. \Cref{lem:end-vec-iso} gives the desired result.
\end{proof}

\begin{remark}
  In particular, if $M$ is a simple $A$-module, then $\End_{A}M$ is a simple
  $k$-algbera.
\end{remark}

\begin{lemma}
  Let $M$ be a simple $A$-module, then $\End_{A}M$ has finite $k$-dimension.
  \leanok
  \lean{end_simple_mod_finite}
  % \uses{thm:wed}
\end{lemma}
\begin{proof}
  By~\cref{thm:wed}, such $D$ and $n$ always exists. Hence we only need to show
  $D^{\opp}$ has finite $k$-dimension. Since $\dim_{k}A=\dim_{k}\Mat_{n}(D)$ are
  both finite, we conclude $D^{\opp}$ is finite as a $k$-vector space by pulling
  back the finiteness along $D \hookrightarrow \Mat_{n}(D)$.
\end{proof}

\begin{remark}
  Note that for all $A$-module $M$, $\End_{\End_{A}M}M$ is a $k$-algebra as
  well, with $k \hookrightarrow \End_{\End_{A}M}M$ given by
  $a \mapsto (x \mapsto a\cdot x)$. Thus, we always have a $k$-algebra
  homomorphism $A \to \End_{\End_{A}M}M$ given by the $A$-action on $M$. When
  $A$ is a simple ring, this map is injective.
\end{remark}

\begin{definition}[Balanced Module]\label{def:balanced-mod}
  \leanok
  \lean{IsBalanced}
  For any ring $A$ and $A$-module $M$, we say $M$ is a
  balanced $A$-module, if the $A$-linear map $A \to \End_{\End_{A}M}M$ is
  surjective.
\end{definition}

\begin{remark}
  Balancedness is invariant under linear isomorphism.
\end{remark}

\begin{lemma}\label{lem:balanced-self}
  For any ring $A$, $A$ is balanced as $A$-module. \leanok
\end{lemma}
\begin{proof}
  If $f \in \End_{\End_{A}M}A$, then the image of $f(1)$ under
  $A \to\End_{\End_{A}}A$ is $f$ again.
\end{proof}

We assume again that $A$ is a finite dimensional simple $k$-algebra.
\begin{lemma}
  Any simple $A$-module is balanced.
  \leanok
  \lean{isBalanced_of_simpleMod}
  % \uses{lem:direct-sum-simple-mod}
\end{lemma}
\begin{proof}
  Indeed, if $M$ is a simple $A$-module, then $A \cong \bigoplus_{i\in\iota} M$
  for some indexing set $\iota$ by~\cref{lem:direct-sum-simple-mod}. Since $A$
  is balanced, $\bigoplus_{i\in\iota}M$ is balanced. Let
  $g \in \End_{\End_{A}M}M$, we can define a corresponding
  $G \in \End_{\End_{\bigoplus_{i}M}}\left(\bigoplus_{i}M\right)$ by sending
  $(v_{i})$ to $(g(v_{i}))$. Since $\bigoplus_{i}M$ is balanced, we know that
  for some $a\in A$, $G$ is the image of $a$ under
  $A \to \End_{\End_{\bigoplus_{i}M}}\left(\bigoplus_{i}M\right)$. Then the
  image of $a$ under $A \to \End_{\End_{A}M}M$ is $g$.
\end{proof}

\begin{lemma}\label{lem:iso-end-end}
  For any simple $A$-module $M$, we have $A \cong \End_{\End_{A}M}M$ as
  $k$-algebras.
\end{lemma}
\begin{proof}
  The canonical map $A \to \End_{\End_{A}M}M$ is both injective and surjective,
  as $M$ is a balanced $A$-module and $A$ is a simple ring.
\end{proof}

%%% Local Variables:
%%% mode: LaTeX
%%% TeX-master: "../print"
%%% End:


%%FILE blueprint/src/content/skolem_noether.tex
\section{Skolem-Noether Theorem}\label{sec:skolem-noether}

Let $K$ be a field, $A$, $B$ be $K$-algebras where $A$ is central simple and
finite $K$-dimensional and $B$ is simple. Let $M$ be a simple $A$-module.

\begin{construction}
  For any $K$-algebra homomorphism $f : B \to A$, we give $M$ a
  $B \otimes_{K}\End_{A}M$-module structure by defining $(b\otimes l)\cdot m$ to
  be $f(b)\cdot l(m)$. To emphasis $f$, we denote $M$ with the
  $B\otimes_{K}\End_{A}M$-module structure by $M^{f}$.
  \leanok
  \lean{IsMod}
\end{construction}

\begin{lemma}
  Let $f : B \to A$ be a $K$-algebra homomorphism, $M^{f}$ is finitely generated
  as a $B\otimes_{K}\End_{A}M$-module.
  \leanok
  \lean{module_inst_findim}
\end{lemma}
\begin{proof}
  Since $M$ is a finite $A$-module and $A$ a finite dimensional $K$-vector
  space, $M$ is a finite dimensional $K$-vector space as well. Suppose
  $S \subseteq M$ generates $M$ as $K$-module, the claim is that $S$ generates
  $M^{f}$ as well. Let $x \in M^{f}$, we write
  $x = \sum_{i}\lambda_{i}\cdot s_{i}$ with $\lambda_{i}\in K$ and $s_{i}\in S$.
  Note that $\lambda_{i}\cdot s_{i} = (\rho(\lambda_{i})\otimes\mathbf{1}_{M})$
  in $M^{f}$ where $\rho : K \to B$ is the map giving $B$ its $K$-algebra
  structure. Hence $x$ is in the span of $S$ in $M^{f}$ as well.
\end{proof}

\begin{remark}
  Given that $B$ is simple, any $k$-algebra homomorphism $f : B\to A$ injective;
  therefore by finite $K$-dimensionality of $A$, $B$ is finite $K$-dimensional
  as well.
\end{remark}

\begin{lemma}
  \label{lem:iso-fg}
  Let $f,g : B \to A$ be two $K$-algebra homomorphisms. Then $M^{f}$ and $M^{g}$
  are isomorphic as $B\otimes_{K}\End_{A}M$-module.
  \leanok
  \lean{iso_fg}
  % \uses{lem:iso-of-dim-eq}
\end{lemma}

\begin{proof}
  By~\cref{lem:iso-of-dim-eq}, it is sufficient to prove
  $\dim_{K}M^{f}=\dim_{K}M^{g}$. But as $K$-vector space, $M^{f}$ and $M^{g}$
  are {\em literally\/} $M$.
\end{proof}

\begin{theorem}[Skolem-Noether]\label{thm:skolem-noether}
  Let $f, g : B \to A$ be two $K$-algebra homomorphism. Then $f$ and $g$ differ
  only by a conjugation. That is there exists a unit $x \in A^{\times}$ such
  that $g = x f x^{-1}$.
  \leanok
  \lean{SkolemNoether'}
  % \uses{lem:exists-simple-mod, lem:iso-fg, lem:iso-end-end}
\end{theorem}

\begin{proof}
  Let $M$ be any simple $A$-module (which exists
  by~\cref{lem:exists-simple-mod}). By~\cref{lem:iso-fg}, we have some
  isomorphism $\phi : M^{f}\cong M^{g}$ as $B\otimes_{K}\End_{A}M$-module. Since
  $M$ is simple, we have that $M$ is a balanced $A$-module
  by~\cref{lem:iso-end-end}. Let $e$ denote the $k$-algerba isomorphism
  $A \cong \End_{\End_{A}M}M$ given by the $A$-action on $M$. Since both $\phi$
  and $\phi^{-1}$ defines an element of $\End_{\End_{A}M}M$, we define
  $a := e^{-1}(\phi)$ and $b := e^{-1}(\phi^{-1})$. Then $ab = 1$ since
  $e(ab)=e(a)\cdot e(b)=\phi\phi^{-1}=1$. We prove that the image of $f$ and
  $afb$ under $e$ are the same; that is for all $x\in B$ and $m\in M$,
  $e(g(x))(m) = e(a f(x) b)(m)$. The right hand side is equal to
  \[
    \begin{aligned}
      e(af(x)b)(m) &=\left(e(a)\circ e(f(x))\circ e(b)\right)(m) \\
                   &=\left(\phi\circ e(f(x))\circ \phi^{-1}\right)(m)\\
                   &=\phi\left(f(x)\cdot\phi^{-1}(m)\right) \\
    \end{aligned}.
  \]
  Similarly, the left hand side is equal to $g(x)\cdot m$. Note that $\phi$ is
  $B\otimes \End_{A}M$-linear. Therefore
  $\phi\left(\left(x\otimes\mathbf{1}\right)\cdot \phi^{-1}(m)\right) = \left(x\otimes\mathbf{1}\right)\cdot m$.
  Unfolding the definition of $M^{f}$ and $M^{g}$, we see this is saying
  precisely $\phi\left(f(x)\cdot\phi^{-1}(m)\right)=g(x)\cdot m$.
\end{proof}



%%% Local Variables:
%%% mode: LaTeX
%%% TeX-master: "../print"
%%% End:


%%FILE blueprint/src/content/unclassified.tex
% \setcounter{chapter}{-1}
\section{A Collection of Useful Lemmas}\label{sec:unclassified}

In section, we collect some lemmas that does not really belong to anywhere.

\subsection{Tensor Product}

\begin{lemma}\label{lem:expand-tensor-in-basis}
  Let $M$ and $N$ be $R$-modules such that $\mathcal{C}_{i\in\iota}$ is a basis
  for $N$, then every elements of $x \in M \otimes_{R} N$ can be uniquely
  written as $\sum_{i\in\iota}m_{i}\otimes \mathcal{C}_{i}$ where only finitely
  many $m_{i}$'s are non-zero
\end{lemma}

\begin{proof}
  Given the basis $\mathcal{C}$, we have $R$-linear isomorphism
  $N \cong\bigoplus_{i\in\iota}R$, hence
  $M\otimes_{R}N \cong \bigoplus_{i\in\iota}(M\otimes_{R}R)\cong\bigoplus_{i\in\iota}M$
  as $R$-modules.
\end{proof}
By switching $M$ and $N$, the symmetric statement goes without saying.

% \noindent\rule{\textwidth}{0.2pt}

\begin{lemma}
  \label{lem:submodule-tensor-submodule}
  Let $K$ be a field, $M$ and $N$ be flat $K$-modules. Suppose $p \subseteq M$
  and $q \subseteq N$ are $K$-submodules, then
  $(p \otimes_{K} N) \sqcap (M \otimes_{K} q) = p \otimes_{K} q$ as
  $K$-submodules.
\end{lemma}

\begin{proof}
  The hard direction is to show
  $(p \otimes_{R} N) \sqcap (M \otimes_{R} q) \le p \otimes_{R} q$. Consider the
  following diagram:
  \begin{center}
    \begin{tikzcd}
      p \otimes_{K} q \ar[r, "u"] \ar[d, "\alpha"] & M \otimes_{K} q \ar[r, "v"] \ar[d, "\beta"] & {}^{M}/_{p} \otimes_{K} q \ar[d, "\gamma"] \\
      p \otimes_{K} N \ar[r, "u'"] & M \otimes_{K} N \ar[r, "v'"] & ^{M}/_{p}
      \otimes_{K} N
    \end{tikzcd}
  \end{center}
  Since $^{M}/_{p}$ is flat, $\gamma$ is injective. Let
  $z \in (p \otimes_{R} N) \sqcap (M \otimes_{R} q) = \im \beta \sqcap \im u'$.
  By abusing notation, replace $z$ with some elements of $M \otimes_{K} q$ and
  continue with $\beta(z)\in\im \beta \sqcap \im u'$. Since
  $v'(\beta(z))=\gamma(v(z))$ and that $\beta(z)\in \im u'$, we conclude that
  $\gamma(v(z))=0$, that is $z\in\ker v=\im u$. We abuse notation again, let
  $z \in p \otimes_{K} q$, we need to show
  $\beta(u(z))\in\im \beta \sqcap \im u'$, but $\beta\circ u=u'\circ\alpha$, we
  finish the proof.
\end{proof}

\subsection{Centralizer and Center}
Let $R$ be a commutative ring and $A$, $B$ be two $R$-algebras. We denote
centralizer of $S\subseteq A$ by $C_{A}S$ and centre of $A$ by $Z(A)$.

\begin{lemma}\label{lem:sup-centralizer}
  Let $S, T$ be two subalgebras of $A$, then
  $C_{A}(S\sqcup T)=C_{A}(S)\sqcap C_{A}(T)$. \leanok
  \lean{Subalgebra.centralizer_sup}
\end{lemma}
This lemma can be generalized to centralizers of arbitrary supremum of
subalgebras.

\begin{lemma}\label{lem:centralizer-inclusion-tensor}
  If we assume $B$ is free as $R$-module, then for any $R$-subalgebra $S$, we
  have that $C_{A\otimes_{R}B}\left(\im\left(S\to A\otimes_{R}B\right)\right)$
  is $C_{A}(S) \otimes_{R} B$
  \leanok
  \lean{centralizer_inclusionRight}
  % \uses{lem:expand-tensor-in-basis}
\end{lemma}
A symmetric statement goes without saying.
\begin{proof}
  Let $w\in C_{A\otimes_{R}B}\left(\im\left(S\to A\otimes_{R}B\right)\right)$.
  Since $B$ is free, we choose an arbitrary basis $\mathcal{B}$;
  by~\cref{lem:expand-tensor-in-basis}, we write
  $w = \sum_{i}m_{i}\otimes_{K}\mathcal{B}_{i}$. It is sufficient to show that
  $m_{i}\in C_{A}(S)$ for all $i$. Let $a \in S$, we need to show that
  $m_{i}\cdot a = a \cdot m_{i}$. Since $w$ is in the centralizer,
  $w \cdot (a\otimes 1) = (a\otimes 1)\cdot w$. Hence we have
  $\sum_{i}(a\cdot m_{i})\otimes\mathcal{B}_{i}=\sum_{i}(m_{i}\cdot a)\otimes\mathcal{B}_{i}$.
  By the uniqueness of~\cref{lem:expand-tensor-in-basis}, we conclude
  $a\cdot m_{i}=m_{i}\cdot a$.
\end{proof}

\begin{remark}
  A useful special case is when $S=A$, then since $C_{A}(A)=Z(A)$, we have
  $C_{A\otimes_{R}B}\left(\im\left(A\to A\otimes_{R} B\right)\right)$ is equal
  to $Z(A)\otimes_{R} B$. Since
  $\im\left(R\otimes_{R}B\to A\otimes_{R}B\right)=\im\left(A\to A\otimes_{R}B\right)$,
  we conclude its centralizer in $A\otimes_{R}B$ is $Z(A)\otimes_{R}B$.
\end{remark}

\begin{corollary}\label{cor:centralizer-tensor-centralizer}
  Assume $R$ is a field. Let $S$ and $T$ be $R$-subalgebras of $A$ and $B$
  respectively. Then $C_{A\otimes_{R}B}\left(S\otimes_{R}T\right)$ is equal to
  $C_{A}(S)\otimes_{R}C_{B}(T)$
  \leanok
  \lean{centralizer_tensor_centralizer}
  % \uses{lem:submodule-tensor-submodule}
\end{corollary}
\begin{proof}
  From~\cref{lem:submodule-tensor-submodule}, $C_{A}(S)\otimes_{R}C_{B}(T)$ is
  equal to
  $\left(C_{A}(S)\otimes_{R} B\right)\sqcap \left(A\otimes_{R}C_{B}(T)\right)$.
  The left hand side $C_{A}(S)\otimes_{R} B$ is equal to
  $C_{A\otimes_{R}B}\left(\im\left(S\to A\otimes_{R}B\right)\right)$ and the
  right hand side is equal to
  $C_{A\otimes_{R}B}\left(\im\left(T\to A\otimes_{R}B\right)\right)$. Hence
  by~\cref{lem:sup-centralizer}, their intersection is equal to
  \[C_{A\otimes_{R}B}\left(\im\left(S\to A\otimes_{R} B\right)\sqcup \im\left(T\to A\otimes_{R}B\right)\right)\]
  This is precisely $C_{A\otimes_{R}B}\left(S\otimes_{R}T\right)$.
\end{proof}

\begin{corollary}
  \label{cor:center-tensor-center}
  Assume $R$ is a field. The centre of $A\otimes_{R} B$ is
  $Z\left(A\right)\otimes_{R}Z\left(B\right)$. 
  \leanok
  \lean{IsCentralSimple.center_tensorProduct}
  % \uses{cor:centralizer-tensor-centralizer}
\end{corollary}
\begin{proof}
  Special case of~\cref{cor:centralizer-tensor-centralizer}.
\end{proof}

\subsection{Some Isomorphisms}

\begin{construction}
  \label{con:self-tensor-opp-iso-end}
  Let $R$ be a commutative ring and $A$ an $R$-algebra. Then we
  have an $R$-algebra homomorphism $A\otimes_{R}A^{\opp} \cong \End_{R}A$ given
  by $a\otimes 1 \mapsto (a\cdot\bullet)$ and
  $1\otimes a \mapsto (\bullet\cdot a)$. When $R$ is a field and $A$ is a finite dimensional central simple algebra, this morphism is an isomorphism by~\cref{cor:alghom-bijective-of-dim-eq}.
\end{construction}

\begin{construction}
  \label{con:end-vec-iso-matrix}
  Let $A$ be an $R$-algebra and $M$ an $A$-module. We have isomorphism
  $\End_{A}\left(M^{n}\right)\cong \Mat_{n}\left(\End_{A}M\right)$ as
  $R$-algebras. For any $f \in \End_{A}\left(M^{n}\right)$, we define a matrix
  $M$ whose $(i,j)$-th entry is
  \[
    x \mapsto f\left(0,\dots,x,\dots,0\right)_{i},
  \]
  where $x$ is at the $j$-th position. On the other hand, if
  $M \in \Mat_{n}\left(\End_{A}M\right)$, we define an $A$-linear map
  $f : M^{n}\to M^{n}$ by
  \[
    v \mapsto \left(\sum_{j}M_{ij}v_{j}\right)_{i}.
  \]
  \leanok
  \lean{endVecAlgEquivMatrixEnd}
\end{construction}

\begin{construction}
  \label{con:matrix-matrix}
  Let $A$ be an $R$-algebra. Then $\Mat_{m}\left(\Mat_{n}(A)\right) \cong \Mat_{mn}(A)$.
  The trick is to think $\Mat_{m}A$ as $\{0,\dots,m-1\} \times \{0,\dots,m-1\} \to A$. Since the indexing set $\{0,\dots,mn-1\}$ bijects with $\left(\{0,\dots,m-1\}\times\{0,\dots,n-1\}\right)$, the isomorphism is just function currying, function uncurrying, precomposing and postcomposing bijections.
  \leanok
  \lean{Matrix.comp}
\end{construction}

\begin{construction}
  \label{con:matrix-tensor-matrix}
  Let $A, B$ be $R$-algebras. Then $\Mat_{mn}\left(A\otimes_{R}B\right)\cong\Mat_{m}(A)\otimes_{R}\Mat_{n}B$ as $K$-algebras. We first construct $R$-algebra isomorphism $A \otimes_{R}\Mat_{n}(R) \cong \Mat_{n}(A)$:
  \[
    \begin{aligned}
      & a \otimes 1 \mapsto \diag{a}~\text{and}~ 1\otimes (m_{ij}) \mapsto (m_{ij})\\
      & \sum_{i,j} m_{ij}\otimes \delta_{ij}\mapsfrom (m_{ij}),
    \end{aligned}
  \]
  where $\diag$ is the diagonal matrix and $\delta_{ij}$ the matrix whose only non-zero entry is at $(i,j)$-th and is equal to $1$.
  Thus $\Mat_{m}(A)\otimes_{R}\Mat_{n}(B)\cong \left(A\otimes_{R}B\right)\otimes_{R} \left(\Mat_{m}(R)\otimes_{R} \Mat_{n}(R)\right)$ as $R$-algebra.
  The Kronecker product gives us an $R$-algebra map $\Mat_{m}(R)\otimes_{R}\Mat_{n}(R) \to \Mat_{mn}(R)$. We want this map to be an isomorphism. By~\cref{lem:simple-ring-iff}, we only need to prove it to be surjective: for all $\delta_{ij} \in \Mat_{mn}(R)$, we interpret $\Mat_{mn}(R)$ as a function $\{0,\dots,m-1\}\times\{0,\dots,n-1\}\to R$, then $\delta_{ij}$ is the image of $\delta_{ab}\otimes \delta_{cd} \in \Mat_{m}(R)\otimes_{R}\Mat_{n}(R)$ where $i = (a, c)$ and $j = (b,d)$. Combine everything together, we see $\Mat_{mn}\left(A\otimes_{R}B\right)$ is isomorphic to $\Mat_{mn}\left(A\otimes_{R} B\right)$ as $R$-algebras.
  \leanok
  \lean{BrauerGroup.matrix_eqv, BrauerGroup.kroneckerMatrixTensor'}
  % \uses{lem:simple-ring-iff}
\end{construction}

%%% Local Variables:
%%% mode: LaTeX
%%% TeX-master: "../print"
%%% End:


%%FILE blueprint/src/content/wed_artin.tex
\section{Wedderburn-Artin Theorem for Simple Rings}\label{sec:wed-artin}

\subsection{Classification of Simple Rings}\label{sec:wed-artin-proof}


\begin{lemma}[minimal ideal of simple rings]\label{lemma:min-ideal-simple-ring}
  Let $A$ be a ring and $I$ a non-trivial minimal left ideal of $A$, then $I$ is
  a simple $A$-module.
  \leanok
  \lean{minimal_ideal_isSimpleModule}
\end{lemma}
\begin{proof}
  Let $J\le I$ be an $A$-submodule of $I$, suppose $J$ is non-trivial, we prove
  that $J=I$. Then the image $J'$ of $J$ under $I \hookrightarrow A$ is a
  non-trivial left ideal of $A$. Since $I \hookrightarrow A$ is injective, it is
  sufficient to prove that $J' = I$. This is because $J'\le I$ and $J' \not< J$.
\end{proof}

\begin{lemma}
  Let $A$ be a simple ring and $I$ a non-trivial left ideal. One can write
  $1 \in A$ as $\sum_{i=0}^{n}x_{i}y_{i}$ for some $x_{i}\in I$ and
  $y_{i} \in A$.
\end{lemma}

\begin{proof}
  Let $I'$ be the two-sided ideal spanned by $I$. Then since $A$ is a simple
  ring, $I' = A$. Thus $1 \in I'$. One can write $1 \in A$ as
  $\sum_{i}a_{i}x_{i}b_{i}$ for some $x_{i}\in I$ and $a_{i},b_{i}\in A$, since
  $I$ is a left ideal $a_{i}x_{i}\in I$ as well.
\end{proof}

Now, we can find the smallest $n$ such that $1\in A$ can be written as
$\sum_{i=0}^{n}x_{i}y_{i}$ for some $x_{i}\in I$ and $y_{i}\in A$. Let us fix
the notations $n$, $x_{i}$ and $y_{i}$

\begin{lemma} The $n$, $x_{i}$ and $y_{i}$ are all non-zero. \leanok
  \lean{Wedderburn_Artin.aux.nxi_ne_zero}
\end{lemma}

\begin{proof}
  If $n$ is $0$, then $1 = 0$ in $A$, but all simple rings are non-trivial. We
  argue by contradiction to prove that all $x_{i}$ and $y_{i}$ are non-zero.
  Assume there exists a $j$ such that $y_{j}\ne0$ implies $x_{j} = 0$. Without
  loss of generality, we assume $j = 0$. Then
  $1 = \sum_{i=0}^{n}x_{i}y_{i}=\sum_{i=1}^{n}x_{i}y_{i}$. This contradicts the
  minimality of $n$.
\end{proof}

\begin{theorem}[Wedderburn]\label{thm:wed}
  Let $A$ be a simple ring and $I$ a non-trivial minimal left ideal. Then there
  exists a non-zero $n \in \mathbb{N}$ such that $A \cong I^{n}$ as $A$-modules.
  \leanok
  \lean{Wedderburn_Artin.aux.equivIdeal}
  \uses{lemma:min-ideal-simple-ring}
  \uses{lemma:min-ideal-simple-ring}
\end{theorem}

\begin{proof}
  We continue to write $1 = \sum_{i=0}^{n}x_{i}y_{i}$ in the shortest possible
  manner. Then we can define an $A$-linear map $g : I^{n}\to A$ by
  $(v_{i})\mapsto \sum v_{i}y_{i}$. Then $g$ is surjective: if $a \in A$, then
  $(ax_{i})$ is mapped to $a$ under $g$. $g$ is injective as well: support
  $g(v_{i})=0=\sum_{i}v_{i}y_{i}$ with $(v_{i})$ not all zero. Without loss of
  generality, we assume $v_{0} \ne 0$, then the ideal $\langle v_{0}\rangle$ is
  equal to $I$ (since $I$ is simple\cref{lemma:min-ideal-simple-ring}). Thus
  $x_{0}\in I = \langle v_{0}\rangle$; implying that $x_{0}=r\cdot v_{0}$ for
  some $r\in A$. Thus
  $1=1 - r\cdot 0 = \sum_{i=0}^{n}x_{i}y_{i}-\sum_{i=0}^{n}r\cdot v_{i}y_{i}$.
  In this way, we cancelled the term at $i=0$, contradicting the minimality of
  $n$. Hence $g$ is an isomorphism.
\end{proof}

\begin{theorem}[Wedderburn-Artin (Ideal)]
  Let $A$ be an Artinian simple ring. There exists a non-zero $n$ and an ideal
  $I \subseteq A$ such that $I$ is simple as an $A$-module and $A \cong I^{n}$
  as $A$-module.
  \label{thm:wed-artin-ideal}
  \leanok
  \lean{Wedderburn_Artin_ideal_version}
  \uses{thm:wed}
\end{theorem}

\begin{proof}
  By~\cref{thm:wed}, we only need a minimal left ideal. Since $A$ is Artinian,
  such ideal exists.
\end{proof}

\begin{theorem}[Wedderburn-Artin (Algebra)] \label{thm:wed-artin-algebra}
 Let $K$ be a field and $B$ an finite dimensional simple algebra over $K$.
  There exists a non-zero $n\in\mathbb{N}$ and a division $K$-algebra $S$ such
  that $B \cong \Mat_{n}(S)$.
  \leanok
  \lean{Wedderburn_Artin_algebra_version}
  \uses{thm:wed-artin-ideal, con:end-vec-iso-matrix}
\end{theorem}

\begin{proof}
  By~\cref{thm:wed-artin-ideal}, we can find a $n$ and a minimal left ideal $I$
  such $A \cong I^{n}$ as $A$-modules. Note that
  ${\left(\End_{B}I\right)}^{\opp}$ is a division ring. Then since
  $B^{\opp}\cong\End_{B}B\cong{\End_{B}(I^{n})}\cong\Mat_{n}\left(\End_{B}I\right)$
  as rings where the final isomorphism is from~\cref{con:end-vec-iso-matrix}, we
  have $e : B\cong \Mat_{n}\left(\End_{B}I\right)^{\opp}$ as rings. We also have
  a $K$-algebra structure on ${\left(\End_{B} I\right)}^{\opp}$ given by
  $(a \cdot f)(x)=f(a\cdot x)$, and this algebra structure promotes the ring
  isomorphism $e$ to a $k$-algebra isomorphism.
\end{proof}

\subsection{Uniqueness of the Classification}\label{sec:wed-artin-unique}
In the previous section, we know that finite dimensional simple $K$-algebra $B$
over is in fact a matrix algebras of a division $K$-algebra $S$. In this
section, we prove that the division algebra $S$ is essentially unique.

\begin{theorem}[Uniqueness of Wedderburn-Artin theorem]\label{thm:wed-artin-uniq}
  Let $B$ be a
  finite-dimensional simple $K$-algebra. Suppose $B$ is isomorphic as
  $k$-algebras to both $\Mat_{n}(D)$ and $\Mat_{n'}(D')$ where $n, n'$ are
  non-zero natural numbers and $D, D'$ are $k$-division algebra, then $n = n'$
  and $D \cong D'$ as $k$-algebras.\label{thm:wed-artin-unique} \leanok
  \lean{Wedderburn_Artin_uniqueness₀}
  \uses{lem:end-simple-iso}
\end{theorem}

\begin{proof}
  Since $D^{n}$ is a simple $B$-module, by~\cref{lem:end-simple-iso}, we see
  that $\End_{A}D^{n}\cong D^{\opp}$ and $\End_{A}D^{n}\cong D'^{\opp}$ as
  $k$-algebras. Thus $D^{\opp}\cong D'^{\opp}$ as $k$-algebras, consequently
  $D\cong D'$ as $k$-algebras as well. Since
  $A \cong \Mat_{n}(D)\cong \Mat_{n'}(D')\cong \Mat_{n'}(D)$ as $k$-algebras and
  $A$ is finite $k$-dimensional, a dimension argument shows that $n=n'$.
\end{proof}

%%% Local Variables:
%%% mode: LaTeX
%%% TeX-master: "../print"
%%% End:

